% ============================================================
%  tez.tex — ITU M.Sc. Thesis Main File
%  "Topological Data Analysis for Recall-Oriented Wearable
%   Fall Detection: Subject-General and Subject-Personalized
%   Evaluation"
%  Author : Kenan Evren BOYABATLI
%  Advisor: Prof. Dr. Atabey KAYGUN
%  Date   : June 2026
% ============================================================

\documentclass[onluarkali,ingilizce,yukseklisans,karton,LisansustuEgitim]{itutez}

% ============================================================
%  THESIS METADATA
% ============================================================
\yazar{Kenan Evren}{BOYABATLI}
\ogrencino{509231205}
\unvan{}

\anabilimdali{Matematik M\"uhendisli\u{g}i Anabilim Dal{\i}}
             {Department of Mathematical Engineering}

\programi{Matematik M\"uhendisli\u{g}i Program{\i}}
         {Mathematical Engineering Programme}

\tarih{HAZ\.IRAN 2026}{JUNE 2026}
\tarihKucuk{Haziran 2026}{June 2026}

\tezyoneticisi{Prof.\ Dr.\ Atabey KAYGUN}
              {\.Istanbul Teknik \"Universitesi}
\tezyoneticisiENG{Prof.\ Dr.\ Atabey KAYGUN}
                 {Istanbul Technical University}

\esdanismani{}{}
\esdanismaniENG{}{}

% ── Turkish title (3 lines) ──────────────────────────────────
\baslik{Giyilebilir Sens\"orlerle D\"u\c{s}me Tespitinde}
       {TDA}
       {}

% ── English title (3 lines) ─────────────────────────────────
\title{Recall-Oriented TDA for}
      {Wearable-Sensor Fall Detection}
      {}

\tezvermetarih{Haziran 2026}{June 2026}
\tezsavunmatarih{Haziran 2026}{June 2026}

% ── Jury (placeholders) ─────────────────────────────────────
\juriBir{Prof.\ Dr.\ Placeholder One}{Istanbul Technical University}
\juriIki{Prof.\ Dr.\ Placeholder Two}{Istanbul Technical University}
\juriUc{Assoc.\ Prof.\ Dr.\ Placeholder Three}{Istanbul Technical University}
\juriDort{}{}
\juriBes{}{}

% ── Fix undefined cls commands ──────────────────────────────
% The Turkish inner cover uses \th@TezYoneticisiTR and
% \th@esdanismanTR which the cls never defines via a public cmd.
% bibnameToC is referenced in thebibliography env but not defined.
\makeatletter
\def\th@TezYoneticisiTR{Prof.\ Dr.\ Atabey KAYGUN}
\def\th@esdanismanTR{}
\makeatother
\providecommand{\bibnameToC}{\if@Ingilizce REFERENCES\else KAYNAKLAR\fi}

% ── Dedication (empty) ──────────────────────────────────────
\ithaf{}

% ── Foreword ────────────────────────────────────────────────
\onsoz{%
% ---------- BEGIN onsoz.tex ----------
I would like to express my sincere gratitude to my advisor, Prof.\ Dr.\ Atabey KAYGUN, for his guidance and support throughout this thesis.

Computing resources used in this work were provided by the National Center for High Performance Computing of Turkey (UHeM) under grant number 4025462026.

\vspace{1.5cm}
\noindent
\begin{tabular*}{\textwidth}{@{\extracolsep{\fill}}l r@{}}
June 2026 & Kenan Evren BOYABATLI \\
\end{tabular*}
% ---------- END onsoz.tex ----------
}

% ── Turkish Abstract ────────────────────────────────────────
\ozet{%
% ---------- BEGIN ozet.tex ----------
Düşmeler; yaşlı bireylerde yaralanmanın, bağımsızlık kaybının ve ölümün başlıca
nedenlerinden biridir. Dünya genelinde yaşlı nüfusun hızla artması, düşme olaylarının
zamanında ve güvenilir biçimde tespit edilmesini bir kamu sağlığı önceliği hâline
getirmektedir. Giyilebilir ataletsel sensörler (ivmeölçer ve jiroskop), bireyin günlük
yaşamını kısıtlamadan sürekli izleme imkânı sunduğundan, otomatik düşme tespiti için
pratik bir donanım altyapısı sağlar. Bu problemde tespit edilemeyen bir düşmenin
(yanlış negatif) maliyeti, yanlış bir alarmın (yanlış pozitif) maliyetinden çok daha
yüksektir; çünkü kaçırılan bir düşme, yerde uzun süre yardımsız kalma ve buna bağlı
ciddi komplikasyonlar anlamına gelebilir. Dolayısıyla düşme tespiti doğası gereği geri
çağırma (recall) odaklı bir sınıflandırma problemidir.

Bu tez, gecikme gömme (delay embedding) tabanlı bir işlem hattı aracılığıyla kısa
giyilebilir ivmeölçer pencerelerine uygulanan kalıcı homolojinin (persistent homology);
değerlendirmesi titiz, öznitelik inşası ilkeli ve uç (edge) cihazlarda çalışabilecek
kadar hafif bir geri-çağırma odaklı düşme tespitini destekleyip desteklemediğini
araştırmaktadır. Çalışmanın merkezindeki iddia, başarının büyük ölçüde sınıflandırıcı
seçiminden değil, sinyali topolojik bir öznitelik uzayına taşıyan temsil haritasından
kaynaklandığıdır.

\vspace{2mm}
\noindent\textbf{Yöntem.}
Önerilen işlem hattı, her üç-eksenli kaydı önce yönelimden bağımsız bir sinyal büyüklük
vektörüne (SMV) dönüştürür. Tüm kayıtlar ortak bir örnekleme frekansı olan 50 Hz'e
yeniden örneklenir ve sabit uzunlukta pencerelere bölünür. Her pencerenin faz uzayı,
Takens'in gömme teoremi anlamında gecikme gömmeyle yeniden inşa edilir; böylece tek
boyutlu skaler sinyalden çok boyutlu bir nokta bulutu elde edilir. Bu nokta bulutu
üzerine bir filtrasyon kurulur ve sıfırıncı ile birinci boyutta kalıcı homoloji
hesaplanır. Elde edilen kalıcılık diyagramlarının (persistence diagram) her biri on
istatistikle özetlenir; bu istatistikler dört sinyal düzeyi tanımlayıcısıyla birleşerek
yirmi dört boyutlu, kompakt bir öznitelik vektörü ($\phi_\lambda$) oluşturur. Bu vektör,
kasıtlı olarak basit tutulan iki sınıflandırıcıya verilir: lojistik regresyon (LR) ve
destek vektör makinesi (SVM). Her ikisi de sınıf dengesizliğini gözetmek üzere dengeli
sınıf ağırlıklarıyla, SVM ayrıca Platt olasılık kalibrasyonuyla eğitilir.

Topolojik temsilin davranışı; pencere uzunluğu, gömme boyutu, gecikme, filtrasyon türü
(Vietoris--Rips, Alpha, SparseRips), uzaklık metriği ve eşik gibi hiperparametrelerden
oluşan bir demete ($\lambda$) bağlıdır. Bu yüksek boyutlu ve koşullu arama uzayı,
ızgara aramasını hesaplama açısından olanaksız kıldığından, hiperparametreler ağaç
yapılı Parzen kestiricisi (TPE) ile yürütülen Bayesçi optimizasyon kullanılarak
seçilir. Her veri kümesi için ayrı bir arama yürütülür ve sınıflandırıcıdan bağımsız
bir amaç fonksiyonu ile en iyi yapılandırma belirlenir.

\vspace{2mm}
\noindent\textbf{Veri Kümeleri ve Değerlendirme.}
Yöntem, üç açık karşılaştırma veri kümesi üzerinde değerlendirilmiştir: MobiFall,
SisFall ve FallAllD veri kümesinin 40 Hz türetilmiş çeşidi olan FAD\_40Hz. Bu veri
kümeleri; sensör konumu, örnekleme frekansı, denek profili ve düşme türleri bakımından
farklılaşarak yöntemin genellenebilirliğini sınamaya olanak tanır. Değerlendirme,
denek-duyarlı bir çerçevede üç protokolle yapılır: rastgele bölünmeye dayanan ve üst
sınır referansı sağlayan naif protokol; her deneğin tamamen dışarıda bırakıldığı,
birincil protokol olan denek-dışı-bırakma (leave-one-subject-out, LOSO) protokolü; ve
her deneğin yalnızca kendi verisiyle modellendiği kişisel protokol. Bu protokolleri,
karar eşiğinin geniş bir aralıkta taranmasına dayanan bir eşik analizi tamamlar. Birincil
başarım ölçütü, geri çağırmayı kesinliğe göre dört kat ağırlıklandıran $F_2$ skorudur;
denekler arası dağılım, $t$-dağılımına dayalı yüzde 95 güven aralıklarıyla raporlanır.
Tüm ölçekleme ve eşik parametreleri yalnızca eğitim katmanından kestirilerek bilgi
sızıntısı önlenir.

\vspace{2mm}
\noindent\textbf{Bulgular.}
LOSO protokolü altında temsil; SisFall'da $0.97$'ye, FAD\_40Hz'de $0.96$'ya ve
MobiFall'da $0.90$'a varan ortalama $F_2$ skoru elde etmiş, geri çağırma değeri ise tüm
veri kümesi--sınıflandırıcı bileşimlerinde $0.95$'in üzerinde kalmıştır. Bu sonuçlar,
temsilin önceden belirlenmiş başarım hedeflerini en az bir sınıflandırıcı için
karşıladığını göstermektedir. Bu etkinliğin kuramsal dayanağı Cohen--Steiner kararlılık
teoremidir; bu teorem, iki kalıcılık diyagramı arasındaki darboğaz (bottleneck)
uzaklığının, girdi sinyallerindeki bozulmanın supremum normuyla sınırlı olduğunu garanti
eder ve dolayısıyla elle tasarlanmış istatistiksel, spektral veya derin öğrenme tabanlı
öznitelik ailelerinde benzeri bulunmayan biçimsel bir gürültü-dayanıklılığı sertifikası
sağlar.

Üç veri kümesi yapısal olarak birbirinden farklı en iyi yapılandırmalara yakınsamıştır:
MobiFall geniş bir pencerede SparseRips kompleksini, SisFall kısa bir pencerede Alpha
kompleksini, FAD\_40Hz ise kısa bir pencerede standart Rips kompleksini seçmiştir. Bu
durum, evrensel bir yapılandırmanın var olmadığını ve en iyi topolojik geometrinin
düşme dinamiklerindeki ve nokta-bulutu yapısındaki gerçek farkları yansıttığını ortaya
koyar. Ayrıca iki sınıflandırıcı arasındaki $F_2$ farkı tüm veri kümelerinde $0.02$'nin
altında kalırken, topolojik hiperparametrelerin ayarlanması $F_2$'yi büyüklük olarak çok
daha geniş bir aralıkta değiştirmektedir. Bu gözlem, başarımın sınıflandırıcı sınıfından
çok temsil haritası tarafından belirlendiği temel iddiasını güçlü biçimde
desteklemektedir.

Son olarak, model yeniden eğitimi gerektirmeyen sonradan (post-hoc) eşik
kişiselleştirmesinin etkisi incelenmiştir. Denek düzeyindeki kazanım dağılımı, tüm veri
kümesi--sınıflandırıcı bileşimlerinde sağa çarpıktır: deneklerin çoğunluğu yalnızca
marjinal fayda sağlarken küçük bir azınlık $F_2$'de $0.10$'dan büyük kazanımlar elde
etmektedir. Bu heterojenlik, tek bir sabit eşik kaymasının tüm denekleri eşit biçimde
iyileştireceği varsayımını çürütmekte ve en uygun eşiğin denek-özgü olduğunu
göstermektedir.

\vspace{2mm}
\noindent\textbf{Sonuç.}
Tezin başlıca katkıları şunlardır: (i) giyilebilir sinyallerden geri-çağırma odaklı
düşme tespiti için kuramsal olarak güvenceli, hesaplama açısından hafif bir topolojik
temsil haritası; (ii) bu temsili üç ayrı veri kümesinde tutarlı biçimde optimize eden
Bayesçi arama yordamı; (iii) naif, LOSO ve kişisel protokolleri bir araya getiren
denek-duyarlı bir değerlendirme çerçevesi; ve (iv) model yeniden eğitimi gerektirmeyen,
çıkarım zamanında uygulanabilen ve denek-özgü heterojenliği ortaya koyan bir eşik
kişiselleştirme analizi. Bulgular bütün olarak, topolojik temsilin basit
sınıflandırıcılarla dahi güçlü ve gürültüye dayanıklı düşme tespiti sağlayabileceğini;
ancak kişiselleştirmenin homojen bir iyileştirme sunmadığını ve denek bazında ele
alınması gerektiğini göstermektedir.

\vspace{3mm}\noindent
\textbf{Anahtar Kelimeler:} düşme tespiti, giyilebilir sensörler, topolojik veri analizi,
kalıcı homoloji, gecikme gömme, geri-çağırma odaklı sınıflandırma.
% ---------- END ozet.tex ----------
}

% ── English Summary ─────────────────────────────────────────
\summary{%
% ---------- BEGIN summary.tex ----------
Falls are a leading cause of injury, loss of independence, and mortality among
the elderly, and wearable inertial sensors offer a practical means of automatic
fall detection. Because a missed fall is far more costly than a false alarm, the
detection problem is inherently recall-oriented. This thesis investigates whether
persistent homology, applied to short wearable accelerometer windows through a
delay-embedding pipeline, can support recall-oriented fall detection that is at
once rigorous in its evaluation, principled in its feature construction, and
lightweight enough for edge deployment.

The proposed pipeline transforms each three-axis recording into an
orientation-independent signal magnitude vector, reconstructs its phase space by
delay embedding in the sense of Takens, and summarises the resulting point cloud
with persistent homology in dimensions zero and one. Ten statistics of each
persistence diagram, together with four signal-level descriptors, form a compact
twenty-four-dimensional feature vector that is passed to deliberately simple
classifiers, namely logistic regression and a support vector machine. The
hyperparameters of the topological representation are selected by Bayesian
optimisation with a tree-structured Parzen estimator. The method is evaluated on
three public benchmark datasets---MobiFall, SisFall, and FAD\_40Hz---under a
subject-aware framework comprising naive, leave-one-subject-out, and
subject-personalised protocols, complemented by a decision-threshold sweep.

Under the leave-one-subject-out protocol the representation attains a mean
$F_2$ score of up to $0.97$ on SisFall, $0.96$ on FAD\_40Hz, and $0.90$ on
MobiFall, with recall above $0.95$ throughout. The difference between the two
classifiers is consistently below $0.02$, whereas tuning the topological
hyperparameters changes $F_2$ by an order of magnitude more, demonstrating that
performance is governed by the representation rather than the classifier. The
optimal topological configuration differs across datasets, reflecting genuine
differences in fall dynamics and point-cloud geometry. Post-hoc threshold
personalisation, which requires no model retraining, yields a strongly
right-skewed distribution of gains: most subjects benefit only marginally while a
small minority gain more than $0.10$ in $F_2$. The Cohen--Steiner stability
theorem provides a formal noise-robustness certificate that distinguishes the
topological representation from competing feature families.

\vspace{3mm}\noindent
\textbf{Keywords:} fall detection, wearable sensors, topological data analysis,
persistent homology, delay embedding, recall-oriented classification.
% ---------- END summary.tex ----------
}

% ── Abbreviations ───────────────────────────────────────────
\kisaltmalistesi{%
% ---------- BEGIN kisaltmalar.tex ----------
\hspace{-3mm}
\begin{tabular}{p{2.2cm}l}
{\bf ADL}   & {\bf:} Activities of Daily Living\\
{\bf CI}    & {\bf:} Confidence Interval\\
{\bf CNN}   & {\bf:} Convolutional Neural Network\\
{\bf CPU}   & {\bf:} Central Processing Unit\\
{\bf CV}    & {\bf:} Cross-Validation\\
{\bf EI}    & {\bf:} Expected Improvement\\
{\bf FAD}   & {\bf:} FallAllD 40~Hz Derived Variant\\
{\bf FN}    & {\bf:} False Negative\\
{\bf FP}    & {\bf:} False Positive\\
{\bf GP}    & {\bf:} Gaussian Process\\
{\bf GPU}   & {\bf:} Graphics Processing Unit\\
{\bf GUDHI} & {\bf:} Geometry Understanding in Higher Dimensions\\
{\bf HAR}   & {\bf:} Human Activity Recognition\\
{\bf HPC}   & {\bf:} High Performance Computing\\
{\bf KKT}   & {\bf:} Karush--Kuhn--Tucker\\
{\bf LOSO}  & {\bf:} Leave-One-Subject-Out\\
{\bf LR}    & {\bf:} Logistic Regression\\
{\bf LSTM}  & {\bf:} Long Short-Term Memory\\
{\bf PD}    & {\bf:} Persistence Diagram\\
{\bf PH}    & {\bf:} Persistent Homology\\
{\bf RBF}   & {\bf:} Radial Basis Function\\
{\bf SMBO}  & {\bf:} Sequential Model-Based Optimisation\\
{\bf SMV}   & {\bf:} Signal Magnitude Vector\\
{\bf SVM}   & {\bf:} Support Vector Machine\\
{\bf TDA}   & {\bf:} Topological Data Analysis\\
{\bf TN}    & {\bf:} True Negative\\
{\bf TP}    & {\bf:} True Positive\\
{\bf TPE}   & {\bf:} Tree-structured Parzen Estimator\\
\end{tabular}
% ---------- END kisaltmalar.tex ----------
}

% ============================================================
%  PACKAGES
% ============================================================

\usepackage{graphicx}
\graphicspath{{./}}

\usepackage{amsmath,amssymb,amsthm,mathtools}

% Theorem environments and math macros (requires amsmath/amsthm)
% ---------- BEGIN defs.tex ----------
% ============================================================
%  defs.tex — Theorem environments and math macros
%  Loaded by tez.tex via \input defs.tex
%  NO \usepackage commands here — all packages are in tez.tex
% ============================================================

% ------------------------------------------------------------
% Theorem environments (numbered within chapters)
% ------------------------------------------------------------
\newtheorem{definition}{Definition}[chapter]
\newtheorem{theorem}{Theorem}[chapter]
\newtheorem{proposition}{Proposition}[chapter]
\newtheorem{lemma}{Lemma}[chapter]
\newtheorem{corollary}{Corollary}[chapter]
\newtheorem{remark}{Remark}[chapter]

% ------------------------------------------------------------
% Math operators
% ------------------------------------------------------------
\DeclareMathOperator{\Rec}{Rec}
\DeclareMathOperator{\Prec}{Prec}
\DeclareMathOperator{\Acc}{Acc}
\DeclareMathOperator{\Var}{Var}
\DeclareMathOperator{\Cov}{Cov}
\DeclareMathOperator{\diam}{diam}
\DeclareMathOperator{\rank}{rank}
\DeclareMathOperator{\im}{im}
\DeclareMathOperator{\sgn}{sgn}
\DeclareMathOperator{\low}{low}
\DeclareMathOperator{\PH}{PH}
\DeclareMathOperator{\Pers}{Pers}
\DeclareMathOperator{\conv}{conv}
\DeclareMathOperator{\EI}{EI}
\DeclareMathOperator{\KL}{KL}
\DeclareMathOperator*{\argmax}{arg\,max}
\DeclareMathOperator*{\argmin}{arg\,min}

% ------------------------------------------------------------
% Math shorthands
% ------------------------------------------------------------
\providecommand{\R}{\mathbb{R}}
\providecommand{\F}{\mathbb{F}}
\providecommand{\Z}{\mathbb{Z}}
\providecommand{\N}{\mathbb{N}}
\providecommand{\E}{\mathbb{E}}
\providecommand{\Prob}{\mathbb{P}}
\providecommand{\1}{\mathbf{1}}
\providecommand{\Sset}{\mathcal{S}}
\providecommand{\W}{\mathcal{W}}
\providecommand{\X}{\mathcal{X}}
\providecommand{\Y}{\mathcal{Y}}
\providecommand{\T}{\mathcal{T}}
\providecommand{\D}{\mathcal{D}}
\providecommand{\PD}{\mathrm{PD}}
\providecommand{\Vect}{\mathbf{Vect}}
\providecommand{\Fcal}{\mathcal{F}}
% ---------- END defs.tex ----------

% booktabs and multirow already loaded by itutez.cls; load
% tabularx and array which are not
\usepackage{tabularx}
\usepackage{array}

\usepackage{algorithm}
\usepackage{algpseudocode}

\usepackage{float}

\usepackage{tikz}
\usetikzlibrary{shapes.geometric,arrows.meta,positioning,fit,backgrounds,decorations.pathreplacing}
\usepackage{adjustbox}

\usepackage{subcaption}  % already loaded by cls; harmless duplicate
\usepackage{enumitem}

% ── Bibliography (biblatex/biber APA style) ─────────────────
\usepackage[style=apa,backend=biber,sortcites=true,
            maxnames=20,minnames=1,dashed=false]{biblatex}
\addbibresource{references.bib}

% Uniform horizontal alignment on every page: the class sets an asymmetric
% binding margin (odd 40 mm / even 26.4 mm left), which makes the text block
% shift left/right between pages. Apply the binding margin on the left of
% every page instead (single-sided layout). Runs after the class hook.
\AtBeginDocument{\setlength{\evensidemargin}{\oddsidemargin}}

% ============================================================
%  DOCUMENT BODY
% ============================================================
\begin{document}

% Number and list headings down to subsection only; paragraphs stay
% unnumbered run-in headings (avoids ugly 5-level "x.y.z.0.w" entries).
\setcounter{secnumdepth}{3}
\setcounter{tocdepth}{3}

% ── Chapters ────────────────────────────────────────────────
% ============================================================
%  BEGIN ch1.tex
% ============================================================
% ─────────────────────────────────────────────────────────────
\phantomsection
\chapter{Introduction}
\label{sec:introduction}

% ─────────────────────────────────────────────────────────────
\section{The Fall Detection Problem}
\label{subsec:motivation}

\paragraph{Problem statement.}
The core problem addressed in this thesis is binary classification of wearable inertial sensor windows: given a fixed-length segment $w$ of an accelerometer (and optionally gyroscope) recording from a person performing daily activities, assign a label $\hat{y}\in\{0,1\}$ indicating whether the segment contains a fall event ($\hat{y}=1$) or an activity of daily living (ADL, $\hat{y}=0$).

Three structural properties of the data make this problem non-trivial and distinguish it from standard time-series classification tasks.

\textbf{Orientation ambiguity.} A wearable sensor is attached to the body in an orientation that varies across individuals, mounting positions, and sessions. The three axis values $\mathbf{x}(t)=(x_1(t),x_2(t),x_3(t))^\top$ therefore depend on a nuisance rotation from the special orthogonal group $\mathrm{SO}(3)$, which is not a fall-relevant quantity. Any feature built from individual axes conflates motion content with placement geometry and cannot generalise across mounting positions.

\textbf{Class imbalance.} Falls are rare events; the positive class (fall) represents between $10\%$ and $30\%$ of windows in the benchmark datasets used here. Standard accuracy-oriented classifiers systematically under-detect falls because the cost of missing a rare event is not reflected in the training objective. The evaluation metric must encode the asymmetric clinical cost explicitly.

\textbf{Inter-subject variability.} Individuals differ substantially in gait patterns, body mass, and fall biomechanics. A model trained on one group of users and evaluated on a held-out user (the subject-general scenario) must generalise across this variability without access to any labelled data from the test subject. Subject-personalisation, which relaxes this requirement by providing per-user data, addresses a complementary question.

The remainder of this chapter documents the clinical significance of the problem, surveys the existing literature, and explains how the approach developed in this thesis addresses each of the three challenges above.

\paragraph{Clinical and societal context.}
Falls are the second leading cause of accidental injury death worldwide, after road
traffic injuries~\parencite{who2021falls}.  Globally, approximately 684,000 fatal falls
occur each year, with adults over the age of 60 accounting for the overwhelming
majority~\parencite{who2021falls}.  An estimated 37.3~million falls annually are severe
enough to require medical attention, representing a sustained and growing burden on
healthcare systems.  As the global population ages---the World Health Organization
projects that the number of people aged 60 and above will reach 2.1~billion by
2050~\parencite{who2022ageing}---the scale of this burden is set to increase substantially
without effective countermeasures.

The medical consequences of a fall extend well beyond the acute injury.  A fall that
goes undetected initiates a ``long lie'': a period during which the fallen person
remains on the ground, unable to summon help, and is exposed to secondary complications
including dehydration, hypothermia, pressure sores, and aspiration pneumonia.  Even
when the acute injury is minor, a long lie lasting more than an hour doubles the
probability of subsequent institutionalisation and is independently associated with
increased one-year mortality~\parencite{wild1981longlie}.  Beyond the physical consequences,
the \emph{fear of falling}---a well-documented post-fall psychological
syndrome~\parencite{tinetti1988fear}---leads to self-imposed activity restriction that
accelerates deconditioning and, paradoxically, further elevates fall risk.  Timely
detection is therefore not merely about comfort; it is a clinical intervention with
measurable outcomes.

Wearable inertial sensors offer a continuous and unobtrusive monitoring modality that
is uniquely positioned to close this detection gap.  A single accelerometer attached
to the waist or carried in a trouser pocket records the three-axis acceleration of the
body's centre of mass throughout the day.  Unlike camera-based or floor-sensor systems,
it imposes no infrastructure requirement and preserves the user's privacy and mobility.
Modern smartphones and smartwatches already contain the necessary hardware and are
carried by a large fraction of the target population, making passive fall detection a
feature that can be deployed at scale at marginal hardware cost.

The contrast between the two kinds of error is the premise that shapes the rest of
this thesis. A missed fall can trigger the long-lie complications described above,
whereas a false alarm costs at most a brief, unnecessary check; the two outcomes are
therefore far from equally serious. For this reason the present work pursues a
\emph{recall-oriented} objective and scores every model with the $F_2$ measure, which
weights recall four times as heavily as precision. In effect the detector is
deliberately tuned to give precedence to the fall event---tolerating the occasional
false alarm rather than risk overlooking a genuine fall---which is the design stance
reflected in the title of this thesis. The formal definition of $F_2$ and the reasoning
behind the four-to-one weighting are developed in the Methodology.

% ─────────────────────────────────────────────────────────────
\section{Related Work}
\label{sec:related}

Wearable fall detection has been approached from three main directions: rule-based
threshold detectors, classical machine-learning pipelines built on handcrafted
features, and end-to-end deep-learning models.  Each direction has produced systems
that achieve strong scores on controlled benchmark datasets, yet each carries
structural limitations that become apparent when the goal shifts from benchmark
performance to robust, personalised deployment.  A fourth line of research applies
persistent homology to time-series classification in domains outside of fall detection.
This section surveys each strand, identifies the specific limitations that motivate
the present work, and closes with an explicit positioning statement.

% ─────────────────────────────────────────────────────────────
\subsection{Threshold-Based Fall Detectors}
\label{subsec:rw-threshold}

\subsubsection{Method family}

The oldest and most widely deployed family of wearable fall detectors operates
directly on the signal magnitude vector
$\mathrm{SMV}(t)=\|\mathbf{a}(t)\|_2$
and raises an alarm whenever $\mathrm{SMV}(t)$ exceeds a fixed scalar threshold~$T$:
\[
  \widehat y(t) = \mathbf{1}\{\mathrm{SMV}(t) > T\}.
\]
Kangas et al.~\parencite{kangas2008fall} established the archetypal version of this
approach on a waist-mounted accelerometer and reported high sensitivity on laboratory
falls.  Subsequent work introduced two-stage detectors that combine an impact
threshold with a post-impact stillness check: a fall is declared only if a large
$\mathrm{SMV}$ spike is followed by a prolonged low-activity period, filtering out
brief vigorous activities that produce high-magnitude excursions without a subsequent
motionless phase~\parencite{igual2013challenges}.  Multi-axis variants replace the scalar
magnitude with component-specific thresholds or combine vertical, horizontal, and
resultant acceleration thresholds in a logical rule~\parencite{delahoz2014survey}.

\subsubsection{Mathematical limitations}

The threshold decision rule $\mathbf{1}\{\mathrm{SMV}(t) > T\}$ has three
fundamental limitations that are mathematical rather than incidental.

\textbf{No probabilistic output.}  The rule produces a hard binary decision with
no associated posterior probability.  It is therefore impossible to apply a
principled recall-precision trade-off: there is no continuous operating parameter
that smoothly exchanges false negatives for false positives.  Adjusting sensitivity
requires changing the hard threshold~$T$, but this adjustment is global---it shifts
the operating point identically for all individuals and all activities.

\textbf{No principled threshold selection.}  The threshold $T$ is chosen by a
combination of empirical cross-validation and clinical convention.  There is no
statistical framework for selecting $T$ that is optimal under a specified loss
function.  In particular, the $F_2$-optimal threshold---which depends on the
class-imbalance ratio in a given deployment population---cannot be identified from
first principles without a probabilistic model.

\textbf{No subject-specific adaptation.}  A single threshold $T$ is applied to all
individuals.  An active person with a vigorous gait produces baseline $\mathrm{SMV}$
values that approach the threshold during normal activity, generating chronic false
alarms.  A frail elderly person with a low-energy fall may not reach the threshold
at all.  The threshold detector has no mechanism to adjust to individual motion
dynamics because it has no subject-level probabilistic model.

\subsubsection{Comparison with the present work}

The present thesis replaces the scalar hard decision with a calibrated posterior
probability $f_\theta(\phi_\lambda(w)) \in [0,1]$.  The decision threshold $\tau$
is then a continuous parameter that can be set to maximize $F_2$ on a per-subject
basis.  The threshold sweep over
$\mathcal{T}=\{0.05,0.10,\ldots,0.95\}$ identifies the subject-optimal $\tau^*(s)$
retrospectively, and the gains---up to $+0.12$ units of $F_2$ on individual
subjects---demonstrate that subject-specific calibration provides value that is
structurally inaccessible to a global threshold detector.

% ─────────────────────────────────────────────────────────────
\subsection{Statistical Feature Methods}
\label{subsec:rw-statistical}

\subsubsection{Method family}

The dominant paradigm in wearable fall detection from the mid-2000s through the
mid-2010s combined handcrafted feature extraction with a standard supervised
classifier.  A window $w$ of length $L$ is mapped to a fixed-dimensional vector
$\phi_{\mathrm{hand}}(w)\in\R^d$ by computing scalar statistics of the raw signal
or its frequency-domain representation.  Typical time-domain features include the
window mean, variance, range, peak value, root mean square, and signal-magnitude
area; typical frequency-domain features include spectral energies in low-frequency
bands, peak frequency, and spectral entropy~\parencite{igual2013challenges,
delahoz2014survey}.  These features are fed to a classifier---most commonly
support vector machines, $k$-nearest neighbours, or decision trees---trained on
labeled windows from a fixed dataset.

Surveys by Igual et al.~\parencite{igual2013challenges} and Delahoz and
Labrador~\parencite{delahoz2014survey} catalogue dozens of systems built on this
template and report accuracy figures in the range $90\%$--$99\%$ on controlled
datasets.  These figures are broadly consistent across methods, suggesting that
the feature choice matters less than the evaluation protocol: studies reporting
the highest scores tend to use random train-test splits ignoring subject membership,
which inflates performance by allowing the model to memorize individual quirks
rather than generalize across persons.

\subsubsection{Mathematical limitations}

\textbf{Absence of a stability certificate.}  The relationship between the
input perturbation $\|w - w'\|$ and the feature perturbation
$\|\phi_{\mathrm{hand}}(w) - \phi_{\mathrm{hand}}(w')\|$ is not bounded in any
useful way for handcrafted statistics.  A small noise perturbation of a signal
window can produce an arbitrarily large change in a feature that includes a
zero-crossing count, an extremum detection, or a threshold-based spectral band
energy.  This instability is not a problem when training and test data come from
the same sensor and subject population, but it becomes significant under sensor
substitution, mounting-position variation, or cross-dataset evaluation.

\textbf{No geometric grounding.}  The feature set is selected by domain intuition
and validated empirically.  There is no underlying theory that explains why the
chosen statistics should separate falls from activities of daily living, nor any
guarantee that the same statistics remain discriminative when the population,
sensor, or activity set changes.  In the language of learning theory, the
inductive bias of the feature extractor is implicit and uncontrolled.

\textbf{Fixed feature dimension with no adaptation.}  The dimension $d$ is fixed by
the feature recipe and cannot be adjusted to the dataset geometry.  Different
datasets exhibit different fall dynamics---controlled laboratory falls versus
naturalistic community falls, elderly populations versus young adults---but the
feature extractor is invariant to these differences.  The hyperparameter of this
approach is the classifier, not the feature extractor itself.

\textbf{Subject-general evaluation conflates variability sources.}  Most papers in
this family report a single aggregate accuracy figure without distinguishing
within-subject repeatability from across-subject generalization.  As a consequence,
it is unclear whether the reported scores reflect a model that generalizes to new
individuals or merely one that memorizes the subjects present in the training fold.

\subsubsection{Comparison with the present work}

The feature map $\phi_\lambda$ of the present thesis differs from handcrafted
features in three respects.

First, the Cohen--Steiner stability theorem provides a formal Lipschitz bound:
small noise perturbations of the input window produce at most equally small
perturbations of the persistence diagram, and hence of $\phi_\lambda$.  This
guarantee is a mathematical theorem, not an empirical observation, and it holds
for any signal, any subject, and any sensor.

Second, the feature extractor is parameterized by $\lambda\in\Lambda$ and the
optimal $\lambda^*_K$ is identified by Bayesian search over a nontrivial search
space.  Different datasets converge to different configurations (SparseRips at 5~s
for MobiFall; Alpha at 1~s for SisFall; Rips at 1~s for FAD\_40Hz), demonstrating
that dataset-specific geometry is captured by the parameterisation.  No handcrafted
feature family has an analogous adaptation mechanism.

Third, the evaluation strictly separates subject-level generalization from
within-subject repeatability by using a leave-one-subject-out protocol with
$R=15$ repetitions, reporting both the mean and the 95\% confidence interval
for each combination of dataset, classifier, and protocol.

% ─────────────────────────────────────────────────────────────
\subsection{Deep Learning Approaches}
\label{subsec:rw-dl}

\subsubsection{Method family}

Convolutional neural networks (CNNs), recurrent networks (LSTMs and GRUs), and
more recently transformer-based architectures have been applied to wearable fall
detection with the aim of learning discriminative representations directly from
raw or minimally processed inertial data.  Wang et al.~\parencite{wang2020falldl}
demonstrated that a multi-scale CNN applied to three-axis accelerometer windows
achieves competitive sensitivity and specificity on MobiFall and SisFall without
manual feature engineering.  Batchuluun et al.~\parencite{batchuluun2022deep} provide
a systematic review of deep-learning fall detection systems published between 2018
and 2022 and identify CNN-LSTM hybrids and attention-based models as the current
state of the art on standard benchmarks, with reported $F_1$ scores above $0.95$
on several controlled datasets.

The apparent performance advantage of deep models over classical methods is partly
attributable to their ability to capture temporal structure across multiple scales
simultaneously, and partly to methodological differences: deep learning papers tend
to use larger datasets and more favorable evaluation protocols than their classical
counterparts, making direct comparison difficult.

\subsubsection{Engineering and scientific limitations}

\textbf{GPU dependency.}  Training a competitive deep network on inertial windows
requires GPU compute for gradient-based optimisation, which is practical in a
research setting but problematic for deployment on the battery-constrained processors
of wearable devices.  Inference may be feasible on a smartphone CPU with quantized
models, but training---and therefore retraining for personalisation---is not.

\textbf{Opacity.}  The features learned by a deep network's intermediate layers are
not interpretable in terms of the physical signal.  There is no explanation for why
a particular window is classified as a fall that a clinician or system designer can
evaluate independently of the training data.  TDA features, by contrast, have a direct
geometric interpretation: H$_0$ statistics measure the spread of the delay-embedded
cloud (cluster count and birth-death lifetime distribution), and H$_1$ statistics
measure the presence of loops (periodic or oscillatory structure).

\textbf{No formal stability guarantee.}  Adversarial-robustness research has
repeatedly demonstrated that small, human-imperceptible perturbations of the input
can reverse the output of a deep classifier.  For a safety-critical fall-detection
system, the absence of a stability certificate is a principled concern: a small
sensor hardware change, a firmware update affecting the analog front-end, or a
change in mounting position can degrade performance in ways that are unpredictable
without re-evaluation.  The Cohen--Steiner bound of TDA provides a certificate that
does not depend on the training distribution.

\textbf{The ``trains once, deploys forever'' paradigm.}  Deep models are trained on
a fixed corpus and deployed without adaptation.  Subject-specific personalisation
requires either fine-tuning---which demands labeled falls from the target
individual---or meta-learning frameworks that introduce significant additional
complexity and still require some labeled falls at deployment time.  In the present
thesis, subject-level adaptation is achieved by a threshold sweep that requires no
labeled falls, no model parameters to update, and no access to the training corpus.

\textbf{Computational cost at scale.}  The pipeline introduced in this thesis
consumed approximately 75~minutes of CPU time for the Bayesian representation search
on a 40-core server.  This cost is paid once per dataset, not per user.
Training a competitive deep model requires comparable or greater wall-clock time
and additionally requires GPU hardware.  For deployments targeting resource-limited
clinical or community settings, the CPU-only profile of the TDA pipeline is a
meaningful engineering advantage.

\subsubsection{Comparison with the present work}

The present thesis does not claim to match or exceed state-of-the-art deep learning
scores on fall detection benchmarks.  The purpose is different: to establish whether
a mathematically well-founded, computationally lightweight, and interpretable
pipeline can achieve strong $F_2$ scores under a rigorous subject-aware evaluation.
The answer is affirmative for two of the three datasets ($F_2 \geq 0.93$ on SisFall
and FAD\_40Hz) and satisfactory for the third ($F_2 \geq 0.88$ on MobiFall).  The
TDA pipeline achieves these results on a CPU without GPU access, with formally
certified feature stability, and with an interpretable feature space---properties
that deep models do not offer jointly.

% ─────────────────────────────────────────────────────────────
\subsection{Topological Data Analysis for Time-Series Classification}
\label{subsec:rw-tda}

\subsubsection{Persistent homology in applied settings}

The use of persistent homology as a feature extractor for time-series classification
has grown substantially since the early 2010s, following the development of
computationally efficient boundary matrix reduction
algorithms~\parencite{edelsbrunner2002topological} and their implementation in
open-source libraries such as GUDHI~\parencite{gudhi} and Ripser.

Pereira and de Mello~\parencite{pereira2015persistent} applied persistence diagrams to
time-series classification in several benchmark domains and showed that topological
features are competitive with or superior to classical statistical features when the
discriminative signal is encoded in the \emph{shape} of short segments rather than
their amplitude or frequency content.  Umeda~\parencite{umeda2017time} extended this line
to multivariate time series and demonstrated that the combination of $H_0$ and $H_1$
features provides complementary information not captured by either alone.

Gidea and Katz~\parencite{gidea2018topological} applied TDA to financial time series,
showing that topological features of sliding-window point clouds detect market
instability with a lead time not accessible to standard spectral methods.  Their
work established the delay-embedding framework---building a point cloud from lagged
copies of a scalar signal and computing its persistent homology---as a broadly
applicable TDA pipeline.  The pipeline introduced in the present thesis follows
this framework but extends it with a systematic search over the embedding and
filtration parameters rather than fixing them a priori.

\subsubsection{TDA for activity recognition and wearable signals}

Ott and Demir~\parencite{ott2017topological} applied TDA features to human activity
recognition (HAR) from smartphone accelerometer data.  Using a fixed sliding window
and pre-selected embedding parameters, they showed that persistence statistics
achieve accuracy comparable to handcrafted features on a standard HAR benchmark.
Xu et al.~\parencite{xu2019topological} similarly used TDA to distinguish activity classes
in a multiclass HAR setting, reporting improvements over baseline statistical
features on several benchmark datasets.

These works establish the viability of persistent homology for wearable inertial
classification, but differ from the present thesis in several ways that limit their
direct applicability to the fall detection problem.

\textbf{Objective.}  Both Ott et al.\ and Xu et al.\ optimise accuracy in a balanced
multiclass setting.  Fall detection is a binary problem with severe class imbalance,
and the appropriate objective is recall-oriented ($F_2$) rather than accuracy.
The distinction matters: a classifier that maximizes accuracy on an imbalanced dataset
can do so by predicting the majority class (non-fall) for almost all windows, achieving
high accuracy at the cost of near-zero recall.

\textbf{Evaluation protocol.}  The HAR studies cited above use random train-test
splits that do not control for subject membership, which inflates performance
estimates for any method by allowing individual-specific patterns to leak from
training to test.  The present thesis strictly enforces subject separation via
leave-one-subject-out cross-validation.

\textbf{Hyperparameter treatment.}  Prior TDA-for-wearables works fix the embedding
dimension $m$, delay $\tau$, and filtration type $c$ either by prior convention or
by a coarse manual grid.  The present thesis treats the full hyperparameter tuple
$\lambda=(w_\mathrm{sec}, m, \tau, r, c, d_\mathrm{met}, \varepsilon, q)$ as the
primary object of study and identifies the optimal tuple per dataset via 200
iterations of Bayesian search.  The finding that the optimal complex type differs
across all three datasets (SparseRips, Alpha, Rips) is a direct consequence of this
systematic search; it could not have been discovered by fixing the complex type to a
convention.

\textbf{Personalisation.}  Neither Ott et al.\ nor Xu et al.\ report subject-level
personalisation results or analyse the distribution of per-subject gains.  The
present thesis treats the distribution of the sweep gain $\Delta F_2(s)$ as a
scientific object in its own right, showing that it is right-skewed and
dataset-dependent and that its coefficient of variation exceeds $50\%$ in all cases.

% ─────────────────────────────────────────────────────────────
\subsection{Positioning of the Present Work}
\label{subsec:rw-positioning}

Table~\ref{tab:rw-comparison} summarises the structural differences between the
four approaches surveyed above and the present thesis across six dimensions.

\begin{table}[h!]
\centering
\caption{Structural comparison of fall detection and TDA-for-wearables approaches.
  \checkmark\ = property satisfied; $\times$ = property absent; $\sim$ = partially
  satisfied.}
\label{tab:rw-comparison}
\footnotesize
\begin{tabularx}{\textwidth}{lXXXXX}
\toprule
\textbf{Property} &
\textbf{Threshold} &
\textbf{Stat.\ features} &
\textbf{Deep learning} &
\textbf{TDA-HAR} &
\textbf{This work} \\
\midrule
Probabilistic output           & $\times$ & \checkmark & \checkmark & \checkmark & \checkmark \\
Formal stability certificate   & $\times$ & $\times$   & $\times$   & \checkmark & \checkmark \\
Dataset-specific tuning        & $\times$ & $\times$   & $\sim$     & $\times$   & \checkmark \\
Subject-aware evaluation       & $\times$ & $\sim$     & $\sim$     & $\times$   & \checkmark \\
Recall-oriented objective      & $\times$ & $\sim$     & $\sim$     & $\times$   & \checkmark \\
CPU-only deployment            & \checkmark & \checkmark & $\times$ & \checkmark & \checkmark \\
Subject-level personalisation  & $\times$ & $\times$   & $\times$   & $\times$   & \checkmark \\
\bottomrule
\end{tabularx}
\end{table}

The table makes the novelty of the present work concrete.  No prior approach combines
a formal stability certificate (Theorem), dataset-specific parameterisation (via
Bayesian search), subject-aware evaluation (LOSO with repetitions), a recall-oriented
objective ($F_2$), and subject-level personalisation (threshold sweep) within a
single unified pipeline.  The closest prior work---TDA for HAR---satisfies the
stability and CPU constraints but lacks all three evaluation properties.

The gap that this thesis directly addresses is the intersection of the last four
rows: a principled, recall-oriented, subject-aware, personalised fall detector that
carries a mathematical stability guarantee and is deployable on consumer hardware
without GPU acceleration.

% ─────────────────────────────────────────────────────────────
\section{Limitations of Existing Approaches}
\label{subsec:limitations-related}

Despite the practical appeal of wearable fall detection, the path from sensor signal
to reliable deployment is obstructed by a set of well-known technical difficulties.
A brief taxonomy of existing approaches clarifies where the gaps lie.

\paragraph{Threshold-based detectors.}
The earliest and most widely deployed family of methods compares the signal magnitude
vector $\|\mathbf{a}(t)\|_2$ against a fixed threshold~$T$: a fall is declared
whenever $\|\mathbf{a}(t)\|_2 > T$.  These methods are computationally trivial and
require no training data.  Their fundamental limitation is the absence of a principled
calibration mechanism.  The threshold~$T$ is set empirically, the output is binary
with no associated uncertainty, and there is no statistical framework for adjusting
the operating point to the false-alarm tolerance of a particular deployment environment.
In practice, fixed-threshold detectors achieve high recall on controlled datasets but
at false-alarm rates that are clinically unacceptable for continuous wear: vigorous
activities of daily living (running, stair descent, jumping) frequently exceed the
same magnitude threshold as a fall impact.

\paragraph{Statistical and spectral feature methods.}
A substantial literature constructs handcrafted feature vectors from short
accelerometer windows---means, variances, spectral energies, signal-magnitude
area, zero-crossing rates---and trains standard classifiers on these
features~\parencite{igual2013challenges, delahoz2014survey}.  These methods improve over
threshold-based approaches by producing calibrated probability outputs and by
exploiting the multi-dimensional structure of the signal beyond a single scalar
statistic.  However, their features are \emph{ad hoc}: selected by domain intuition
rather than derived from a principled geometric theory of the signal.  There is
no formal stability guarantee---small perturbations of the input signal can cause
arbitrary changes in the feature vector if the perturbation happens to cross a
feature threshold.  As a result, the transferability of trained models across
sensors, sampling rates, and subject populations is empirically fragile.

\paragraph{Deep learning approaches.}
Convolutional and recurrent neural networks have achieved competitive performance on
fall detection benchmarks by learning end-to-end feature representations from raw
inertial data~\parencite{batchuluun2022deep, wang2020falldl}.  The gains over handcrafted
features are real, but the engineering trade-offs are severe.  Training requires large
labelled datasets, GPU compute, and hyperparameter searches whose cost scales with
model complexity.  The resulting models are opaque: the learned features are not
interpretable, and there is no geometric or mathematical explanation for why a
particular signal window is classified as a fall.  Deployment on edge hardware
(smartwatches, embedded sensors) is complicated by memory footprint and energy
consumption.  Most critically for the personalization question, deep models are
trained once on a fixed corpus and deployed without adaptation---the ``trains once,
deploys forever'' paradigm---which systematically ignores the substantial variation
in motion dynamics across individuals.

\paragraph{The personalization gap.}
Across all three families, a common structural weakness emerges: the decision rule is
calibrated at the population level and applied uniformly to every individual.  This
is statistically equivalent to assuming that the optimal decision boundary is the
same for all subjects.  In reality, individuals differ in gait dynamics, body geometry,
sensor placement geometry, and fall style, so the optimal threshold for a classifier
applied to subject~$s$ will in general differ from the population-level optimum.
Recognising this, some recent works report \emph{personalized} evaluation results, but
typically through isolated case studies rather than a systematic statistical
characterisation.  The question of whether subject-level adaptation is uniformly
beneficial, or only selectively so, remains open.

% ─────────────────────────────────────────────────────────────
\section{Topological Data Analysis for Wearable Signals}
\label{subsec:tda-motivation}

The limitations catalogued in Section~\ref{subsec:limitations-related} share a common structural origin: they operate on representations---axis-wise features, spectral energies, learned embeddings---that do not capture the \emph{geometric shape} of the signal trajectory in phase space, and none of them carries a mathematical certificate of stability with respect to sensor noise. Topological data analysis (TDA), and persistent homology in particular~\parencite{edelsbrunner2010computational, carlsson2009topology}, directly addresses both deficits.

\paragraph{From orientation to geometry.}
The approach begins with a projection that eliminates the orientation ambiguity. The signal magnitude vector $\mathrm{SMV}(t)=\|\mathbf{x}(t)\|_2$ is invariant under any rotation $\mathbf{R}\in\mathrm{SO}(3)$: $\|\mathbf{R}\mathbf{x}(t)\|_2=\|\mathbf{x}(t)\|_2$. This single projection removes the dependence on sensor mounting without any calibration or alignment procedure, and it preserves the full energy of the signal. The resulting scalar series $\mathrm{SMV}(t)$ is the input to all subsequent processing.

\paragraph{From scalar time series to point cloud.}
A short window $w$ of the scalar series is lifted to a finite point cloud $P_\lambda(w)\subset\R^m$ by \emph{delay embedding}~\parencite{takens1981detecting}: successive overlapping $m$-tuples of lagged values form the coordinates of each point. The Takens theorem guarantees that, for $m$ exceeding twice the intrinsic dimension of the underlying dynamical system, this embedding is generically injective---the point cloud faithfully represents the phase-space trajectory of the system that produced the signal. The classification problem is now one of comparing geometric objects: two point clouds derived from two windows of the same length.

A fall event and a daily-activity window produce point clouds of qualitatively different topology. A fall window traces a brief, non-recurrent excursion in $\R^m$: the point cloud is elongated, cluster-like, with strong zeroth-degree ($H_0$) topological structure. A walking or stair-climbing window traces a roughly periodic orbit: the point cloud contains closed loops and exhibits richer first-degree ($H_1$) structure. This topological distinction is visible at the level of individual windows and is independent of the raw amplitude or energy of the signal---it is a \emph{shape} difference, not a \emph{size} difference.

\paragraph{Persistent homology: tracking topology across scales.}
Persistent homology quantifies the topological shape of the point cloud by constructing a nested sequence of simplicial complexes---the \emph{filtration}---indexed by a scale parameter, and tracking how topological features (connected components, loops, voids) are born and destroyed as the scale increases. The result is a \emph{persistence diagram}: a multiset of $(b,d)$ pairs recording the birth scale $b$ and death scale $d$ of each feature. A feature with large persistence $d-b$ is genuine; a feature with small persistence is noise.

Two properties of persistent homology are decisive for the fall-detection setting.

\textbf{Provable stability.}  The Cohen--Steiner stability theorem~\parencite{cohen2007stability}
states that the bottleneck distance between two persistence diagrams is bounded above
by the supremum norm of the perturbation between the corresponding input functions:
\[
  d_B(D(f),\,D(g)) \;\leq\; \|f - g\|_\infty.
\]
Applied to inertial windows, this guarantees that sensor noise, quantisation artefacts, or minor recording-condition variations produce at most correspondingly small perturbations of the topological descriptor. This is a \emph{mathematical theorem}, not an empirical observation, and it holds independently of the training dataset, the subject population, or the sensor hardware. The rule-based and classical feature methods reviewed in Section~\ref{subsec:limitations-related} possess no analogous certificate; deep-learning methods possess stability guarantees only under specific architectural assumptions that are rarely verified in practice.

\textbf{Fixed-dimensional vectorisation.}  The persistence diagram is converted into
a fixed 24-dimensional feature vector $\phi_\lambda(w)\in\R^{24}$ by computing
summary statistics of the birth-death pairs in homological degrees~$k=0$ and~$k=1$.
The dimension 24 is independent of the window length, sampling rate, or number of
simplices in the filtration.  This uniformity allows a single classifier to be applied
across windows of different lengths and datasets recorded at different rates---a
property that is not shared by spectral features (dimension grows with signal length)
or deep-network embeddings (dimension is a design parameter tied to model architecture).

\paragraph{Existing applications and the distinguishing features of this work.}
TDA has been applied to time-series classification across several
domains~\parencite{gidea2018topological, pereira2015persistent, umeda2017time}, and
to activity recognition from inertial
sensors~\parencite{xu2019topological, ott2017topological}.  The distinguishing features
of the present work are: (i) an explicit $\mathrm{SO}(3)$-invariant preprocessing chain before the topological analysis; (ii) a recall-oriented objective ($F_2$) that is propagated consistently through hyperparameter search, threshold selection, and evaluation; (iii) systematic Bayesian optimisation of the full hyperparameter tuple $\lambda$ over multiple benchmark datasets; and (iv) a statistical treatment of subject-level personalisation as a distributional question across the full subject population, rather than a single-subject case study.

% ─────────────────────────────────────────────────────────────
\section{The Threshold Decision as a First-Class Component}
\label{subsec:threshold-intro}

Wearable fall detection is an instance of binary classification under class imbalance:
falls are rare relative to activities of daily living, so the positive class is
significantly underrepresented in any realistic deployment.  Under such imbalance,
the choice of decision threshold $\tau$ in the rule
$\widehat{y}(w)=\mathbf{1}\{f_\theta(\phi_\lambda(w))\geq\tau\}$
is not a cosmetic detail but a critical design parameter that directly controls the
precision-recall trade-off.

As established in Section~\ref{subsec:motivation}, the thesis prioritises \emph{recall}
through the recall-weighted $F_2$ score. Under class imbalance the decision threshold
$\tau$ is the concrete lever that realises this preference: lowering $\tau$ raises
recall at the expense of precision, so the operating point must be chosen deliberately
rather than left at the default $0.5$.

Beyond the population-level threshold, the thesis investigates \emph{subject-level
threshold personalisation}.  The key observation is that the posterior probability
function $f_\theta\circ\phi_{\lambda^*}$ may be systematically miscalibrated
differently for different subjects: a subject whose fall windows produce
lower-than-average posterior scores needs a lower threshold to achieve high recall,
while a subject whose ADL windows produce elevated scores needs a higher threshold to
control false alarms.  A sweep of $\tau$ over the grid
$\mathcal{T}=\{0.05,0.10,\ldots,0.95\}$ identifies the per-subject optimal threshold
$\tau^*(s)$ retrospectively.

The computational cost of this sweep is $O(|\mathcal{T}|\cdot|W_s|)$ comparisons,
where $|\mathcal{T}|=19$ and $|W_s|$ is the window count for subject~$s$.  This
amounts to less than one second of wall time for all subjects in any of the three
datasets, against a Bayesian representation search that consumes approximately
75~minutes per dataset.  The ratio exceeds $10^3$: threshold personalisation is a
computationally negligible add-on to an already-trained model.  It requires no model
retraining, no access to the original training data, and no transmission of data to
a server---it is executable on-device at inference time from a small per-subject
calibration set.  This combination of properties---near-zero cost, privacy
preservation, and measurable $F_2$ gain---makes the threshold decision a first-class
engineering contribution of the thesis, not merely an evaluation detail.

% ─────────────────────────────────────────────────────────────
\section{Contributions}
\label{subsec:contributions}

The thesis makes the following five contributions.

\textbf{Contribution 1: A mathematically grounded TDA pipeline for wearable fall
detection.}  We construct a feature map $\phi_\lambda\colon\R^{n\times 3}\to\R^{24}$
that transforms a raw three-axis accelerometer window into a fixed 24-dimensional
topological descriptor via delay embedding, simplicial filtration, persistent homology
computation, and vectorization.  The map is parameterised by a hyperparameter tuple
$\lambda=(w_\mathrm{sec},m,\tau,r,c,d_\mathrm{met},\varepsilon,q)$ whose components
control the geometric construction.  The Cohen--Steiner stability theorem provides a
formal noise-robustness certificate for the resulting features.

\textbf{Contribution 2: A classifier-agnostic Bayesian hyperparameter optimisation
framework.}  A single Optuna/TPE study~\parencite{optuna} with a classifier-averaged
$F_2$ objective identifies $\lambda^*_K$ for each dataset~$K$ in 200 trials.  The
averaged objective ensures that the recovered representation supports both logistic
regression and support vector machines without being over-fitted to either.  We show
that this optimisation identifies dataset-specific configurations (H2) and that the
variation in $F_2$ across trials---up to 0.18 units on MobiFall---far exceeds the
variation attributable to classifier choice (H3).

\textbf{Contribution 3: A three-protocol subject-aware evaluation framework.}  Three
evaluation protocols are reported jointly for each dataset and classifier: the Naive
(random split, upper bound), the General (leave-one-subject-out, primary), and the
Personal (per-subject split, idealised personalisation reference).  The LOSO protocol
with $R=15$ independent repetitions provides repetition-level confidence intervals via
a $t$-statistic, and the joint reporting of all three protocols makes the
train--test leakage implicit in naive evaluation visible quantitatively.

\textbf{Contribution 4: A statistical characterisation of threshold heterogeneity
as a distributional phenomenon.}  The subject-level sweep gain
$\Delta F_2(s) = F_2(\tau^*(s), s) - F_2(\hat\tau_0(s), s)$ is analysed as a
distribution rather than a scalar average.  We show that the distribution is
right-skewed and dataset-dependent: MobiFall exhibits a mean gain of ${\approx}0.04$--$0.05$
with a coefficient of variation exceeding $50\%$, while SisFall exhibits a mean gain
near zero for most subjects.  This distributional analysis confirms Hypothesis~4
and replaces the anecdotal treatment of personalisation common in prior work.

\textbf{Contribution 5: A deployment pathway to consumer devices.}  The pipeline
is CPU-only, requires no GPU, and executes persistent homology at the window sizes
considered in tens of milliseconds on a desktop-class CPU core---a profile compatible
with ARM Cortex-A smartphone processors and modern smartwatch SoCs.  The threshold
sweep personalisation runs in under one second on device.  We describe a prospective
on-device calibration procedure that requires only a short wearing period and no
server-side computation, and identify one-class anomaly detection from ADL data as
the natural extension that eliminates the cold-start problem entirely.

% ─────────────────────────────────────────────────────────────
\section{Thesis Organisation}
\label{subsec:organisation}

The remainder of the thesis is organised as follows.

\textbf{Chapter~2 (Methodology)} provides the mathematical foundations.  It
introduces the signal-magnitude-vector transformation and delay embedding, develops
the theory of simplicial complexes, persistent homology, and the Cohen--Steiner
stability theorem, describes the feature vectorization and the Bayesian optimisation
procedure, and defines the three evaluation protocols and the statistical inference
framework used throughout the thesis.

\textbf{Chapter~3 (Experiments)} describes the three benchmark datasets, the
experimental setup, and the numerical results: dataset-specific optimal configurations,
full performance tables for all classifier-protocol combinations, and the threshold
sweep analysis.

\textbf{Chapter~4 (Results and Discussion)} interprets the experimental outcomes
against the four research hypotheses, presents the per-subject threshold gain
distributions, synthesises findings across datasets, discusses the deployment
pathway, and identifies limitations and directions for future work.

\textbf{Chapter~5 (Conclusion)} summarises the contributions at three levels:
mathematical, statistical, and engineering.

% ─────────────────────────────────────────────────────────────
\section*{Notation}

Throughout the thesis, $\R$ denotes the real line, $\R^d$ the $d$-dimensional
Euclidean space, and $\mathbf{1}\{\cdot\}$ the indicator function.  Scalars are
denoted by lowercase letters, vectors by bold lowercase, and sets by calligraphic
uppercase.  The symbol $\|\cdot\|_2$ denotes the Euclidean norm and
$\|\cdot\|_\infty$ the supremum norm.  All probabilistic statements are with
respect to the randomness of the class-balancing seed unless stated otherwise.
% ============================================================
%  END ch1.tex
% ============================================================

% ============================================================
%  BEGIN ch3.tex
% ============================================================
% ═════════════════════════════════════════════════════════════
\phantomsection
\chapter{Methodology}
\label{sec:methodology}

This chapter develops the mathematical machinery underlying the fall-detection pipeline, independently of any particular dataset or experimental configuration. The exposition follows the order in which the pipeline actually uses each concept---not the order of mathematical sophistication---so that every construction is introduced only once it is needed by the one before it.

The starting point is a wearable inertial sensor that delivers a discrete, vector-valued signal $\mathbf{x}(t)\in\R^3$ at successive time steps. Three structural challenges arise before any classifier can be applied.

\textbf{Orientation.} The three axis values depend on how the sensor is mounted---a rotation from $\mathrm{SO}(3)$ that varies across wearers and sessions. Any feature built from individual axes is therefore confounded by placement geometry. The first design choice is to project onto the Euclidean norm $\mathrm{SMV}(t)=\|\mathbf{x}(t)\|_2$, which is provably invariant under any element of $\mathrm{SO}(3)$.

\textbf{Temporal geometry.} Once the signal is orientation-invariant, the question is how to represent a short window of it. A fall event and a walking episode produce scalar sequences that differ not primarily in amplitude statistics, but in \emph{geometric shape}: a fall window traces a transient, non-recurrent excursion, while an ADL window traces a roughly periodic orbit. To exploit this distinction, the scalar window is lifted to a finite point cloud in $\R^m$ via delay embedding. The classification problem then becomes one of comparing geometric objects---a setting where topological data analysis is naturally effective.

\textbf{Stability.} A useful feature map must be stable: small sensor noise should not produce large feature perturbations. Persistent homology provides this guarantee through the Cohen--Steiner stability theorem, which bounds the perturbation of the topological descriptor by the supremum norm of the input perturbation. No fixed handcrafted feature family possesses a comparable mathematical certificate.

The methodology is organised as follows. Section~\ref{subsec:signal_rep} (\emph{Time Series Embedding}) processes the raw sensor signal---orientation, sampling rate, windowing---into a single scalar window $w\in\R^L$, independently of any downstream geometric or classification choice; every signal event the pipeline handles falls under this one umbrella. Section~\ref{subsec:embedding} (\emph{The Geometric Problem}) asks why a scalar window should be treated as anything other than a sequence of numbers, and answers by lifting $w$ to a finite point cloud $P_\lambda(w)\subset\R^m$ via delay embedding---the move from a time-series problem to a geometric one. Sections~\ref{subsec:topological_prelim}--\ref{subsec:vectorization} develop the topological machinery that studies the resulting point cloud: simplicial complexes, homology groups, persistence modules, the stability theorem, and their reduction to a fixed feature vector $\phi_\lambda(w)\in\R^{24}$. Section~\ref{subsec:formal} (\emph{Problem Redefinition: The Classification Problem}) takes this feature vector and states, only now, what "solving fall detection" formally means: a classification problem over $\phi_\lambda$, governed by the hyperparameter tuple $\lambda$ that collects every open parameter fixed in the preceding sections, evaluated by a metric chosen for this problem's specific cost asymmetry. Section~\ref{subsec:classifiers} defines the two classifier families that consume $\phi_\lambda$ and produce calibrated fall probabilities. Section~\ref{subsec:bayesopt} develops the Bayesian search framework that will be used in the Experiments chapter to find the optimal $\lambda^*$ by treating the entire pipeline as a black-box objective. Section~\ref{subsec:eval_framework} defines the cross-validation protocols used to evaluate it. Figure~\ref{fig:pipeline-overview} shows the resulting pipeline from raw signal to decision.

% ── Pipeline overview figure ──────────────────────────────────
\begin{figure}[H]
\centering
\resizebox{\linewidth}{!}{%
\begin{tikzpicture}[
  node distance=0pt,
  box/.style={draw, rounded corners=3pt, fill=blue!6,
              minimum height=1.4cm, text width=2.8cm,
              align=center, font=\large},
  arr/.style={->, >=Stealth, line width=1.4pt},
  oplbl/.style={font=\normalsize\itshape, fill=white,
                inner sep=2pt, align=center},
  chalbl/.style={font=\normalsize, text=black!60, align=center}
]
% ── Data-object nodes ────────────────────────────────────────
\node[box] (raw) at (0,0)
  {$\mathbf{x}(t)\in\R^3$\\[2pt]\footnotesize raw sensor};
\node[box] (smv) at (5,0)
  {$w\in\R^L$\\[2pt]\footnotesize SMV window};
\node[box] (pc)  at (10,0)
  {$P_\lambda(w)\subset\R^m$\\[2pt]\footnotesize point cloud};
\node[box] (ph)  at (15,0)
  {$\{D_k\}_{k=0,1}$\\[2pt]\footnotesize persistence diagrams};
\node[box] (dec) at (20,0)
  {$\hat{y}\in\{0,1\}$\\[2pt]\footnotesize decision};

% ── Arrows with operation labels ─────────────────────────────
\draw[arr] (raw) -- node[oplbl,above=4pt]{SO(3)\\projection} (smv);
\draw[arr] (smv) -- node[oplbl,above=4pt]{delay\\embedding} (pc);
\draw[arr] (pc)  -- node[oplbl,above=4pt]{filtration\\$+$ PH} (ph);
\draw[arr] (ph)  -- node[oplbl,above=4pt]{vectorise\\$+$ classify} (dec);

% ── Challenge annotations below ──────────────────────────────
% Brace under raw → smv
\draw[decorate,decoration={brace,amplitude=7pt,mirror}]
  ([yshift=-8pt]raw.south west) -- ([yshift=-8pt]smv.south east)
  node[chalbl,midway,below=10pt]{Orientation};

% Brace under smv → pc
\draw[decorate,decoration={brace,amplitude=7pt,mirror}]
  ([yshift=-8pt]smv.south west) -- ([yshift=-8pt]pc.south east)
  node[chalbl,midway,below=10pt]{Temporal geometry};

% Brace under pc → dec
\draw[decorate,decoration={brace,amplitude=7pt,mirror}]
  ([yshift=-8pt]pc.south west) -- ([yshift=-8pt]dec.south east)
  node[chalbl,midway,below=10pt]{Stability (Cohen--Steiner)};

\end{tikzpicture}%
}%
\caption{%
  Fall-detection pipeline overview. Each node is a mathematical object;
  each edge is a map. The three structural challenges (orientation,
  temporal geometry, stability) addressed in this chapter are annotated
  below the corresponding stages.}
\label{fig:pipeline-overview}
\end{figure}

% ═════════════════════════════════════════════════════════════
\section{Time Series Embedding}
\label{subsec:signal_rep}

Every event the raw sensor signal can produce---a fall, a step, a stumble, a period of rest---first has to be turned into a single, comparable object: a fixed-length, orientation-invariant scalar window. This section builds that object without reference to any particular dataset, sampling rate, or downstream classifier; the result is a scalar window $w\in\R^L$ ready to be treated as a geometric trajectory in the next section.

\subsection{Discrete signal and orientation-invariant projection}
\label{subsubsec:discrete-signal}

A wearable inertial sensor delivers a discrete, vector-valued signal $\mathbf{x}(t)\in\R^3$ whose three coordinates record acceleration along the sensor's own axes. A single recording sampled at rate $f_s$ is the finite discrete signal
\[
  \mathbf{x}\colon \{0,1,\ldots,T\}\longrightarrow\R^3.
\]
These axes are fixed relative to the sensor housing, not to the body: a rotation $\mathbf{R}\in\mathrm{SO}(3)$ that varies across individuals, mounting positions, and even recording sessions maps the body-frame signal to the sensor-frame signal. Axis-wise representations therefore conflate motion content with this nuisance orientation.

\begin{definition}[Signal magnitude vector]
\label{def:smv}
The \emph{signal magnitude vector} of a three-axis signal $\mathbf{x}\colon\Z_{\geq 0}\to\R^3$ is the scalar process
\begin{equation}
  \mathrm{SMV}(t)=\|\mathbf{x}(t)\|_2=\sqrt{x_1(t)^2+x_2(t)^2+x_3(t)^2}.
  \label{eq:smv}
\end{equation}
\end{definition}

\begin{proposition}[$\mathrm{SO}(3)$-invariance]
\label{prop:smv-invariance}
For every rotation $\mathbf{R}\in\mathrm{SO}(3)$ and every $t\geq 0$, $\|\mathbf{R}\mathbf{x}(t)\|_2=\|\mathbf{x}(t)\|_2$. Hence the map $\mathbf{x}\mapsto\mathrm{SMV}$ is invariant under reorientation of the sensor frame.
\end{proposition}

\begin{proof}
$\|\mathbf{R}\mathbf{x}\|_2^2=\mathbf{x}^\top\mathbf{R}^\top\mathbf{R}\mathbf{x}=\mathbf{x}^\top\mathbf{x}=\|\mathbf{x}\|_2^2$, using $\mathbf{R}^\top\mathbf{R}=I$.
\end{proof}

The accelerometer and the gyroscope are processed independently, producing two scalar time series $\mathrm{SMV}_{\mathrm{acc}}$ and $\mathrm{SMV}_{\mathrm{gyr}}$ per recording. All subsequent analysis is carried out on these scalar processes; the loss of directional information is the cost of orientation invariance. Write $s(t)=\mathrm{SMV}(t)$ for the resulting scalar series. A binary label $y\in\{0,1\}$ and a subject index $s\in\Sset$ are associated with each recording. The subject-aware structure is essential: independence of training and test samples is asserted at the level of subjects, not individual windows.

\begin{remark}[From orientation to geometry]
\label{rem:so3-geometry}
Once the signal is reduced to the orientation-invariant scalar series $\mathrm{SMV}(t)$, the classification question becomes: what is the \emph{geometric shape} of a short window of this series? A fall event traces a transient, non-recurrent excursion when the window is later lifted to a delay-embedded space $\R^m$ (Section~\ref{subsec:embedding}); a daily activity traces a roughly periodic orbit. The two shapes are topologically distinct---different numbers of connected components and loops---and persistent homology is designed to detect precisely this distinction in a provably stable way. This is the reason the present section builds towards a \emph{scalar window}, not directly towards a feature vector: shape only becomes visible once a window is treated as a geometric trajectory, which is the subject of Section~\ref{subsec:embedding}.
\end{remark}

\subsection{Spectral sufficiency}
\label{subsubsec:rate}

The fall-discriminative content of inertial wearable signals is concentrated in a low-frequency band. The pre-impact acceleration peak and the impact transient are dominated by frequencies below $20$~Hz; energy above this band consists predominantly of sensor noise and mechanical vibration. This observation provides a theoretical lower bound on the sampling rate required to preserve fall-relevant signal content.

\begin{theorem}[Nyquist--Shannon~\parencite{shannon1949communication}]
\label{thm:nyquist}
A real-valued continuous-time signal whose Fourier transform is supported in $[-B,B]$ is exactly recovered from its samples taken at any rate $f_s\geq 2B$.
\end{theorem}

With $B=20$~Hz as the spectral upper limit of discriminative content, any sampling rate $f_s\geq 40$~Hz is sufficient to preserve all fall-relevant information. The specific target rate used in the experiments and its justification for each dataset are given in Section~\ref{subsec:datasets}.

\subsection{Window operator and extraction}
\label{subsubsec:windows}

For a window length $L\in\N$ and a starting index $t_0\in\{0,\ldots,T-L\}$, the \emph{window operator} extracts a contiguous segment of the scalar series:
\begin{equation}
  W_{t_0,L}\colon\R^{T+1}\longrightarrow\R^{L},
  \qquad
  W_{t_0,L}(s)=(s_{t_0},\,s_{t_0+1},\,\ldots,\,s_{t_0+L-1}).
  \label{eq:window-op}
\end{equation}
The window length is determined by a hyperparameter $w_{\mathrm{sec}}>0$ measuring window duration in seconds and the sampling rate $f_s$ of the specific dataset:
\[
  L:=\lfloor w_{\mathrm{sec}}\cdot f_s\rfloor.
\]

Fall and non-fall trials are windowed by different rules to reflect the structural difference between their respective signal regimes.

For a fall trial $\mathbf{x}$, a single window is extracted around the index of the global SMV maximum:
\begin{equation}
  t^*=\argmax_{t}\,\mathrm{SMV}_{\mathrm{acc}}(t),\qquad
  w^*=W_{t^*-\lfloor L/2\rfloor,\,L}(\mathbf{x}).
  \label{eq:fall-window}
\end{equation}
The centring strategy guarantees that the impact transient is contained in the window regardless of the recording length.

For an ADL trial, consecutive non-overlapping windows of length $L$ are extracted from the start of the recording with stride $L$:
\begin{equation}
  w_j=W_{jL,\,L}(\mathbf{x}),\qquad j=0,1,\ldots,\lfloor T/L\rfloor-1.
  \label{eq:adl-window}
\end{equation}
This produces multiple windows per ADL trial, which is appropriate because daily activities are quasi-stationary and substantially longer than fall events.

The output of this section is a scalar window $w\in\R^L$: orientation-invariant, spectrally justified, and cropped to a fixed temporal support. Nothing so far has used any property of $w$ beyond it being a sequence of $L$ real numbers. The next section asks what is gained by treating that same sequence as tracing a trajectory in a geometric space.

% ═════════════════════════════════════════════════════════════
\section{The Geometric Problem}
\label{subsec:embedding}

A scalar window $w\in\R^L$ carries no explicit geometric structure; it is just a list of numbers. Yet the qualitative difference between a fall and a daily activity is a difference of \emph{shape}: a fall window traces a transient, non-recurrent excursion, while an ADL window traces a roughly periodic orbit (Remark~\ref{rem:so3-geometry}). Scalar summary statistics of $w$---mean, variance, spectral energy---do not see this distinction directly; they are blind to the order in which values occur relative to one another. To make shape a computable quantity, $w$ must first be turned into a \emph{geometric} object: a finite point cloud in $\R^m$. This is the move from a time-series problem to a geometric problem, and it is achieved by delay embedding.

\subsection{Definition}
\label{subsubsec:delay_def}

\begin{definition}[Delay embedding]
\label{def:delay}
Let $s=(s_0,s_1,\ldots,s_{L-1})\in\R^L$ be a finite scalar sequence. For integers $m\geq 2$ (embedding dimension), $\tau\geq 1$ (lag), and $r\geq 1$ (stride), the \emph{$(m,\tau,r)$-delay embedding} of $s$ is the finite point cloud
\begin{equation}
  P_{m,\tau,r}(s)=\bigl\{\mathbf{p}_i\in\R^m:i=0,1,\ldots,N-1\bigr\},
  \label{eq:delay-emb}
\end{equation}
where
\[
  \mathbf{p}_i=\bigl(s_{ir},\,s_{ir+\tau},\,s_{ir+2\tau},\,\ldots,\,s_{ir+(m-1)\tau}\bigr),
\]
and $N=\lfloor(L-(m-1)\tau-1)/r\rfloor+1$ is the number of embedded points.
\end{definition}

Each point $\mathbf{p}_i$ collects $m$ samples of $s$ separated by lag $\tau$, and consecutive points are separated by stride $r$. As $i$ increases the sliding window traces a discrete trajectory in $\R^m$. The cardinality $N$ scales linearly in $L$ and inversely in $r$, which is decisive for the computational cost of the subsequent homology computation. Applying this construction to a window $w$ produces the point cloud $P_\lambda(w):=P_{m,\tau,r}(w)\subset\R^m$ that is the object of study for the remainder of this chapter.

\begin{figure}[!ht]
\centering
\includegraphics[width=\textwidth]{figures/fig_delay_embedding}
\caption{Left: synthetic ADL (walking) and fall accelerometer windows. Right: the
  corresponding two-dimensional delay embeddings ($m=2$, $\tau=3$).  The ADL
  trajectory forms a compact, roughly periodic arc; the fall trajectory produces a
  sparser, elongated point cloud due to the large-amplitude impact spike.  This
  geometric difference motivates the use of persistent homology to distinguish the
  two signal classes.}
\label{fig:delay-embedding}
\end{figure}

\subsection{The Takens embedding theorem}
\label{subsubsec:takens}

The use of delay coordinates is motivated by a foundational theorem of dynamical systems theory. Let $M$ be a smooth compact $d$-dimensional manifold and $\Phi\colon M\to M$ a smooth diffeomorphism (the time-$\Delta t$ flow map of an autonomous dynamical system). Let $h\colon M\to\R$ be a smooth scalar observable. The associated \emph{delay map} of order $m$ is
\[
  \Psi^{h,\Phi}_m\colon M\longrightarrow\R^m,
  \quad
  x\longmapsto\bigl(h(x),h(\Phi(x)),h(\Phi^2(x)),\ldots,h(\Phi^{m-1}(x))\bigr).
\]

\begin{theorem}[Takens~\parencite{takens1981detecting}]
\label{thm:takens}
Let $M$ be a smooth compact $d$-dimensional manifold. For $m\geq 2d+1$, it is a generic property of the pair $(\Phi,h)$ that $\Psi^{h,\Phi}_m\colon M\to\R^m$ is a smooth embedding, i.e., a smooth injective immersion onto its image. In particular, $\Psi^{h,\Phi}_m(M)$ is a smooth $d$-dimensional submanifold of $\R^m$ diffeomorphic to $M$.
\end{theorem}

\emph{Generic} means that the set of pairs $(\Phi,h)$ for which the conclusion fails is contained in a meagre subset of $C^k(M,M)\times C^k(M,\R)$ in the appropriate $C^k$-topology, i.e., a countable union of nowhere-dense sets. Theorem~\ref{thm:takens} states that, for $m$ sufficiently large relative to the intrinsic dimension of the attractor, the reconstructed trajectory is topologically and differentially faithful to the underlying state space; topological invariants computed on $P_{m,\tau,r}(s)$ therefore reflect intrinsic properties of the dynamical system that produced the signal.

\begin{theorem}[Fractal extension, Sauer--Yorke--Casdagli~\parencite{sauer1991embedology}]
\label{thm:syc}
Let $A\subseteq\R^N$ be a compact invariant set of a dynamical system with box-counting (Minkowski) dimension $\dim_B(A)=d_A$. For $m>2d_A$, almost every pair $(\Phi,h)$ (with respect to prevalence in $C^k(\R^N,\R^N)\times C^k(\R^N,\R)$) gives an injective delay map $\Psi^{h,\Phi}_m$ on $A$.
\end{theorem}

Theorem~\ref{thm:syc} relaxes the smooth-manifold assumption of Theorem~\ref{thm:takens}: fractal invariant sets, which are typical for real chaotic systems, can be embedded provided $m$ exceeds twice their box-counting dimension. This is the relevant statement for inertial sensor signals, whose attractors need not be smooth manifolds.

\begin{remark}[Theoretical grounding of the parameter search]
\label{rem:embedding-guidance}
Theorems~\ref{thm:takens} and~\ref{thm:syc} provide positive guidance for the hyperparameter search over $(m,\tau)$. The embedding dimension $m$ must exceed twice the intrinsic dimension of the signal's attractor: for a falling body (a relatively low-dimensional mechanical system) this places the minimum at $m\approx 3$--$5$; for a periodic ADL attractor (a limit cycle, intrinsic dimension $1$) the threshold is $m\geq 3$. Small values of $m$ are therefore ruled out on theoretical grounds before any data is seen. The lag $\tau$ controls the unfolding of the trajectory: too small a $\tau$ compresses all coordinates into a near-diagonal line; too large a $\tau$ destroys temporal correlation. Standard heuristics suggest choosing $\tau$ near the first local minimum of the autocorrelation function, which for fall signals is typically in the range of $2$--$10$ samples at $50$~Hz. The Bayesian search over $(m,\tau)$ in Section~\ref{subsec:bayesopt} can therefore be interpreted as searching within a theoretically motivated region of the hyperparameter space.
\end{remark}

\subsection{Geometric interpretation for fall detection}
\label{subsubsec:geometry}

The geometric content of the delay embedding is what makes it useful for fall detection.

During a \emph{fall}, the SMV signal undergoes a brief high-amplitude transient: a sharp pre-impact rise, a peak at ground contact, and a partial decay. In the reconstruction space the trajectory makes a rapid excursion into an extreme region and does not return to its starting neighbourhood within the window. The resulting point cloud is \emph{elongated and non-recurrent}: a one-dimensional or quasi-one-dimensional sweep with few or no closed loops.

During a \emph{daily activity} such as walking, the SMV signal is approximately periodic with period $T_{\mathrm{walk}}\approx 1$~s. In the reconstruction space the trajectory traces a structured, roughly periodic orbit. The corresponding point cloud is \emph{recurrent and loop-forming}.

This structural dichotomy is precisely the geometry detected by persistent homology in degrees $k=0$ and $k=1$. The zeroth-degree summary $H_0$ records connected components; the first-degree summary $H_1$ records independent one-dimensional cycles. Fall windows are expected to exhibit pronounced $H_0$ behaviour (cluster spread along the excursion) and weak $H_1$ behaviour; ADL windows are expected to exhibit richer $H_1$ structure. The empirical statistics tabulated in Section~\ref{subsec:vectorization} are designed to quantify these expectations.

\begin{remark}[$H_0$ and $H_1$ as complementary detectors]
\label{rem:h0-h1-complementary}
The two homological degrees serve distinct and complementary roles. $H_0$ persistence captures the \emph{spread} and \emph{clustering structure} of the point cloud: a long-lived $H_0$ bar indicates two clusters separated by a gap, which corresponds to the two-phase structure of a fall window (pre-impact cluster and post-impact cluster). $H_1$ persistence captures \emph{loop structure}: a long-lived $H_1$ bar indicates a closed cycle in the point cloud, which arises from the periodic orbit of an ADL window. A classifier operating on both $H_0$ and $H_1$ statistics therefore has access to evidence from two geometrically independent sources. In datasets where $H_0$ alone is informative (falls produce widely separated clusters), $H_1$ acts as a complementary discriminator that helps when $H_0$ evidence is ambiguous. The relative contribution of each degree is examined empirically in the experiments.
\end{remark}

\subsection{Filtration scale control}
\label{subsubsec:eps}

For the Rips and SparseRips families, the filtration is truncated at a data-adaptive scale to bound the size of the simplicial complex. Define the multiset of inter-point distances
\[
  D(P)=\bigl\{d_{\mathrm{met}}(\mathbf{p},\mathbf{p}'):\mathbf{p},\mathbf{p}'\in P,\;\mathbf{p}\neq\mathbf{p}'\bigr\}
\]
and set
\begin{equation}
  \varepsilon_{\max}=\mathrm{Quantile}_q(D(P)),
  \qquad
  q\in\{0.20,0.40,0.60,0.80\}.
  \label{eq:eps_max}
\end{equation}
The percentile $q$ controls the scale at which the filtration is truncated: a low percentile concentrates the analysis on fine-scale local geometry; a high percentile admits coarser global structure at the cost of additional simplices. The Alpha complex requires no such cutoff because its size is intrinsically polynomial in $|P|$ via the Delaunay triangulation.

The delay embedding produces a finite point cloud $P_\lambda(w)\subset\R^m$ for every window $w$. The classification problem has therefore been converted into a purely geometric one: distinguish point clouds by shape. Topological data analysis studies point clouds through exactly this lens---persistent homology---which is developed in the next three sections.

% ═════════════════════════════════════════════════════════════
\section{Topological Preliminaries}
\label{subsec:topological_prelim}

The persistent homology of the filtration depends on objects from point-set and algebraic topology that are not standard in signal processing. This section establishes the necessary background.

\subsection{Topological spaces and homotopy}

\begin{definition}[Topological space]
\label{def:topspace}
A \emph{topological space} is a pair $(X,\mathcal{T})$ where $X$ is a set and $\mathcal{T}\subseteq 2^X$ is a family of subsets, the \emph{open sets}, satisfying
\begin{enumerate}[label=(\roman*)]
  \item $\emptyset\in\mathcal{T}$ and $X\in\mathcal{T}$;
  \item $\mathcal{T}$ is closed under arbitrary unions;
  \item $\mathcal{T}$ is closed under finite intersections.
\end{enumerate}
A map $f\colon(X,\mathcal{T}_X)\to(Y,\mathcal{T}_Y)$ is \emph{continuous} if $f^{-1}(U)\in\mathcal{T}_X$ for every $U\in\mathcal{T}_Y$. A continuous bijection with continuous inverse is a \emph{homeomorphism}, written $X\cong Y$.
\end{definition}

Every metric space $(X,d)$ carries the \emph{metric topology}, generated by the open balls $B(x,r)=\{y\in X:d(x,y)<r\}$.

\begin{definition}[Homotopy and homotopy equivalence]
\label{def:homotopy}
Continuous maps $f,g\colon X\to Y$ are \emph{homotopic} ($f\simeq g$) if there exists a continuous map $H\colon X\times[0,1]\to Y$ with $H(\cdot,0)=f$ and $H(\cdot,1)=g$. A continuous map $f\colon X\to Y$ is a \emph{homotopy equivalence} if there exists $g\colon Y\to X$ with $g\circ f\simeq\mathrm{id}_X$ and $f\circ g\simeq\mathrm{id}_Y$. A space is \emph{contractible} if it is homotopy equivalent to a point.
\end{definition}

Homotopy equivalence is strictly coarser than homeomorphism. Every convex subset of $\R^m$ is contractible (via the straight-line homotopy to any interior point). Homology is a homotopy invariant: $X\simeq Y$ implies $H_k(X;\F)\cong H_k(Y;\F)$ for all $k\geq 0$.

\subsection{Abstract and geometric simplicial complexes}

\begin{definition}[Abstract simplicial complex]
\label{def:simplicial-abstract}
Let $V$ be a finite set. An \emph{abstract simplicial complex} on $V$ is a collection $K\subseteq 2^V\setminus\{\emptyset\}$ closed under non-empty subset: $\sigma\in K$ and $\emptyset\neq\rho\subseteq\sigma$ imply $\rho\in K$. An element $\sigma\in K$ with $|\sigma|=k+1$ is a \emph{$k$-simplex} of dimension $k$. The \emph{$k$-skeleton} of $K$ is $K^{(k)}=\{\sigma\in K:\dim\sigma\leq k\}$.
\end{definition}

\begin{definition}[Standard $k$-simplex]
\label{def:standard-simplex}
The \emph{standard $k$-simplex} is
\[
  \Delta^k=\conv\{e_0,e_1,\ldots,e_k\}
  =\Bigl\{(t_0,\ldots,t_k)\in\R^{k+1}:t_i\geq 0,\;\sum_{i=0}^kt_i=1\Bigr\},
\]
where $\{e_i\}$ is the standard basis of $\R^{k+1}$. Its open interior is $\mathrm{int}(\Delta^k)=\{t\in\Delta^k:t_i>0\;\forall i\}$.
\end{definition}

\begin{definition}[Geometric realisation]
\label{def:geo-real}
Let $K$ be an abstract simplicial complex with vertex set $V$. An injection $\iota\colon V\hookrightarrow\R^N$ is in \emph{general position} for $K$ if no $d+2$ vertices lie in a common affine subspace of dimension $\leq d=\dim K$. The \emph{geometric realisation} of $K$ is
\[
  |K|=\bigcup_{\sigma\in K}\conv\bigl\{\iota(v):v\in\sigma\bigr\}\subset\R^N,
\]
equipped with the subspace topology of $\R^N$.
\end{definition}

The homeomorphism type of $|K|$ does not depend on the choice of general-position embedding. Geometric realisation is \emph{functorial}: a simplicial map $f\colon K\to L$ (a vertex map sending each simplex of $K$ to a simplex of $L$) induces a continuous map $|f|\colon|K|\to|L|$ by barycentric interpolation.

% ═════════════════════════════════════════════════════════════
\section{Simplicial Homology}
\label{subsec:homology}

Fix a field $\F$. The computational implementation uses $\F=\F_2$, the field with two elements; the algebraic framework is identical over any field, modulo orientation conventions.

\subsection{Chains and the boundary operator}

An \emph{oriented $k$-simplex} is an equivalence class $[v_0,\ldots,v_k]$ of ordered tuples of distinct vertices, modulo even permutations: $[v_{\pi(0)},\ldots,v_{\pi(k)}]=\sgn(\pi)\,[v_0,\ldots,v_k]$. Each simplex of positive dimension carries exactly two orientations; over $\F_2$, where $-1=1$, orientation is irrelevant.

\begin{definition}[Chain group]
\label{def:chains}
The \emph{$k$-th chain group} $C_k(K;\F)$ is the $\F$-vector space freely generated by the oriented $k$-simplices of $K$, with the identification $-\sigma=\bar\sigma$ for the two orientations of the same underlying simplex.
\end{definition}

\begin{definition}[Boundary operator]
\label{def:boundary}
The \emph{$k$-th boundary} $\partial_k\colon C_k(K;\F)\to C_{k-1}(K;\F)$ is the $\F$-linear map defined on basis elements by
\begin{equation}
  \partial_k\bigl([v_0,v_1,\ldots,v_k]\bigr)
  =\sum_{i=0}^k(-1)^i[v_0,\ldots,\widehat{v_i},\ldots,v_k],
  \label{eq:boundary}
\end{equation}
where $\widehat{v_i}$ denotes omission of $v_i$.
\end{definition}

\begin{lemma}[Fundamental relation $\partial^2=0$]
\label{lem:boundary-squared}
For all $k\geq 1$, $\partial_{k-1}\circ\partial_k=0$.
\end{lemma}

\begin{proof}
It suffices to verify the identity on a basis element $\sigma=[v_0,\ldots,v_k]$. Two applications of \eqref{eq:boundary} give
\begin{multline*}
  \partial_{k-1}(\partial_k\sigma)
  =\sum_{j<i}(-1)^{i+j}[v_0,\ldots,\widehat{v_j},\ldots,\widehat{v_i},\ldots,v_k]\\
  +\sum_{i<j}(-1)^{i+j-1}[v_0,\ldots,\widehat{v_i},\ldots,\widehat{v_j},\ldots,v_k],
\end{multline*}
where in the second sum the omitted vertex $v_j$ has been shifted by one position because $v_i$ with $i<j$ was already deleted. For each unordered pair $\{i,j\}$ with $i<j$ the two contributions cancel exactly, yielding zero.
\end{proof}

Geometrically, Lemma~\ref{lem:boundary-squared} expresses the fact that the boundary of a simplex is itself a closed surface and hence has no boundary.

\subsection{Homology groups and Betti numbers}

Lemma~\ref{lem:boundary-squared} implies $\im(\partial_{k+1})\subseteq\ker(\partial_k)$. The quotient is well defined.

\begin{definition}[Homology and Betti numbers]
\label{def:homology}
The \emph{cycle group} is $Z_k=\ker\partial_k\subseteq C_k(K;\F)$; the \emph{boundary group} is $B_k=\im\partial_{k+1}\subseteq Z_k$. The \emph{$k$-th homology} is the quotient
\[
  H_k(K;\F)=Z_k/B_k,
\]
and the \emph{$k$-th Betti number} is $\beta_k=\dim_\F H_k(K;\F)$.
\end{definition}

The Betti numbers admit elementary geometric interpretations. $\beta_0$ counts the connected components of $|K|$; $\beta_1$ counts the number of independent one-dimensional cycles (loops not bounding any 2-chain); $\beta_2$ counts the number of enclosed two-dimensional voids.

\begin{theorem}[Euler--Poincar\'e formula]
\label{thm:euler}
Let $f_k=|\{\sigma\in K:\dim\sigma=k\}|$. The Euler characteristic of $K$, defined by
\[
  \chi(K)=\sum_{k\geq 0}(-1)^k f_k,
\]
satisfies
\begin{equation}
  \chi(K)=\sum_{k\geq 0}(-1)^k\beta_k.
  \label{eq:euler}
\end{equation}
\end{theorem}

\begin{proof}
The rank--nullity theorem applied to $\partial_k$ gives $\dim C_k=\dim Z_k+\dim B_{k-1}$ for all $k\geq 0$, with $B_{-1}=0$. Summing with alternating signs and rearranging telescopes the boundary-group terms; the result is~\eqref{eq:euler}.
\end{proof}

The formula links a combinatorial datum ($f_k$) to homotopy invariants ($\beta_k$) and serves as a consistency check on any homology computation.

\subsection{Functoriality}

\begin{definition}[Simplicial and chain maps]
\label{def:simplicial-map}
A \emph{simplicial map} $f\colon K\to L$ is a vertex map $f\colon V(K)\to V(L)$ with $f(\sigma)\in L$ for every $\sigma\in K$. The induced \emph{chain map} $f_\#\colon C_k(K;\F)\to C_k(L;\F)$ is defined on basis elements by
\[
  f_\#([v_0,\ldots,v_k])=
  \begin{cases}
    [f(v_0),\ldots,f(v_k)] & \text{if } f(v_0),\ldots,f(v_k)\text{ are distinct,}\\
    0 & \text{otherwise.}
  \end{cases}
\]
\end{definition}

One checks that $f_\#\partial_k^K=\partial_k^L f_\#$, i.e., chain maps commute with boundaries. Consequently, $f_\#$ maps cycles to cycles and boundaries to boundaries, and descends to a well-defined linear map
\[
  H_k(f)\colon H_k(K;\F)\longrightarrow H_k(L;\F),\quad [z]\longmapsto[f_\#(z)].
\]
Composition is preserved, $H_k(g\circ f)=H_k(g)\circ H_k(f)$, and identities are preserved, $H_k(\mathrm{id})=\mathrm{id}$. In categorical language, $H_k(\,\cdot\,;\F)$ is a functor from the category of simplicial complexes to the category $\Vect_\F$ of $\F$-vector spaces.

This functoriality is the bridge to persistent homology. The inclusions $K_i\hookrightarrow K_j$ for $i\leq j$ in any filtration are simplicial maps, and the induced maps $H_k(K_i;\F)\to H_k(K_j;\F)$ are the structure maps of the persistence module defined below.

% ═════════════════════════════════════════════════════════════
\section{Persistent Homology}
\label{subsec:ph}

\subsection{Filtrations and complex families}
\label{subsubsec:filtration}

\begin{definition}[Filtration]
\label{def:filtration}
A \emph{filtration} of a simplicial complex $K$ is a nested family of subcomplexes
\[
  \emptyset=K_0\subseteq K_1\subseteq\cdots\subseteq K_n=K,
\]
indexed by a non-decreasing real sequence $\varepsilon_0\leq\varepsilon_1\leq\cdots\leq\varepsilon_n$ of \emph{filtration values}.
\end{definition}

For a finite point cloud $P\subset\R^m$ equipped with metric $d_{\mathrm{met}}$, the methodology uses three filtered complex families.

\paragraph{Vietoris--Rips complex.}
\begin{equation}
  \mathcal{R}_\varepsilon(P)=\bigl\{\sigma\subseteq P:\diam_{d_{\mathrm{met}}}(\sigma)\leq\varepsilon\bigr\}.
  \label{eq:rips}
\end{equation}
The Rips complex is a \emph{flag complex}: every clique of its 1-skeleton is filled by a higher-dimensional simplex. The filtration is therefore determined entirely by the order in which edges appear.

\paragraph{Alpha complex.}
For the Euclidean metric, the Voronoi cell of $p\in P$ is
\[
  V_p=\bigl\{x\in\R^m:\|x-p\|_2\leq\|x-p'\|_2\;\forall p'\in P\bigr\}.
\]
The \emph{Alpha complex} is the nerve of the restricted cover $\{B(p,\varepsilon)\cap V_p\}_{p\in P}$:
\begin{equation}
  \mathcal{A}_\varepsilon(P)=\Bigl\{\sigma\subseteq P:\bigcap_{p\in\sigma}\bigl(B(p,\varepsilon)\cap V_p\bigr)\neq\emptyset\Bigr\}.
  \label{eq:alpha}
\end{equation}
$\mathcal{A}_\varepsilon(P)$ is a subcomplex of the Delaunay triangulation, hence has at most $O(|P|^{\lceil m/2\rceil})$ simplices, dramatically less than the worst-case Rips bound of $2^{|P|}$.

\paragraph{Sparse Rips complex.}
The SparseRips construction~\parencite{cavanna2015geometric} replaces $P$ by a greedy \emph{maxmin} landmark subset $L\subseteq P$: starting from an arbitrary seed, each successive landmark is the point of $P$ farthest from the existing landmarks. The persistent homology of the resulting weighted complex is provably a $(1+\varepsilon)$-approximation of that of $\mathcal{R}_\varepsilon(P)$ in the interleaving distance of Definition~\ref{def:interleaving}.

\begin{theorem}[Nerve theorem~\parencite{borsuk1948imbedding}]
\label{thm:nerve}
Let $\mathcal{U}=\{U_i\}_{i\in I}$ be a finite open cover of a paracompact space $X$ such that every non-empty finite intersection $U_{i_0}\cap\cdots\cap U_{i_p}$ is contractible. Then $|N(\mathcal{U})|\simeq X$, where $N(\mathcal{U})$ is the abstract simplicial complex with simplices $\{i_0,\ldots,i_p\}$ for which the corresponding intersection is non-empty.
\end{theorem}

The Alpha complex satisfies the hypothesis of Theorem~\ref{thm:nerve}: each $B(p,\varepsilon)\cap V_p$ is the intersection of a ball and a convex Voronoi cell, hence convex and therefore contractible, and finite intersections of convex subsets of $\R^m$ are convex and contractible. Consequently $|\mathcal{A}_\varepsilon(P)|\simeq\bigcup_{p\in P}B(p,\varepsilon)$: the Alpha complex faithfully encodes the topology of the union of balls without introducing simplices beyond the Delaunay triangulation.

\begin{figure}[!ht]
\centering
\includegraphics[width=\textwidth]{figures/fig_rips_filtration}
\caption{Rips$\,(X,2\varepsilon)$ filtration of a six-point cloud at four scales.
  At $\varepsilon=0$ only vertices are present ($\beta_0=6$, $\beta_1=0$).  As
  $\varepsilon$ increases, edges appear and connected components merge ($\beta_0$
  decreases).  A loop ($H_1$ class) is born when the polygon is complete and dies
  when triangles fill it.  Features far from the diagonal in the persistence diagram
  have long lifetime and are noise-robust (Cohen--Steiner stability,
  Theorem~\ref{thm:stability}).}
\label{fig:rips-filtration}
\end{figure}

\subsection{Persistence modules}
\label{subsubsec:persmods}

Regard the linearly ordered set $(\R,\leq)$ as a category: objects are real numbers, and there is a unique morphism $r\to s$ whenever $r\leq s$. Let $\Vect_\F$ be the category of $\F$-vector spaces and $\F$-linear maps.

\begin{definition}[Persistence module]
\label{def:persmod}
A \emph{(one-parameter) persistence module} over $\F$ is a functor
\[
  \mathbb{V}\colon(\R,\leq)\longrightarrow\Vect_\F.
\]
Concretely, $\mathbb{V}$ assigns a vector space $V_r$ to each $r\in\R$ and a linear map $\rho_{r,s}\colon V_r\to V_s$ (the \emph{structure map}) to each $r\leq s$, with $\rho_{r,r}=\mathrm{id}$ and $\rho_{s,t}\circ\rho_{r,s}=\rho_{r,t}$.
\end{definition}

A morphism $\Phi\colon\mathbb{U}\to\mathbb{V}$ is a natural transformation: linear maps $\Phi_r\colon U_r\to V_r$ with $\Phi_s\circ\rho^\mathbb{U}_{r,s}=\rho^\mathbb{V}_{r,s}\circ\Phi_r$ for all $r\leq s$. Persistence modules with such morphisms form a category $\mathbf{Pers}_\F$.

Applying $H_k(\,\cdot\,;\F)$ levelwise to a filtration and to each inclusion $K_i\hookrightarrow K_j$ yields a persistence module $\mathbb{V}_k=\{H_k(K_i;\F)\}$: the \emph{$k$-th persistence module of the filtration}. This is the central algebraic invariant studied by persistent homology.

\begin{definition}[Interleaving distance]
\label{def:interleaving}
For $\varepsilon\geq 0$, the \emph{$\varepsilon$-shift} of $\mathbb{V}$ is the persistence module $\mathbb{V}[\varepsilon]_r:=V_{r+\varepsilon}$. Modules $\mathbb{U},\mathbb{V}$ are \emph{$\varepsilon$-interleaved} if there exist morphisms $\Phi\colon\mathbb{U}\to\mathbb{V}[\varepsilon]$ and $\Psi\colon\mathbb{V}\to\mathbb{U}[\varepsilon]$ with $\Psi[\varepsilon]\circ\Phi=\rho^\mathbb{U}_{r,r+2\varepsilon}$ and $\Phi[\varepsilon]\circ\Psi=\rho^\mathbb{V}_{r,r+2\varepsilon}$ for all $r$. The \emph{interleaving distance} is
\[
  d_I(\mathbb{U},\mathbb{V})=\inf\bigl\{\varepsilon\geq 0:\mathbb{U}\text{ and }\mathbb{V}\text{ are }\varepsilon\text{-interleaved}\bigr\}.
\]
\end{definition}

The interleaving distance is an extended pseudometric on persistence modules; it descends to a metric on the set of isomorphism classes.

\subsection{The structure theorem and persistence diagrams}
\label{subsubsec:structure}

\begin{definition}[Interval module]
\label{def:interval-mod}
For $b<d$ in $\R\cup\{+\infty\}$, the \emph{interval module} $\F[b,d)$ is the persistence module with $V_r=\F$ for $b\leq r<d$, $V_r=0$ otherwise, and identity structure maps between non-zero positions.
\end{definition}

\begin{theorem}[Structure theorem, Crawley-Boevey~\parencite{crawley2015decomposition}]
\label{thm:structure}
Every pointwise finite-dimensional persistence module $\mathbb{V}$ over a field $\F$ decomposes uniquely, up to isomorphism and reordering, as a direct sum of interval modules:
\begin{equation}
  \mathbb{V}\;\cong\;\bigoplus_{i\in I}\F[b_i,d_i).
  \label{eq:structure}
\end{equation}
The multiset $\{(b_i,d_i)\}_{i\in I}$ is a complete isomorphism invariant of $\mathbb{V}$.
\end{theorem}

The special case for tame modules over filtered finite complexes was established by Zomorodian and Carlsson~\parencite{edelsbrunner2002topological}. Each interval $(b_i,d_i)$ represents a topological feature of degree $k$ that is \emph{born} at scale $b_i$ and \emph{dies} at scale $d_i$; its \emph{persistence} is $p_i=d_i-b_i$. The diagonal is added by convention so that the diagram is a multiset in
\[
  \{(b,d)\in[\,-\infty,\infty\,]^2:b\leq d\}\cup\{\text{diagonal with infinite multiplicity}\}.
\]

\begin{definition}[Persistence diagram]
\label{def:pd}
The \emph{$k$-th persistence diagram} of a filtration is
\begin{equation}
  \PD_k=\bigl\{(b_i,d_i):b_i<d_i<+\infty\bigr\}_{i\in I_k}.
  \label{eq:pd}
\end{equation}
Features with $d_i=+\infty$ are excluded from the present feature vectorisation because their infinite lifetime requires special handling in any scalar summary.
\end{definition}

\subsection{Distances and stability}
\label{subsubsec:pd-distances}

\begin{definition}[Bottleneck and Wasserstein distances]
\label{def:bottleneck}
A \emph{partial matching} between diagrams $D_1,D_2$ is a bijection $\gamma\colon D_1\cup\Delta\to D_2\cup\Delta$ that is the identity outside a finite set, where $\Delta=\{(x,x):x\in\R\}$ is the diagonal. The \emph{$p$-Wasserstein distance} ($1\leq p<\infty$) is
\[
  W_p(D_1,D_2)=\inf_\gamma\Bigl(\sum_{x\in D_1\cup\Delta}\|x-\gamma(x)\|_\infty^p\Bigr)^{1/p},
\]
and the \emph{bottleneck distance} is
\[
  d_B(D_1,D_2)=\inf_\gamma\sup_{x\in D_1\cup\Delta}\|x-\gamma(x)\|_\infty=\lim_{p\to\infty}W_p(D_1,D_2).
\]
\end{definition}

\begin{theorem}[Stability of persistence diagrams, Cohen-Steiner--Edelsbrunner--Harer~\parencite{cohen2007stability}]
\label{thm:stability}
Let $f,g\colon X\to\R$ be tame continuous functions on a triangulable compact space $X$, and let $D_f,D_g$ be the persistence diagrams of their sublevel-set filtrations. Then
\[
  d_B(D_f,D_g)\leq\|f-g\|_\infty.
\]
\end{theorem}

\begin{theorem}[Algebraic stability, Chazal et al.~\parencite{chazal2016structure}]
\label{thm:isometry}
For any two pointwise finite-dimensional persistence modules $\mathbb{U},\mathbb{V}$ over a field $\F$,
\[
  d_B\bigl(\PD(\mathbb{U}),\PD(\mathbb{V})\bigr)=d_I(\mathbb{U},\mathbb{V}).
\]
\end{theorem}

Theorem~\ref{thm:isometry} upgrades stability from an inequality to an isometry: the bottleneck distance on diagrams equals the interleaving distance on modules, providing a complete algebraic characterisation of the geometry of persistence.

For point cloud filtrations, the metric stability statement of practical interest is the following.

\begin{corollary}[Metric stability]
\label{cor:metric-stability}
Let $P,Q\subset\R^m$ be finite point sets and let
\[
  d_H(P,Q)=\max\Bigl(\max_{p\in P}\min_{q\in Q}\|p-q\|_2,\,\max_{q\in Q}\min_{p\in P}\|p-q\|_2\Bigr)
\]
be their Hausdorff distance. Then for every $k\geq 0$,
\[
  d_B\bigl(\PD_k(\mathcal{R}(P)),\,\PD_k(\mathcal{R}(Q))\bigr)\leq 2\,d_H(P,Q).
\]
\end{corollary}

\begin{proof}[Proof sketch]
If $d_H(P,Q)\leq\delta$, then each simplex of $\mathcal{R}_\varepsilon(P)$ lifts to a simplex of $\mathcal{R}_{\varepsilon+2\delta}(Q)$ by replacing each vertex by a point of $Q$ within distance $\delta$ and applying the triangle inequality to the diameter. By symmetry the two Rips modules are $2\delta$-interleaved, hence $d_B\leq 2\delta$ by Theorem~\ref{thm:isometry}.
\end{proof}

The practical content is that bounded sensor noise of amplitude $\delta$ perturbs each embedded point by at most $\delta$ in $\ell_2$, hence the persistence diagrams of two realisations of the same fall event differ by at most $2\delta$ in bottleneck distance. Features with persistence $p_i\gg 2\delta$ are noise-robust; near-diagonal features with $p_i\approx 0$ are noise-sensitive.

\begin{figure}[!ht]
\centering
\includegraphics[width=0.85\textwidth]{figures/fig_persistence_diagram}
\caption{Persistence diagrams ($H_0$ circles, $H_1$ squares) for delay-embedded ADL
  and fall windows.  Points far from the diagonal represent long-lived topological
  features.  The fall window typically exhibits a high-persistence $H_1$ feature
  (loop) absent or short-lived in the ADL window---a geometric signature of the
  sharp impact trajectory in the delay-embedded space.  Triangles and diamonds
  denote essential classes (infinite death, plotted at the top margin).}
\label{fig:persistence-diagram}
\end{figure}

\subsection{Computational reduction}
\label{subsubsec:reduction}

Persistence diagrams are computed by \emph{boundary matrix reduction}. Order the simplices of the filtration as $\sigma_1,\sigma_2,\ldots,\sigma_n$ so that $\dim\sigma_i\leq\dim\sigma_{i+1}$ and simplices entering at smaller scale precede those at larger scale (consistent with filtration order). The boundary matrix $D\in\F^{n\times n}$ has $D_{i,j}=1$ if $\sigma_i\in\partial\sigma_j$ and $0$ otherwise. The \emph{pivot} of column $j$ is $\low(D,j)=\max\{i:D_{i,j}\neq 0\}$, undefined if column $j$ is zero.

\begin{algorithm}[H]
\caption{Standard persistence algorithm (column reduction over $\F_2$).}
\label{alg:reduction}
\begin{algorithmic}[1]
\Require Boundary matrix $D\in\{0,1\}^{n\times n}$ in filtration order
\Ensure Reduced matrix $R$ with unique pivots; persistence pairs read from $R$
\State $R\leftarrow D$
\For{$j=1$ \textbf{to} $n$}
  \While{$\exists\;i<j$ with $\low(R,i)=\low(R,j)\neq\bot$}
    \State $R[\,\cdot\,,j]\leftarrow R[\,\cdot\,,j]+R[\,\cdot\,,i]$ \Comment{column addition over $\F_2$}
  \EndWhile
\EndFor
\State \textbf{Pairs:} $\{(\low(R,j),\,j):R[\,\cdot\,,j]\neq 0\}$
\end{algorithmic}
\end{algorithm}

\begin{proposition}[Uniqueness of pairing]
\label{prop:reduction-unique}
The pairing produced by Algorithm~\ref{alg:reduction} is independent of the order in which column additions are performed; it depends only on the filtration.
\end{proposition}

The proposition is a corollary of the structure theorem~\ref{thm:structure}. The worst-case complexity of Algorithm~\ref{alg:reduction} is $O(n^3)$; the GUDHI library~\parencite{gudhi} implements optimised variants (twist reduction, clearing, dualisation to cohomology) with near-linear average-case complexity on the point cloud sizes considered here.

% ═════════════════════════════════════════════════════════════
\section{Feature Vectorisation}
\label{subsec:vectorization}

Persistence diagrams are complete isomorphism invariants of the filtration but cannot be passed directly to standard classifiers, which require fixed-length input vectors. Two obstacles obstruct direct use:
\begin{enumerate}[label=(\roman*)]
  \item the cardinality $|I_k|$ varies across windows and configurations;
  \item the diagram space, under the bottleneck or Wasserstein metric, is not a Hilbert space.
\end{enumerate}
A variety of vectorisation schemes circumvent these obstacles: persistence images~\parencite{adams2017persistence}, persistence landscapes~\parencite{bubenik2015statistical}, and others. The present methodology uses a fixed set of scalar statistics, motivated by the requirement that the feature dimension be constant across $\lambda$ to permit a uniform Bayesian search.

\subsection{Persistence-based statistics}
\label{subsubsec:stats-topo}

For each homological degree $k\in\{0,1\}$, the persistence diagram $\PD_k$ yields a finite multiset of birth-death pairs. A fixed set of scalar statistics is computed from these pairs for each degree, motivated by the requirement that the feature dimension be constant across all hyperparameter configurations $\lambda$ to permit a uniform classifier interface. The specific statistics chosen and their definitions are listed in Section~\ref{subsec:features} of the Experiments chapter.

\subsection{Signal-level features}
\label{subsubsec:signal-feats}

Four additional scalar features are computed directly from $\mathrm{SMV}_{\mathrm{acc}}$ on the window:
\begin{equation}
  f_{21}=\max_t\mathrm{SMV},\quad
  f_{22}=\mathrm{std}(\mathrm{SMV}),\quad
  f_{23}=\overline{\mathrm{SMV}},\quad
  f_{24}=\max_t\mathrm{SMV}-\min_t\mathrm{SMV}.
  \label{eq:sig-feats}
\end{equation}

\subsection{Complete feature vector}
\label{subsubsec:full-vec}

The complete feature vector under configuration $\lambda$ is
\begin{equation}
  \phi_\lambda(w)=\bigl(f_1^{(0)},\ldots,f_{10}^{(0)},\,f_1^{(1)},\ldots,f_{10}^{(1)},\,f_{21},f_{22},f_{23},f_{24}\bigr)\in\R^{24}.
  \label{eq:full-vec}
\end{equation}
The fixed dimension is a methodological constraint imposed by the requirement of a uniform classifier interface across $\lambda$.

The point cloud has now been converted into a fixed-dimensional vector $\phi_\lambda(w)\in\R^{24}$ carrying a stability certificate inherited from Theorem~\ref{thm:stability}. What remains is to turn this vector into a fall/no-fall decision, which requires stating---only now that every ingredient exists---what "solving fall detection" formally means.

% ═════════════════════════════════════════════════════════════
\section{Problem Redefinition: The Classification Problem}
\label{subsec:formal}

The preceding sections built a representation map $w\mapsto\phi_\lambda(w)\in\R^{24}$ out of raw signal processing (Section~\ref{subsec:signal_rep}) and topological data analysis (Sections~\ref{subsec:embedding}--\ref{subsec:vectorization}), without reference to labels, subjects, or a classifier. This section redefines the original problem---detect a fall---as a formal classification problem stated over $\phi_\lambda$: what the hyperparameter tuple $\lambda$ collects, what a classifier is required to do, and which metric judges success.

Let $(\X,\Y,\Sset)$ be a random triple drawn from an unknown joint distribution $\Prob$, where $\X$ is a vector-valued inertial recording, $\Y\in\{0,1\}$ the binary fall label ($\Y=1$ denotes a fall), and $\Sset$ takes values in a finite set of subject identifiers. The marginal class prior $\pi:=\Prob(\Y=1)$ is strictly less than $1/2$, reflecting the empirical fact that falls are rare minority events.

\subsection{Performance metric selection}
\label{subsubsec:metric-selection}

Before defining the classifier, it is necessary to fix what "success" means, since the choice of performance metric is not a presentational detail: it is a design parameter that shapes every step of the pipeline, from the Bayesian objective to the decision threshold. This subsection surveys the standard binary classification metrics, identifies the structural limitations of each for the fall-detection setting, and derives the $F_2$ score as the appropriate choice.

\paragraph{Confusion matrix and derived rates.}
For a binary classifier producing predictions $\widehat y_i\in\{0,1\}$ against labels $y_i\in\{0,1\}$ on a test set of size $N$, define the four confusion counts:
\[
  \mathrm{TP}=|\{i:y_i=1,\widehat y_i=1\}|,\quad
  \mathrm{FP}=|\{i:y_i=0,\widehat y_i=1\}|,\quad
  \mathrm{FN}=|\{i:y_i=1,\widehat y_i=0\}|,\quad
  \mathrm{TN}=|\{i:y_i=0,\widehat y_i=0\}|.
\]
The four primary derived rates are:
\begin{align*}
  \Acc  &= \frac{\mathrm{TP}+\mathrm{TN}}{N}, &
  \Rec    &= \frac{\mathrm{TP}}{\mathrm{TP}+\mathrm{FN}}, \\
  \Prec &= \frac{\mathrm{TP}}{\mathrm{TP}+\mathrm{FP}}, &
  F_1                &= \frac{2\,\Prec\cdot\Rec}{\Prec+\Rec}.
\end{align*}
All four are reported in result tables for completeness.

\paragraph{Why accuracy fails under class imbalance.}
Falls are rare: in the three datasets used here the class prior $\pi=\Prob(\Y=1)$ ranges from approximately $0.10$ to $0.30$. A trivial classifier that predicts $\widehat y=0$ for every window achieves accuracy $1-\pi \geq 0.70$ while detecting zero falls. Accuracy cannot distinguish this degenerate solution from a genuinely useful detector; it is therefore an inappropriate objective whenever $\pi \ll 1/2$.

\paragraph{Why recall alone is insufficient.}
Recall measures the fraction of true falls detected, so maximising recall directly incentivises predicting $\widehat y=1$ for every window---a classifier with $\mathrm{Recall}=1$ and arbitrarily many false alarms. In a deployed system, chronic false alarms erode user trust and lead to the device being removed, defeating the purpose of detection. A useful metric must penalise false alarms to some degree.

\paragraph{Why $F_1$ does not weight recall sufficiently.}
The $F_1$ score is the harmonic mean of precision and recall, treating them symmetrically. Under $F_1$, a unit gain in recall and a unit gain in precision contribute equally to the score. For fall detection this symmetry is inappropriate: a missed fall exposes a vulnerable person to the long-lie complications described in the Introduction, while a false alarm costs at most an unnecessary check. The asymmetry in clinical consequence calls for a metric that values recall more heavily than precision.

\paragraph{The $F_\beta$ family and the choice $\beta=2$.}
The $F_\beta$ score generalises $F_1$ by introducing a weight $\beta>0$ that controls the relative importance of recall:
\begin{equation}
  F_\beta = (1+\beta^2)\,\frac{\Prec\cdot\Rec}{\beta^2\,\Prec+\Rec}.
  \label{eq:Fbeta}
\end{equation}
At $\beta=1$ the metric reduces to $F_1$. For $\beta>1$ recall is weighted more heavily; for $\beta<1$ precision dominates. The primary objective metric of this thesis is the $F_\beta$-score at $\beta=2$:
\begin{equation}
  F_2=\frac{5\,\Prec\cdot\Rec}{4\,\Prec+\Rec}.
  \label{eq:F2}
\end{equation}

\begin{proposition}[Recall weighting under $F_\beta$]
\label{prop:F2-recall-weight}
For fixed precision $\Prec=p$ and recall $\Rec=r$,
\[
  \frac{\partial F_\beta}{\partial r}\bigg|_{\beta=2}\;=\;\frac{20\,p^2}{(4p+r)^2}.
  \qquad
  \frac{\partial F_\beta}{\partial p}\bigg|_{\beta=2}\;=\;\frac{5\,r^2}{(4p+r)^2}.
\]
Consequently, near $p=r$ the marginal $F_2$-gain per unit of recall is four times the marginal gain per unit of precision.
\end{proposition}

\begin{proof}
Direct differentiation of~\eqref{eq:F2} and the substitution $p=r$ give the ratio $4$ between the two partial derivatives.
\end{proof}

This four-to-one ratio is a deliberate design choice reflecting the clinical cost asymmetry: missing a fall is judged approximately four times as costly as issuing a false alarm. The $F_2$ score is used both as the Bayesian objective and as the headline metric in all result tables.

\begin{remark}[Non-triviality of the baseline]
\label{rem:f2-baseline}
A trivial all-positive classifier achieves $\Rec=1$ and $\Prec=\pi$, hence $F_2=5\pi/(4\pi+1)$. When $\pi<0.30$ this gives $F_2<0.56$, a readily beaten bound. The recall-oriented objective does not degenerate to a trivial solution.
\end{remark}

\begin{remark}[Metric as a pipeline-level commitment]
\label{rem:metric-commitment}
The choice of $F_2$ is not applied only at evaluation time. It propagates into every stage of the pipeline: the Bayesian search maximises $F_2$ in the inner cross-validation, the decision threshold $\hat\tau_0$ is selected to maximise $F_2$ on the training fold, and the subject-level personalisation sweep searches for the $F_2$-optimal threshold per subject. A consistent metric across all stages ensures that the pipeline is trained, selected, and evaluated under a single coherent objective.
\end{remark}

\subsection{The hyperparameter tuple}
\label{subsubsec:lambda}

The representation map $\phi_\lambda$ built in Sections~\ref{subsec:signal_rep}--\ref{subsec:vectorization} depends throughout on a tuple of open parameters
\begin{equation}
  \lambda=\bigl(w_{\mathrm{sec}},\;m,\;\tau,\;r,\;c,\;d_{\mathrm{met}},\;\varepsilon,\;q\bigr)\in\Lambda.
  \label{eq:lambda}
\end{equation}
Each component was introduced individually as it first appeared; this subsection collects them into the single tuple that Section~\ref{subsec:bayesopt} optimises jointly. The parameters fall into three conceptually distinct groups according to the stage of the pipeline they control.

\paragraph{Windowing parameter.} The scalar $w_{\mathrm{sec}}>0$ (Section~\ref{subsubsec:windows}) sets the duration of the analysis window in seconds. It determines how much temporal context is included in each classification unit. A short window captures only the impact transient; a long window may include pre-fall and post-fall motion.

\paragraph{Embedding parameters (geometrically meaningful).} The triple $(m,\tau,r)$ of Definition~\ref{def:delay} controls the delay embedding that lifts the scalar window to a point cloud in $\R^m$.
\begin{itemize}[leftmargin=1.5em]
  \item $m\geq 2$ is the \emph{embedding dimension}: the dimension of the reconstructed phase space. By the Takens theorem (Theorem~\ref{thm:takens}), $m$ must exceed twice the intrinsic dimension of the underlying attractor for the embedding to be injective.
  \item $\tau\geq 1$ is the \emph{lag}: the number of time steps separating successive coordinates of each embedded point. It controls the time scale at which the geometry of the trajectory is unfolded.
  \item $r\geq 1$ is the \emph{stride}: the step between successive embedded points. It subsamples the trajectory and directly controls the size of the point cloud, hence the cost of the subsequent homology computation.
\end{itemize}

\paragraph{Filtration parameters.} The quadruple $(c,d_{\mathrm{met}},\varepsilon,q)$ of Section~\ref{subsubsec:filtration} governs the construction of the filtered simplicial complex from the point cloud.
\begin{itemize}[leftmargin=1.5em]
  \item $c$ selects the \emph{complex family}: the rule by which simplices are added as the filtration scale grows. Different families capture different aspects of the point cloud geometry at different computational costs.
  \item $d_{\mathrm{met}}$ is the \emph{metric} on $\R^m$ used to compute inter-point distances.
  \item $\varepsilon\in(0,1)$ is a \emph{filtration-scale percentile}: the filtration is truncated at the $\varepsilon$-quantile of the pairwise distance distribution, bounding the complex size.
  \item $q$ is a \emph{landmark fraction} used by approximate complex constructions to subsample the point cloud before building the complex.
\end{itemize}

\begin{remark}[Geometric vs.\ algorithmic parameters]
\label{rem:param-roles}
The embedding parameters $(m,\tau,r)$ have direct geometric interpretations: they determine the shape and resolution of the reconstructed phase-space trajectory. The filtration parameters $(c,d_{\mathrm{met}},\varepsilon,q)$ are algorithmic choices that affect how the geometry of the point cloud is discretised and at what computational cost. Both groups interact: the optimal filtration parameters depend on the geometry of the point cloud, which depends on the embedding. This coupling is why $\lambda$ is optimised jointly rather than sequentially. The specific candidate values for each parameter and the search procedure are given in Section~\ref{subsubsec:bayes-search-space}.
\end{remark}

\subsection{The representation problem}
\label{subsubsec:repr-problem}

Restated formally, the representation map built in the preceding sections had to map a variable-length scalar window $w\in\R^L$ to a fixed-dimensional vector that a standard classifier can consume, subject to three constraints. It is worth verifying that $\phi_\lambda$ indeed satisfies all three, since a representation failing any one of them would be unusable regardless of classifier quality.

\textbf{Fixed dimension.} The window length $L=\lfloor w_{\mathrm{sec}}\cdot f_s\rfloor$ varies with the hyperparameter $w_{\mathrm{sec}}$ and the dataset's sampling rate $f_s$. A representation that depends on $L$---such as a spectral feature vector or a raw sample sequence---cannot be used directly when comparing configurations with different $w_{\mathrm{sec}}$, or when the same model must be applied across datasets recorded at different rates. Section~\ref{subsec:vectorization} confirms that $\phi_\lambda(w)\in\R^{24}$ regardless of $L$: the dimension depends only on the number of homological degrees considered ($k=0,1$) and the number of statistics per degree (ten each, plus four signal-level features), not on $w_{\mathrm{sec}}$, $f_s$, $L$, or the number of simplices in the filtration.

\textbf{Geometric faithfulness.} The representation should capture the information that discriminates falls from ADL windows. As argued in Section~\ref{subsec:embedding}, this information is primarily geometric: falls and ADL windows produce point clouds of different topological shape in the delay-embedded space. A representation built from scalar statistics of the raw time series---means, variances, spectral energies---does not directly encode this geometric structure and has no principled connection to the shape of the trajectory. $\phi_\lambda$ is built instead from persistence diagrams (Sections~\ref{subsec:topological_prelim}--\ref{subsec:ph}), which are a complete invariant of that shape.

\textbf{Stability.} A representation is stable if small perturbations of the input window produce small perturbations of the output vector. Without stability, sensor noise, quantisation artefacts, or minor changes in signal conditioning can produce arbitrarily large feature changes, making the classifier unreliable across devices and sessions. Stability is guaranteed by the Cohen--Steiner theorem (Theorem~\ref{thm:stability}): the bottleneck distance between two persistence diagrams is bounded by the supremum norm of the perturbation between the corresponding input functions. No fixed handcrafted feature family possesses a comparable mathematical certificate.

\begin{remark}[The representation is the thesis contribution, not the classifier]
\label{rem:repr-is-contribution}
The map $\phi_\lambda$ is the object of study. It is parameterised by $\lambda\in\Lambda$, and the central empirical question is: what is the best $\lambda^*$ for a given dataset, and does the optimal representation generalise across subjects? The classifier $f_\theta$ is deliberately kept simple---logistic regression and linear SVM---so that any observed performance difference can be attributed to the representation, not to classifier complexity. This design choice separates the two contributions cleanly.
\end{remark}

\subsection{The classifier}
\label{subsubsec:classifier-formal}

A probabilistic classifier is a map
\[
  f_\theta\colon\R^{24}\longrightarrow[0,1],
\]
where $f_\theta(\phi_\lambda(w))$ is the estimated posterior fall probability $\widehat\Prob(\Y=1\mid w;\theta,\lambda)$. The model class is restricted to two linear families:
\[
  \Fcal_{\mathrm{lr}}=\{\text{balanced logistic regression with } \ell_1\text{ or } \ell_2\text{ penalty}\},
\]
\[
  \Fcal_{\mathrm{svm}}=\{\text{balanced } C\text{-SVM with Platt-calibrated probabilities}\}.
\]
Both families produce calibrated probability outputs, which is a prerequisite for principled threshold selection. The mathematical development of each family---objective functions, duality, Platt calibration---is given in Section~\ref{subsec:classifiers}.

For a threshold $\tau_0\in(0,1)$ the \emph{subject-general decision rule} is
\begin{equation}
  \widehat y(w)=\1\bigl\{f_\theta(\phi_\lambda(w))\geq\tau_0\bigr\}.
  \label{eq:dec-general}
\end{equation}
For subject-personalised evaluation a subject-specific threshold $\tau_s$ replaces $\tau_0$:
\begin{equation}
  \widehat y_s(w)=\1\bigl\{f_\theta(\phi_\lambda(w))\geq\tau_s\bigr\}.
  \label{eq:dec-personal}
\end{equation}
The trained threshold $\hat\tau_0$ is the $F_2$-optimal value on the training fold (Section~\ref{subsubsec:threshold}). The personalised threshold $\tau_s$ must be estimated from data disjoint from the test segment of subject $s$ to avoid leakage.

% ═════════════════════════════════════════════════════════════
\section{Classification Models}
\label{subsec:classifiers}

Two linear classifier families are considered, $\Fcal_{\mathrm{lr}}$ and $\Fcal_{\mathrm{svm}}$. Both are applied with balanced class weighting and, for SVM, with Platt probability calibration.

\subsection{Logistic regression}
\label{subsubsec:lr}

Logistic regression models the conditional fall probability by
\begin{equation}
  \Prob(Y=1\mid\mathbf{z};\boldsymbol\theta,b)=\sigma(\boldsymbol\theta^\top\mathbf{z}+b),
  \qquad
  \sigma(u)=\frac{1}{1+e^{-u}},
  \label{eq:lr-sigmoid}
\end{equation}
where $\mathbf{z}\in\R^{24}$ is the standardised feature vector and $(\boldsymbol\theta,b)\in\R^{24}\times\R$ are the parameters. Given training data $\{(\mathbf{z}_i,y_i)\}_{i=1}^n$ the regularised negative log-likelihood is
\begin{equation}
  \mathcal{L}_\Omega(\boldsymbol\theta,b)
  =-\frac{1}{n}\sum_{i=1}^n\Bigl[y_i\log\sigma_i+(1-y_i)\log(1-\sigma_i)\Bigr]
  +\frac{1}{C}\,\Omega(\boldsymbol\theta),
  \label{eq:lr-loss}
\end{equation}
where $\sigma_i=\sigma(\boldsymbol\theta^\top\mathbf{z}_i+b)$, $C>0$ is the inverse regularisation strength, and $\Omega(\boldsymbol\theta)\in\{\tfrac12\|\boldsymbol\theta\|_2^2,\,\|\boldsymbol\theta\|_1\}$ selects the $\ell_2$ or $\ell_1$ penalty. The $\ell_2$ objective is strictly convex and smooth and is solved by L-BFGS or coordinate descent; the $\ell_1$ objective is convex and non-smooth and is solved by liblinear coordinate descent in the implementation.

\begin{proposition}[Convexity and uniqueness for $\ell_2$]
\label{prop:lr-convex}
With $\Omega(\boldsymbol\theta)=\tfrac12\|\boldsymbol\theta\|_2^2$, the objective~\eqref{eq:lr-loss} is strictly convex in $(\boldsymbol\theta,b)$ on $\R^{25}$, hence admits a unique global minimiser.
\end{proposition}

\begin{proof}
The negative log-likelihood is convex because $-\log\sigma(u)$ is convex in $u$ and the maps $u\mapsto\boldsymbol\theta^\top\mathbf{z}+b$ are affine. The $\ell_2$ penalty is strictly convex. The sum of a convex and a strictly convex function is strictly convex.
\end{proof}

The predicted probability is $f_\theta(\mathbf{z})=\sigma(\boldsymbol\theta^\top\mathbf{z}+b)\in(0,1)$.

\subsection{Support vector machine}
\label{subsubsec:svm}

The $C$-support vector classifier seeks a hyperplane $\{\boldsymbol{w}^\top\mathbf{z}+b=0\}$ with maximum geometric margin. With slack variables $\xi_i\geq 0$ the primal soft-margin problem is
\begin{equation}
  \min_{\boldsymbol w,b,\boldsymbol\xi}\;
  \tfrac12\|\boldsymbol w\|_2^2+C\sum_{i=1}^n\xi_i
  \quad\text{s.t.}\quad
  y_i(\boldsymbol w^\top\mathbf{z}_i+b)\geq 1-\xi_i,\;
  \xi_i\geq 0,\;
  y_i\in\{-1,+1\}.
  \label{eq:svm-primal}
\end{equation}
The Lagrangian is
\[
  L(\boldsymbol w,b,\boldsymbol\xi,\boldsymbol\alpha,\boldsymbol\mu)
  =\tfrac12\|\boldsymbol w\|^2+C\sum_i\xi_i
  -\sum_i\alpha_i\bigl[y_i(\boldsymbol w^\top\mathbf{z}_i+b)-1+\xi_i\bigr]
  -\sum_i\mu_i\xi_i,
\]
with $\alpha_i,\mu_i\geq 0$. The KKT stationarity conditions yield
\[
  \nabla_{\boldsymbol w}L=0\;\Rightarrow\;\boldsymbol w=\sum_i\alpha_iy_i\mathbf{z}_i,
  \quad
  \partial_bL=0\;\Rightarrow\;\sum_i\alpha_iy_i=0,
  \quad
  \partial_{\xi_i}L=0\;\Rightarrow\;\alpha_i+\mu_i=C.
\]
Substitution into $L$ produces the dual quadratic programme
\begin{equation}
  \max_{\boldsymbol\alpha}\;\sum_i\alpha_i-\tfrac12\sum_{i,j}\alpha_i\alpha_jy_iy_j\mathbf{z}_i^\top\mathbf{z}_j
  \quad\text{s.t.}\quad
  0\leq\alpha_i\leq C,\;\sum_i\alpha_iy_i=0.
  \label{eq:svm-dual}
\end{equation}
The primal solution recovers $\boldsymbol w^*=\sum_i\alpha_i^*y_i\mathbf{z}_i$ and the decision value for a new point is
\begin{equation}
  f(\mathbf{z})=\sum_i\alpha_i^*y_i\mathbf{z}_i^\top\mathbf{z}+b^*.
  \label{eq:svm-decision}
\end{equation}
Points with $\alpha_i^*>0$ are the \emph{support vectors}. With $\kappa(\mathbf z,\mathbf z')=\mathbf z^\top\mathbf z'$ replaced by the RBF kernel
\[
  \kappa_{\mathrm{rbf}}(\mathbf z,\mathbf z')=\exp\bigl(-\gamma\|\mathbf z-\mathbf z'\|_2^2\bigr),
\]
the same dual applies in a reproducing kernel Hilbert space.

\subsection{Platt score calibration}
\label{subsubsec:platt}

The SVM decision value $f(\mathbf z)\in\R$ is not a probability. Platt scaling~\parencite{platt1999probabilistic} fits a logistic function to the decision values on a held-out calibration set:
\begin{equation}
  \Prob(Y=1\mid\mathbf{z})=\sigma\bigl(Af(\mathbf{z})+B\bigr),
  \qquad A,B\in\R,
  \label{eq:platt}
\end{equation}
with $(A,B)$ estimated by minimising the cross-entropy loss against Laplace-smoothed targets
\[
  t_i=\begin{cases}\dfrac{n_++1}{n_++2}&\text{if }y_i=+1,\\[4pt]\dfrac{1}{n_-+2}&\text{if }y_i=-1,\end{cases}
\]
where $n_\pm$ are the counts of positive and negative calibration examples. The smoothing prevents degenerate $0$ or $1$ predictions when the margin is large.

\subsection{Class-weighted empirical risk}
\label{subsubsec:classweight}

Falls are minority events; with $n_0$ negative and $n_1$ positive training examples ($n=n_0+n_1$) the \emph{balanced class weights} are
\begin{equation}
  w_c=\frac{n}{2\,n_c},\qquad c\in\{0,1\}.
  \label{eq:cw}
\end{equation}
This satisfies $w_0n_0=w_1n_1=n/2$. For logistic regression, the weighted log-likelihood replaces $\mathcal{L}_\Omega$ by
\[
  \mathcal{L}_\Omega^{\mathrm{w}}(\boldsymbol\theta,b)
  =-\frac{1}{n}\sum_{i=1}^n w_{y_i}\Bigl[y_i\log\sigma_i+(1-y_i)\log(1-\sigma_i)\Bigr]
  +\frac{1}{C}\Omega(\boldsymbol\theta).
\]
For C-SVM, the scalar $C$ in~\eqref{eq:svm-primal} is replaced by $C_+=Cw_1$ for positive examples and $C_-=Cw_0$ for negative examples, biasing the model toward higher recall.

\subsection{Decision threshold}
\label{subsubsec:threshold}

A probabilistic classifier $f_\theta\colon\R^{24}\to[0,1]$ assigns each window a posterior fall probability. A binary prediction is produced by thresholding at a trained value $\hat\tau_0$:
\begin{equation}
  \widehat y(w)=\1\bigl\{f_\theta(\phi_\lambda(w))\geq\hat\tau_0\bigr\}.
  \label{eq:dec-thresh}
\end{equation}
The trained threshold $\hat\tau_0$ is the $F_2$-optimal value on the training fold:
\[
  \hat\tau_0 = \argmax_{\tau\in(0,1)}\,F_2^{\mathrm{train}}(\tau).
\]
This PR-curve selection is performed once per cross-validation fold using the training-fold probability scores and labels. The resulting threshold is typically more conservative than the nominal $\tau_0=0.50$: the class-imbalanced training distribution places the optimal precision--recall trade-off above the $50\,\%$ probability level, so that high-probability predictions are required before a fall is declared. The trained threshold is used for the subject-general decision rule; subject-level personalisation replaces $\hat\tau_0$ with a subject-specific threshold $\tau_s$ as defined in Section~\ref{subsec:formal}.

% ═════════════════════════════════════════════════════════════
\section{Bayesian Hyperparameter Optimisation}
\label{subsec:bayesopt}

The representation map $\phi_\lambda$ depends on the hyperparameter tuple $\lambda\in\Lambda$ of~\eqref{eq:lambda}, and the classifier depends on its own tuple $\theta\in\Theta$ (regularisation strength, penalty, kernel). Joint optimisation over $\Lambda\times\Theta$ is necessary because the two tuples interact: the geometry of the topological representation determines which classifier hyperparameters are optimal. The optimisation is performed by Bayesian search with a tree-structured Parzen estimator (TPE), implemented in Optuna~\parencite{optuna}.

\subsection{The optimisation problem}
\label{subsubsec:opt-problem}

Let $\mathcal{D}^{\mathrm{train}}_K$ be the labelled feature corpus of dataset $K$ from which the test subject has been removed (LOSO outer split, Section~\ref{subsubsec:protocols}). Let
\[
  \mathrm{CV}_K\colon\Lambda\times\Theta\longrightarrow\R,
  \qquad
  (\lambda,\theta)\longmapsto\widehat{F}_{2,K}(\lambda,\theta)
\]
denote the inner cross-validation estimator of $\E[F_2\mid K,\lambda,\theta]$ obtained by stratified $k$-fold CV on $\mathcal{D}^{\mathrm{train}}_K$. The optimisation problem solved at each outer fold is
\begin{equation}
  \lambda^*_K,\theta^*_K\;\in\;\argmax_{(\lambda,\theta)\in\Lambda\times\Theta}\mathrm{CV}_K(\lambda,\theta).
  \label{eq:bayes-objective}
\end{equation}
The objective is non-differentiable, noisy (because CV estimates fluctuate with the random fold partition and stochastic class-balancing), and expensive to evaluate (each call requires constructing point clouds, computing persistent homology, fitting the classifier, and performing the CV averaging). These properties make grid search infeasible and motivate a sequential model-based approach.

\subsection{Sequential model-based optimisation}
\label{subsubsec:smbo}

Sequential model-based optimisation (SMBO)~\parencite{hutter2011sequential} alternates between fitting a probabilistic surrogate to past evaluations and proposing new evaluations that maximise an acquisition function under the surrogate. Let $y(\lambda,\theta)=\mathrm{CV}_K(\lambda,\theta)+\eta$ denote one noisy evaluation of the objective, with $\eta$ centred. After $t$ evaluations the data set is
\[
  \mathcal{H}_t=\bigl\{((\lambda_i,\theta_i),y_i):i=1,\ldots,t\bigr\}.
\]
A surrogate $p(y\mid\lambda,\theta,\mathcal{H}_t)$ is a probability model over the objective conditional on the search point and the history. The standard acquisition function is the \emph{expected improvement} (EI) over a reference value $y^*$:
\begin{equation}
  \EI_{y^*}(\lambda,\theta)=\E\bigl[\max(y^*-y,0)\bigm|\lambda,\theta,\mathcal{H}_t\bigr].
  \label{eq:ei}
\end{equation}
(Here EI is stated for a minimisation problem with $y$ the negative $\widehat F_2$; sign conventions are equivalent under negation.) The next search point is
\begin{equation}
  (\lambda_{t+1},\theta_{t+1})\in\argmax_{(\lambda,\theta)\in\Lambda\times\Theta}\EI_{y^*}(\lambda,\theta).
  \label{eq:smbo-step}
\end{equation}
The surrogate is updated and the loop continues for a fixed budget. The textbook surrogate is a Gaussian process posterior; TPE replaces it with a non-parametric density model that is better suited to discrete and conditional search spaces.

\subsection{Tree-structured Parzen estimator}
\label{subsubsec:tpe}

TPE~\parencite{bergstra2011algorithms} models the posterior in an inverted form. Instead of modelling $p(y\mid\lambda,\theta)$ directly, it splits the observed history at a quantile $\gamma\in(0,1)$ and models the density of the inputs given two quality regimes,
\[
  p(\lambda,\theta\mid y\leq y^*)=\ell(\lambda,\theta),
  \qquad
  p(\lambda,\theta\mid y>y^*)=g(\lambda,\theta),
\]
where $y^*$ is the empirical $\gamma$-quantile of the observed $\{y_i\}$. The two density estimates are constructed from the corresponding subsets of $\mathcal{H}_t$ by kernel density estimation: each observation contributes a kernel centred at its parameter value, with bandwidth determined adaptively. Conditional and discrete hyperparameters are handled by a hierarchical tree structure on $\Lambda\times\Theta$, with each node carrying its own one-dimensional kernel.

\begin{proposition}[EI under TPE is the ratio of densities]
\label{prop:tpe-ratio}
Under the TPE generative model, the expected improvement $\EI_{y^*}(\lambda,\theta)$ is proportional to the ratio of the conditional densities:
\[
  \EI_{y^*}(\lambda,\theta)\;\propto\;\left(\gamma+\frac{g(\lambda,\theta)}{\ell(\lambda,\theta)}\,(1-\gamma)\right)^{-1}.
\]
Consequently, maximising EI is equivalent to maximising the ratio $\ell(\lambda,\theta)/g(\lambda,\theta)$.
\end{proposition}

\begin{proof}[Proof sketch]
With $p(y\leq y^*)=\gamma$, Bayes' rule gives
\[
  p(\lambda,\theta)=\gamma\,\ell+(1-\gamma)\,g,
  \qquad
  p(y\leq y^*\mid\lambda,\theta)=\frac{\gamma\,\ell}{\gamma\,\ell+(1-\gamma)\,g}.
\]
Substitution into $\EI_{y^*}=\int_{-\infty}^{y^*}(y^*-y)p(y\mid\lambda,\theta)\,dy$ and elementary manipulation yields the displayed proportionality~\parencite{bergstra2011algorithms}.
\end{proof}

\subsection{Comparison with alternative search strategies}
\label{subsubsec:alternatives}

Four search paradigms are available for the optimisation of~\eqref{eq:bayes-objective}. The
present methodology adopts TPE-SMBO; this subsection provides a quantitative justification.

\paragraph{Grid search.}
Under the conditional structure of $\Lambda$ specified in
Section~\ref{subsubsec:bayes-search-space}, the total number of distinct
configurations is
\begin{equation}
  |\Lambda_{\mathrm{grid}}|
  = n_w\cdot n_m\cdot n_\tau\cdot n_r\,
    \bigl(1 + 2\,n_d\,n_\varepsilon\bigr)
  = 9\cdot 4\cdot 5\cdot 3\cdot(1+2\cdot 4\cdot 4)
  = 17{,}820,
  \label{eq:grid-count}
\end{equation}
where the factor $1$ accounts for the Alpha branch (metric and filtration-scale parameters
inactive) and $2\cdot 16=32$ for the Rips and SparseRips branches. Combined with the
classifier grids, $|\Lambda_{\mathrm{grid}}|\times(|\Theta_{\mathrm{lr}}|+|\Theta_{\mathrm{svm}}|)
= 17{,}820\times 18 = 320{,}760$ inner-CV evaluations are needed \emph{per outer LOSO fold}.
For a dataset with $|\Sset|=S$ outer LOSO folds, a conservative estimate of $5$~s per inner-CV
call gives approximately
\[
  320{,}760\times S\times 5\;\mathrm{s},
\]
which for $S=24$ already exceeds $43$~days, making full grid search computationally infeasible on the available hardware.

\paragraph{Random search.}
Random search~\parencite{bergstra2012random} draws configurations independently and uniformly
from $\Lambda\times\Theta$ without reference to past evaluations. It is
statistically consistent under infinite budget and, under the low effective-dimensionality
hypothesis, outperforms grid search at equal sample count by avoiding wasted evaluations
on uninformative axes. For the present budget $T=200$, however, random search explores
only $200/17{,}820\approx 1.1\%$ of the grid and provides no mechanism for directing
subsequent trials toward promising regions identified by earlier ones.

\paragraph{Gaussian-process Bayesian optimisation.}
The canonical SMBO surrogate is a Gaussian process (GP), which models the objective as a
GP-distributed random function and updates its posterior via Bayes' rule after each
evaluation. GP-SMBO is appropriate when the search space is a compact subset of $\R^d$
with a natural metric. It is unsuitable here for two structural reasons. First, no natural
stationary kernel is defined on the conditional, mixed-type space $\Lambda\times\Theta$
without auxiliary embeddings whose inductive bias is difficult to justify. Second, fitting
and inverting the GP covariance matrix costs $O(t^3)$ in the evaluation count $t$; for
$T=200$ iterations this overhead is non-negligible and grows quadratically with the
complexity of the correlation model.

\paragraph{TPE (adopted strategy).}
TPE~\parencite{bergstra2011algorithms} resolves both difficulties. The conditional structure of
$\Lambda$ is encoded natively in the tree; each leaf node carries an independent
univariate kernel density estimator, so per-iteration update cost is $O(t)$.
Proposition~\ref{prop:tpe-ratio} provides the theoretical guarantee that maximising
$\ell/g$ is equivalent to maximising EI under the TPE generative model. Empirically,
TPE has outperformed GP-EI on machine-learning hyperparameter benchmarks at comparable
budgets~\parencite{bergstra2011algorithms,bergstra2012random}, and its implementation in
Optuna~\parencite{optuna} handles mixed, conditional, and discrete spaces transparently.

\subsection{Search space}
\label{subsubsec:bayes-search-space}

The search space $\Lambda$ is a Cartesian product with one factor per component of the hyperparameter tuple~\eqref{eq:lambda}: a finite candidate set for the windowing parameter $w_{\mathrm{sec}}$, for the embedding triple $(m,\tau,r)$, and for the filtration quadruple $(c,d_{\mathrm{met}},\varepsilon,q)$. Several conditional constraints are imposed by the geometry itself rather than chosen arbitrarily: the Alpha complex requires the Euclidean metric, which deactivates the metric and landmark-fraction axes for that branch, since its complex size is already polynomially bounded (Section~\ref{subsubsec:filtration}) and needs no distance-based truncation. The classifier search space $\Theta$ is similarly a finite grid over each classifier family's own hyperparameters (regularisation strength, penalty or kernel choice) from Section~\ref{subsec:classifiers}. The specific candidate values used in this thesis --- identical across all three datasets, with one narrower exception for the Personal protocol to mitigate overfitting in smaller per-subject training sets --- are given where they are applied, in Table~\ref{tab:search_space} and Section~\ref{sec:exp:classifiers} of the Experiments chapter.

\subsection{Outer LOSO and inner cross-validation}
\label{subsubsec:nested-cv}

The selection procedure is the nested protocol
\begin{enumerate}[label=(\arabic*)]
  \item \textbf{Outer.} For each subject $s\in\Sset_K$, define $\mathcal{D}^{\mathrm{train}}_K(s)$ as the union of all windows of subjects $\Sset_K\setminus\{s\}$.
  \item \textbf{Inner.} Run TPE for a fixed budget of $T$ iterations on $\mathcal{D}^{\mathrm{train}}_K(s)$ using stratified $k$-fold CV; record $\lambda^*_K(s),\theta^*_K(s)$.
  \item \textbf{Retrain and test.} Fit a final model on $\mathcal{D}^{\mathrm{train}}_K(s)$ with $(\lambda^*_K(s),\theta^*_K(s))$ and evaluate on the held-out subject $s$.
\end{enumerate}
The outer LOSO ensures that no information from the test subject is used to choose any model component; the inner CV ensures that the surrogate objective is an unbiased estimator of expected performance on unseen subjects from the training pool.

\begin{remark}[Computational budget]
\label{rem:bayes-budget}
The budget $T=200$ was chosen empirically as the smallest value at which the running best objective stabilises across independent runs of the search on the three datasets. Larger budgets do not produce reliably different optima, given the noise level of the inner CV estimator.
\end{remark}

\begin{remark}[Implementation decomposition]
\label{rem:impl-decomp}
Equation~\eqref{eq:bayes-objective} states a joint optimisation over $\Lambda\times\Theta$.
In practice this joint search is decomposed into two sequential stages to reduce
computational cost.
\emph{Stage~1 (TDA search, Optuna):} A single TPE study is run per dataset for
$T=200$ iterations.
The objective is the average $F_2$ of a logistic regression proxy and a linear SVM
proxy, both with fixed hyperparameters, evaluated by three-fold CV on a balanced
training subset.
Using lightweight proxies keeps each function evaluation fast, so the budget is spent
exploring $\Lambda$ rather than the joint space; the averaged objective ensures the
recovered $\lambda^*$ is not overfitted to a single classifier's inductive bias.
\emph{Stage~2 (classifier search, \texttt{GridSearchCV}):} Given $\lambda^*$,
classifier hyperparameters $\theta^*$ are selected by exhaustive grid search with
three-fold CV on the features extracted under $\lambda^*$.
This decomposition approximates the joint optimum while keeping computation tractable.
\end{remark}

\begin{remark}[Classifier-agnostic TDA search]
\label{rem:agnostic-search}
Rather than running a separate Optuna study for each classifier family, a single
study is run with a classifier-agnostic objective.
Define
\[
  \lambda^* \;\in\; \argmax_{\lambda\in\Lambda}\;
  \frac{1}{2}\Bigl(
    \mathrm{CV}_K^{\mathrm{LR}}(\lambda) +
    \mathrm{CV}_K^{\mathrm{SVM}}(\lambda)
  \Bigr),
\]
where $\mathrm{CV}_K^f(\lambda)$ denotes the $K$-fold cross-validated $F_2$ obtained
by classifier $f$ on features extracted under $\lambda$, with $f$'s hyperparameters
fixed at proxy values.
The averaged objective recovers a configuration $\lambda^*$ whose features support
good performance across classifier families simultaneously.
The map $\phi_{\lambda^*}\colon\R^{L\times 2}\to\R^{24}$ must therefore produce an
image that is both approximately linearly separable (beneficial for LR) and
geometrically well-structured in pairwise distance (beneficial for SVM with a
nonlinear kernel).
A representation satisfying both conditions is a stronger finding than one tuned
to a single classifier's inductive bias.
\end{remark}

% ═════════════════════════════════════════════════════════════
\section{Evaluation Framework}
\label{subsec:eval_framework}

\subsection{Performance metrics}
\label{subsubsec:metric}

The choice of $F_2$ as the primary metric---including the inadequacy of Accuracy, Recall alone, and $F_1$ for the fall-detection setting, and the formal recall-weighting property of Proposition~\ref{prop:F2-recall-weight}---is developed in full in Section~\ref{subsubsec:metric-selection}. All four standard rates ($\Acc$, $\Rec$, $\Prec$, $F_1$) are reported in result tables for completeness; $F_2$ is used as the objective throughout the Experiments chapter.

\subsection{Evaluation protocols}
\label{subsubsec:protocols}

Three cross-validation protocols are defined, each measuring a different level of generalisation. They differ in how subject identity is respected during the train/test split. Let $\Sset$ denote the finite set of subject identifiers in a dataset, and let $|\Sset|=S$ be the number of subjects.

\paragraph{Naive protocol (subject-unaware random split).}
The feature matrix is partitioned uniformly at random into training ($70\%$) and test ($30\%$) sets with no constraint on subject membership. A subject may contribute windows to both partitions simultaneously. This violates the independence of training and test subjects and systematically overestimates generalisation performance; it is reported only as an upper-bound reference. The split is repeated ten times with independent seeds; results are means over repetitions.

\paragraph{General protocol (leave-one-subject-out cross-validation).}
For each subject $s\in\Sset$, the test fold is defined as all windows of subject $s$:
\[
  \mathcal{D}^{\mathrm{test}}(s) = \{(w,y):\mathrm{subject}(w)=s\},
  \qquad
  \mathcal{D}^{\mathrm{train}}(s) = \{(w,y):\mathrm{subject}(w)\neq s\}.
\]
The model is trained on $\mathcal{D}^{\mathrm{train}}(s)$ with the Bayesian-optimal hyperparameters and the $F_2$-optimal trained threshold (Section~\ref{subsubsec:threshold}), and evaluated on $\mathcal{D}^{\mathrm{test}}(s)$. This leave-one-subject-out (LOSO) scheme is the standard protocol for measuring generalisation to a previously unseen individual. Because the class-balancing procedure is stochastic, each LOSO fold is repeated $R=15$ times with independent random seeds; reported quantities are means or appropriate aggregates over the resulting $S\times R$ collection of evaluations. The General protocol is the principal protocol of this thesis.

\begin{remark}[Why LOSO rather than random $k$-fold]
\label{rem:loso-rationale}
Standard $k$-fold cross-validation assigns windows to folds at random, which means that windows from the same subject appear in both training and test folds. Because windows from the same subject are highly correlated---they share the same gait pattern, body mass, and recording conditions---random-fold performance is systematically optimistic and does not reflect deployment performance on a new user. LOSO eliminates this leakage by construction: the test fold is always disjoint from the training fold at the subject level.
\end{remark}

\paragraph{Personal protocol (within-subject cross-validation).}
For each $s\in\Sset$, a model is trained and tested exclusively on windows of subject $s$ via a $70/30$ random split, repeated $R=15$ times. The personal protocol represents the idealised setting in which a per-user calibration period is available before deployment; it provides an empirical ceiling for subject-specific personalisation.

\paragraph{Decision-threshold sweep (retrospective heterogeneity diagnostic).}
Using the General-protocol trained classifier, the decision threshold $\tau$ is swept over a uniform grid $\mathcal{T}\subset(0,1)$ for each held-out subject. The subject-level retrospective gain is
\begin{equation}
  \Delta F_2(s)=\max_{\tau\in\mathcal{T}}F_2(\tau,s)-F_2(\hat\tau_0(s),s),
  \label{eq:dF2}
\end{equation}
where $\hat\tau_0(s)$ is the PR-curve trained threshold for subject $s$. The population-level average
\[
  \overline{\Delta F_2}=\frac{1}{S}\sum_{s\in\Sset}\Delta F_2(s)
\]
quantifies the gap between the trained threshold and the subject-specific oracle. Because the oracle threshold $\tau^*(s)=\argmax_\tau F_2(\tau,s)$ is identified on the same test windows on which it is evaluated, the sweep is a retrospective diagnostic of inter-subject heterogeneity, not a valid prospective personalisation method.

\subsection{Statistical inference for performance measurement}
\label{subsubsec:stat-inference}

Every protocol above produces a collection of performance estimates rather than a single number: repeated evaluations under stochastic class-balancing, and one evaluation per held-out subject. Reporting only the mean or the maximum of such a collection discards information about how much that number could plausibly have varied. This subsection defines, in general and problem-independent terms, the machinery used to quantify that variability and to compare two configurations of the pipeline. Only the general method is stated here; the specific repetition counts, resampling budgets, and significance levels used in this thesis are given where the machinery is applied in the Experiments chapter.

Two sources of variability are distinguished. The \emph{within-subject} (or repetition-level) variability arises because a stochastic step of the pipeline---typically class-balancing---is repeated $R$ times with independent seeds for the same train/test split; the \emph{across-subject} variability is the spread of mean performance across the subject pool under the LOSO protocol of Section~\ref{subsubsec:protocols}. The two are reported separately rather than pooled, because they answer different questions: the first asks how reproducible a single subject's evaluation is, the second asks how much performance differs from one subject to the next.

\paragraph{Repetition-level confidence intervals.}
For each $(\text{subject } s,\text{repetition } j)$ pair let $x_{s,j}\in[0,1]$ denote the performance metric of interest. Define the per-repetition study-level mean
\[
  \bar X_j=|\Sset_K|^{-1}\sum_{s\in\Sset_K}x_{s,j},
  \qquad j=1,\ldots,R,
\]
the overall study mean $\overline{X}_K=R^{-1}\sum_j\bar X_j$, and the sample standard deviation of $\{\bar X_j\}$,
\[
  S_R=\Bigl(\frac{1}{R-1}\sum_{j=1}^R(\bar X_j-\overline{X}_K)^2\Bigr)^{1/2}.
\]

\begin{theorem}[Repetition-level $t$ confidence interval]
\label{thm:t-ci}
If the per-repetition study means $\bar X_j$ are exchangeable and approximately Gaussian, then a $1-\alpha$ confidence interval for the limiting expectation $\E[\bar X]$ as $R\to\infty$ is
\begin{equation}
  \overline{X}_K\;\pm\;t_{\alpha/2,\,R-1}\;\frac{S_R}{\sqrt R},
  \label{eq:t-ci}
\end{equation}
where $t_{\alpha/2,R-1}$ is the upper $\alpha/2$ quantile of the $t$-distribution with $R-1$ degrees of freedom.
\end{theorem}

The interval~\eqref{eq:t-ci} captures the repetition-to-repetition stochastic variability of the class-balancing procedure; it does not capture subject-level heterogeneity, which requires the paired and bootstrap tools below.

\paragraph{Subject-level paired comparisons.}
Many of the comparisons of interest are paired at the subject level: two classifiers, two protocols, or two hyperparameter configurations evaluated on the same subject pool. Let
\[
  \delta_s=A_s-B_s\in\R,\qquad s\in\Sset_K,
\]
denote the per-subject paired difference, with the sign convention that positive $\delta_s$ favours $A$. The summaries of interest are the mean $\bar\delta$, the standard deviation $S_\delta$, and the proportion $\widehat{\Prob}(\delta_s>0)$.

\paragraph{Bootstrap confidence intervals.}
For statistics whose finite-sample distribution is difficult to derive analytically, the nonparametric paired bootstrap is used. Let $\theta(\{\delta_s\})$ denote a real-valued statistic of the paired differences (e.g., the subject mean, median, or interquartile range). Generate $B$ bootstrap resamples $\{\delta_s^{*,b}\}_{s\in\Sset_K}$ by sampling $|\Sset_K|$ subjects with replacement from $\Sset_K$ and copying their $\delta_s$ values, and compute
\[
  \theta_b^*=\theta(\{\delta_s^{*,b}\}),\quad b=1,\ldots,B.
\]
The $1-\alpha$ percentile confidence interval is
\begin{equation}
  \bigl[\theta^*_{(\alpha/2)},\,\theta^*_{(1-\alpha/2)}\bigr],
  \label{eq:bootstrap}
\end{equation}
where $\theta^*_{(p)}$ is the $p$-quantile of $\{\theta_b^*\}$.

\paragraph{Permutation tests.}
For exact small-sample inference on the sign pattern of paired differences, the paired sign-flipping permutation test is used. Under the exchangeability null
\[
  H_0\colon\;\delta_s\;\overset{d}{=}\;-\delta_s\qquad\text{for all } s,
\]
the $2^{|\Sset_K|}$ sign-flipped vectors $\{\epsilon_s\delta_s\}$ with $\epsilon_s\in\{-1,+1\}$ are exchangeable in distribution. The exact two-sided $p$-value of an observed test statistic $T_{\mathrm{obs}}=T(\delta)$ is
\begin{equation}
  p=\frac{1}{2^{|\Sset_K|}}\sum_{\epsilon\in\{-1,+1\}^{|\Sset_K|}}\1\bigl\{|T(\epsilon\odot\delta)|\geq|T_{\mathrm{obs}}|\bigr\}.
  \label{eq:permutation}
\end{equation}
When $|\Sset_K|$ is small enough that the sum is computed exactly, the test is exact; otherwise it is estimated by sampling a large number of random sign vectors, which is a Monte Carlo approximation rather than a change of method.

\begin{remark}[Selection bias of headline maxima]
\label{rem:headline-bias}
The maximum of a finite collection of noisy estimators---such as the best single-subject score, or the best value found during hyperparameter search---is biased upward by the selection event itself: the bias grows with the number of competing candidates and with the variance of each estimator. Reporting subject means, the confidence intervals above, and the joint distribution of subject-level differences avoids this bias and yields estimands that are directly interpretable as expectations over the population from which the subjects were drawn. Headline maxima are therefore not used as primary evidence anywhere in this thesis.
\end{remark}
% ============================================================
%  END ch3.tex
% ============================================================

% ============================================================
%  BEGIN ch4.tex
% ============================================================
% ─────────────────────────────────────────────────────────────
\phantomsection
\chapter{Experiments}
\label{sec:experiments}

Chapter~\ref{sec:methodology} defines a parameterised pipeline
\[
  \phi_\lambda\colon w \;\longmapsto\; \phi_\lambda(w)\in\R^{24},
\]
mapping any scalar inertial window $w$ to a fixed feature vector via delay embedding, persistent homology, and vectorisation. The pipeline is controlled by a hyperparameter tuple $\lambda\in\Lambda$ whose optimal value is unknown a priori and is dataset-dependent. This chapter answers the question: \emph{what is $\lambda^*$ for each dataset, and what performance does it yield?}

The experimental strategy is a single-loop search:
\[
  \lambda^*,\theta^* \;=\; \argmax_{(\lambda,\theta)\,\in\,\Lambda\times\Theta}\; \widehat F_2^{\,\mathrm{CV}}(\lambda,\theta),
\]
where $\widehat F_2^{\,\mathrm{CV}}$ is the cross-validated $F_2$ score of the full pipeline evaluated on a given dataset. The search is conducted by Bayesian optimisation with a Tree-structured Parzen Estimator (TPE), which acts as the computational \emph{engine}: it proposes successive $(\lambda,\theta)$ configurations, evaluates the objective, and focuses the search on promising regions of the joint space. All other sections of this chapter exist to set up the input to this engine---the datasets that provide the objective-function evaluations, the preprocessing that makes the data comparable across sessions, and the feature description that defines the search space. Section~\ref{subsec:results} reports what the engine found.

Figure~\ref{fig:pipeline} gives an overview of the complete pipeline from raw sensor measurement to binary fall decision.

% ── Pipeline diagram ──────────────────────────────────────────
\begin{figure}[ht]
\centering
\begin{tikzpicture}[
  node distance = 0.55cm and 0.65cm,
  box/.style    = {rectangle, draw=black, rounded corners=3pt,
                   minimum width=1.80cm, minimum height=0.72cm,
                   align=center, font=\small, fill=white, thick},
  dbox/.style   = {rectangle, draw=black!60, rounded corners=3pt,
                   minimum width=1.80cm, minimum height=0.72cm,
                   align=center, font=\small, fill=gray!8,
                   dashed, thick},
  optbox/.style = {rectangle, draw=blue!55!black, rounded corners=5pt,
                   minimum width=4.6cm, minimum height=1.20cm,
                   align=center, font=\small, fill=blue!5, thick},
  arr/.style    = {->, >=Stealth, thick},
  oarr/.style   = {->, >=Stealth, thick, dashed, blue!55!black},
  lbl/.style    = {font=\footnotesize\bfseries, text=black!60}
]

%% ─── Row 1 : Preprocessing ──────────────────────────────────────
\node[box]                    (sensor)  {Sensor\\Data};
\node[box, right=of sensor]   (smv)     {SMV$(t)$};
\node[box, right=of smv]      (resamp)  {Resample\\50\,Hz};
\node[box, right=of resamp]   (win)     {Window\\Extract};

\draw[arr] (sensor)--(smv);
\draw[arr] (smv)--(resamp);
\draw[arr] (resamp)--(win);

%% ─── Row 2 : Feature Extraction (lambda-dependent) ──────────────
%% Placed directly below Row 1 (column-aligned: emb↔sensor, filt↔smv, …)
\node[dbox, below=1.20cm of sensor]  (emb)   {Delay\\Embed\\$(m,\tau,r)$};
\node[dbox, right=of emb]            (filt)  {Filtration\\$(c,\varepsilon)$};
\node[dbox, right=of filt]           (ph)    {Pers.\,Hom.\\$H_0,H_1$};
\node[dbox, right=of ph]             (vec)   {Vectorise\\$\phi_\lambda\!\in\!\R^{24}$};

\draw[arr] (emb)--(filt);
\draw[arr] (filt)--(ph);
\draw[arr] (ph)--(vec);

%% ─── Row 3 : Classification ──────────────────────────────────────
%% p(w) and tau_0 merged into one box to keep 4+4+4 column symmetry
\node[box, below=1.20cm of emb]  (scale)   {Std\,Scale};
\node[box, right=of scale]       (clf)     {Classifier\\(LR/SVM)};
\node[box, right=of clf]         (thrbox)  {$\hat{p}(w)\!\geq\!\tau_0$};
\node[box, right=of thrbox]      (dec)     {Decision\\Fall/ADL};

\draw[arr] (scale)--(clf);
\draw[arr] (clf)--(thrbox);
\draw[arr] (thrbox)--(dec);

%% ─── Row-to-row center connectors ────────────────────────────────
%% smv and filt are in column 2 (same x); straight vertical arrow.
\draw[arr] (smv.south) -- (filt.north);
%% filt and clf are in column 2 (same x); straight vertical arrow.
\draw[arr] (filt.south) -- (clf.north);

%% ─── Bayesian Optimisation box (left, aligned with Row 2) ───────
\node[optbox, left=0.75cm of emb]  (bayes)
  {\textbf{Bayesian Opt.\ (TPE,}\\
   $T{=}200$~\textbf{iterations)}\\[2pt]
   $\lambda^*=\argmax_{\lambda\in\Lambda}\widehat{F}_2^{\,\mathrm{CV}}(\lambda)$};

%% Horizontal arrow: bayes.east -> emb.west (same y-level)
\draw[oarr] (bayes.east) --
  node[above, font=\footnotesize\itshape]{$\lambda^*$}
  (emb.west);

%% ─── Phase background rectangles ──────────────────────────────
\begin{pgfonlayer}{background}
  \node[rounded corners=5pt, fill=cyan!7, draw=cyan!35, inner sep=5pt,
        label={[lbl]above:Preprocessing},
        fit=(sensor)(win)] {};
  \node[rounded corners=5pt, fill=orange!5, draw=orange!25, inner sep=5pt,
        label={[lbl, text=orange!65!black]above:%
          Feature Extraction~($\lambda$-dep.${}^{\dagger}$)},
        fit=(emb)(vec)] {};
  \node[rounded corners=5pt, fill=green!4, draw=green!28, inner sep=5pt,
        label={[lbl]above:Classification},
        fit=(scale)(dec)] {};
\end{pgfonlayer}

\node[font=\footnotesize, text=black!45,
      below=0.30cm of scale, anchor=north]
  {${}^{\dagger}$Theoretical basis: Methodology section.};

\end{tikzpicture}
\caption{Complete fall detection pipeline from raw inertial data to binary decision.
  \emph{Preprocessing}: SMV computation, resampling to 50~Hz, window extraction.
  \emph{Feature extraction} (dashed, $\lambda$-dependent): delay embedding,
  filtration construction, persistent homology ($H_0$, $H_1$), and
  vectorisation to $\phi_\lambda\in\R^{24}$.
  \emph{Classification}: standardisation, probabilistic classifier, and
  threshold comparison $\hat{p}(w)\geq\tau_0$ for binary decision.
  The \emph{Bayesian optimisation} box (left, blue) selects $\lambda^*$ by
  maximising cross-validated $F_2$ over training subjects via 200~TPE iterations;
  row-to-row connectors run through column~2 (SMV\,$\to$\,Filtration\,$\to$\,Classifier),
  preserving the 4\,$\times$\,3 grid symmetry.}
\label{fig:pipeline}
\end{figure}

% ─────────────────────────────────────────────────────────────
\section{Datasets}
\label{subsec:datasets}

Three publicly available benchmark datasets were used: MobiFall~\parencite{mobifall}, SisFall~\parencite{sisfall}, and FallAllD in its 40~Hz variant (hereafter FAD\_40Hz)~\parencite{fallalld}. Each dataset records inertial sensor signals from human subjects performing scripted fall events and daily activities. Together they provide diversity in sensor placement, subject age, and activity protocol, which is important for assessing whether the method generalises across acquisition conditions. Table~\ref{tab:dataset_summary} summarises the key characteristics of each dataset as used in this work.

\begin{table}[h!]
\centering
\caption{Summary of the three datasets as used in the experiments.}
\label{tab:dataset_summary}
\setlength{\tabcolsep}{4pt}\footnotesize
\begin{tabular}{lcccc}
\toprule
\textbf{Dataset} & \textbf{Subjects} & \textbf{Feature matrix}
  & \textbf{Resampling (Hz)} & \textbf{Placement} \\
\midrule
MobiFall   &  9 & $1808 \times 24$ & $87.5 \to 50$ & Trouser pocket \\
SisFall    & 24 & $7738 \times 24$ & $200  \to 50$ & Waist \\
FAD\_40Hz  & 13 & $9481 \times 24$ & $40   \to 50$ & Waist \\
\midrule
\textbf{Total} & \textbf{46} & $\mathbf{19{,}027 \times 24}$ & --- & --- \\
\bottomrule
\end{tabular}
\end{table}

\subsection{Common target sampling rate}
\label{subsubsec:sampling-rate}

All three datasets were resampled to a common target rate of $f_s=50$~Hz before any further processing. The theoretical lower bound for this choice is given by the Nyquist--Shannon theorem (Theorem~\ref{thm:nyquist} in the Methodology): with discriminative fall content concentrated below $20$~Hz, any rate $f_s\geq 40$~Hz is sufficient. The target $f_s=50$~Hz satisfies this bound with a $25\%$ margin. Two additional practical considerations support this choice.

\paragraph{Integer-ratio resampling.} For SisFall the conversion ratio $200{:}50=4{:}1$ is an integer, reducing resampling to decimation by four. For FAD\_40Hz the ratio $40{:}50=4{:}5$ and for MobiFall the ratio $87.5{:}50=7{:}4$; both admit finite-impulse-response polyphase resamplers at moderate cost.

\paragraph{Computational tractability.} The persistent homology computation dominates pipeline cost. A five-second window contains $L=250$ samples at $50$~Hz versus $L=1000$ samples at $200$~Hz, reducing the point cloud size by a factor of four and making a $200$-iteration Bayesian search across three datasets feasible on the available hardware.

\subsection{MobiFall}
\label{subsubsec:mobifall}

MobiFall v2.0 contains recordings from volunteers who carried a Samsung Galaxy S3 smartphone in their trouser pocket while performing four scripted fall types --- forward fall (hands first), forward fall onto knees, backward fall from a chair, and lateral fall from standing --- as well as nine daily activities such as walking, jogging, stair use, and car entry~\parencite{mobifall}. The nominal recording rate is 87.5~Hz. The device captures seven sensor channels: three-axis linear acceleration, three-axis angular velocity (gyroscope), three-axis magnetic field, barometric pressure, and device orientation expressed as azimuth, pitch, and roll angles.

For this study, only the accelerometer and gyroscope channels were retained. The magnetometer, barometric pressure signal, and orientation estimates were excluded for two reasons. First, these channels are absent or recorded at incompatible rates in the other two datasets, so retaining them would have prevented a methodologically consistent preprocessing pipeline across all three sources. Second, the magnetometer and barometric sensor do not directly encode the inertial dynamics of fall events; including them would increase feature dimensionality without a principled contribution to the classification objective.

The full v2.0 cohort contains participants at different stages of data collection; not all subjects have both fall and daily-activity recordings. Nine subjects (identifiers 2, 3, 4, 5, 7, 8, 9, 10, and 11) were included in this study. The remaining participants were excluded because they performed fall trials only, without the paired daily-activity recordings required for leave-one-subject-out evaluation: a held-out subject who has no daily-activity windows cannot contribute a meaningful test set to the General protocol. After windowing, the retained cohort yields a feature matrix of $1808 \times 24$.

\subsection{SisFall}
\label{subsubsec:sisfall}

SisFall is a waist-worn dataset capturing 15 fall types and 19 daily-activity types across a mixed group of young and older adults~\parencite{sisfall}. The sensor unit is mounted at the waist and houses two accelerometers and one gyroscope: the ADXL345 (dynamic range $\pm 16\,g$, 13-bit resolution, used primarily for fall detection), the MMA8451Q (dynamic range $\pm 8\,g$, 14-bit resolution), and the ITG3200 gyroscope (dynamic range $\pm 2000$\,°/s, 16-bit resolution). All channels are recorded at 200~Hz. The dataset encompasses 23 young adults and 15 older adults.

Of the two accelerometers, only the ADXL345 was used alongside the gyroscope. The ADXL345 was preferred because its higher dynamic range ($\pm 16\,g$) makes it better suited to capturing the impulsive acceleration peaks that characterise fall impacts; the MMA8451Q, with half the range, is susceptible to saturation during vigorous motions. Retaining a single accelerometer also maintains sensor-channel consistency with MobiFall and FAD\_40Hz, each of which provides only one accelerometer. Including the MMA8451Q in parallel would have introduced an asymmetry in feature construction across datasets and complicated the cross-dataset comparison.

Twenty-four subjects were retained: the 23 young adults and one older adult (subject identifier 29) who had sufficient trial counts. The remaining 14 older adult participants were excluded due to an insufficient number of recorded trials to support stable leave-one-subject-out evaluation. The original 200~Hz recording rate was downsampled to 50~Hz. The resulting feature matrix has dimensions $7738 \times 24$, making SisFall the largest dataset by subject count.

\subsection{FallAllD (FAD\_40Hz)}
\label{subsubsec:fallalld}

\paragraph{Original FallAllD dataset.}
FallAllD is a large-scale open dataset of human falls and activities of daily living collected from 15 subjects~\parencite{fallalld}. In its original release, recordings are obtained simultaneously from three sensor devices worn at the waist, wrist, and neck, each capturing four modalities: a three-axis accelerometer, a three-axis gyroscope, a three-axis magnetometer, and a barometric pressure sensor. The inertial sensors (accelerometer and gyroscope) are sampled at 238~Hz, the magnetometer at 80~Hz, and the barometer at 10~Hz. The complete activity label set comprises 135 categories, of which 44 are activities of daily living (labelled with identifiers 1--44) and 91 are distinct fall scenarios (identifiers 101--135) differing in direction, mechanism, and recovery condition. Across all subjects, sensor positions, and activities, the original release contains 26{,}420 individual recording files.

\paragraph{Selection of the 40~Hz derived variant.}
Using the original 238~Hz release directly in the present study would have imposed prohibitive computational costs. The Bayesian hyperparameter search runs 200 iterations per dataset, and each iteration requires extracting topological features --- which involve constructing and computing persistent homology of a point cloud --- from every window in the dataset under the current configuration. At 238~Hz, the number of samples per window is approximately 4.8 times larger than at 50~Hz, substantially increasing both the size of the delay-embedded point cloud and the cost of the homology computation. Furthermore, the four-modality, three-position structure of the original release creates a sensor fusion problem that falls outside the scope of the present methodology, which targets single-position inertial data.

For these reasons, the publicly available 40~Hz derived variant (FAD\_40Hz) was used instead~\parencite{fallalld}. This derived dataset has been pre-processed as follows relative to the original release: (i)~the inertial signals are downsampled from 238~Hz to 40~Hz; (ii)~only the accelerometer and gyroscope channels are retained, with the magnetometer and barometric sensors removed; (iii)~the neck-worn device data is excluded, leaving the waist and wrist positions; (iv)~the activity label set is filtered to a curated subset of fall and ADL types relevant to binary detection tasks. These modifications reduce the data volume by an order of magnitude and eliminate the multi-rate sensor-fusion problem imposed by the original release, while preserving the inertial channels and the subject-level label structure required by the present pipeline.

The 40~Hz signals were subsequently resampled to the common target rate of 50~Hz as described in Section~\ref{subsec:preprocessing}. After activity filtering and subject selection, 13 of the original 15 subjects contributed data; the two excluded subjects had an insufficient number of trials for stable leave-one-subject-out evaluation. The resulting feature matrix is $9{,}481 \times 24$, the largest in the study by sample count.

% ─────────────────────────────────────────────────────────────
\section{Data Processing and Preprocessing}
\label{subsec:preprocessing}

Before any feature computation, all sensor recordings were passed through a common preprocessing pipeline consisting of four sequential steps. This pipeline transforms raw three-axis measurements into a compact, orientation-independent scalar signal, resamples it to a common frequency, and partitions it into fixed-length windows ready for feature extraction. Figure~\ref{fig:pipeline} places this block at the start of the data path.

\textbf{Step 1 --- Signal Magnitude Vector.}
For each sensor (accelerometer and gyroscope separately), the three-axis measurement vector $\mathbf{x}(t) \in \R^3$ at time $t$ is reduced to its Euclidean norm:
\[
  \mathrm{SMV}(t) = \bigl\|\mathbf{x}(t)\bigr\|_2 = \sqrt{x_1(t)^2 + x_2(t)^2 + x_3(t)^2}.
\]
This produces one scalar time series per sensor that is invariant to device orientation on the body, eliminating the need for explicit attitude estimation.

\textbf{Step 2 --- Resampling to 50~Hz.}
The three datasets were originally recorded at 87.5~Hz (MobiFall), 200~Hz (SisFall), and 40~Hz (FAD\_40Hz). All signals were resampled to a common rate of 50~Hz, but the strategies differ by dataset because the resampling direction and the ratio of original to target rate are not uniform.

For SisFall (200~Hz $\to$ 50~Hz), the ratio is the integer 4:1, so integer-ratio decimation was applied: the SMV signal was subsampled by retaining every fourth sample. This produces a uniformly-spaced 50~Hz output without any interpolation step and without spectral distortion within the passband of interest, since the Nyquist frequency of the target rate (25~Hz) is well below the dominant frequency of fall events.

For FAD\_40Hz (40~Hz $\to$ 50~Hz), upsampling by a factor of 5:4 was required. The resampled signal was obtained by constructing a new time grid at 50~Hz and computing signal values at each new grid point by interpolating over the flanking 40~Hz samples. Because the ratio 5:4 is rational and the upsampling factor is moderate, this approach introduces no perceptible temporal smearing in the SMV magnitude signal.

For MobiFall (87.5~Hz $\to$ 50~Hz), the ratio 87.5:50 = 7:4 is rational but non-integer in both directions, so neither pure decimation nor pure upsampling is sufficient. A rational sample-rate conversion was applied, conceptually equivalent to first upsampling by a factor of 4 and then downsampling by a factor of 7. The resampling was performed on the scalar SMV signal rather than on the raw three-axis components, so the non-linear magnitude operation is not affected by the interpolation kernel.

In all three cases the resampling target of 50~Hz was selected on methodological grounds that are discussed in the Methodology section.

\textbf{Step 3 --- Window extraction.}
For each fall trial, a single window centred on the highest-magnitude point of the SMV signal was extracted. For daily-activity trials, consecutive non-overlapping windows of the same length were taken. The window length is a hyperparameter optimised during the experiment and is described in Section~\ref{subsec:hparam}.

\textbf{Step 4 --- Standard scaling.}
Each of the 24 features in the final feature vector was scaled to zero mean and unit variance. The scaling parameters were estimated from the training data only and applied to the test data, ensuring that no information from test subjects influences the normalisation.

\begin{figure}[H]
\centering
\begin{tikzpicture}[
  node distance = 0.5cm,
  io/.style       = {rectangle, draw=black!65, rounded corners=4pt,
                     minimum width=9cm, text width=8.8cm, minimum height=0.65cm,
                     align=center, font=\small, fill=gray!10, thick},
  flowstep/.style = {rectangle, draw=black, rounded corners=3pt,
                     minimum width=9cm, text width=8.8cm, minimum height=0.65cm,
                     align=center, font=\small, fill=white, thick},
  arr/.style      = {->, >=Stealth, thick}
]
\node[io] (in)
  {Raw inertial recording:\;$\mathrm{Acc}(t)\in\R^3$,\;$\mathrm{Gyr}(t)\in\R^3$};
\node[flowstep, below=0.5cm of in] (smv)
  {\textbf{Step 1\;---\;SMV:}\;$\mathrm{SMV}(t)=\|\mathbf{x}(t)\|_2$\quad(per sensor)};
\node[flowstep, below=0.5cm of smv] (res)
  {\textbf{Step 2\;---\;Resample to 50\,Hz}\;(MobiFall 87.5, SisFall 200, FAD 40\,Hz)};
\node[flowstep, below=0.5cm of res] (win)
  {\textbf{Step 3\;---\;Window Extraction}\;(length $L$;\;fall: centred on peak,\;ADL: stride $L$)};
\node[io, below=0.5cm of win] (out)
  {Labelled windows $\{(w_i,\,y_i,\,s_i)\}$};
\draw[arr] (in)--(smv);
\draw[arr] (smv)--(res);
\draw[arr] (res)--(win);
\draw[arr] (win)--(out);
\end{tikzpicture}
\caption{Preprocessing pipeline (Steps 1--3). Detailed description of each step
  is given in Section~\ref{subsec:preprocessing}.}
\label{fig:preprocessing}
\end{figure}

% ─────────────────────────────────────────────────────────────
\section{Feature Description}
\label{subsec:features}

Each window is transformed into a 24-dimensional feature vector. The features are divided into two groups: 20 persistence-based features that capture the topological shape of the motion trajectory, and 4 signal-level features that capture the gross magnitude behaviour of the accelerometer signal. The theoretical derivation of the persistence-based features is presented in the Methodology section; here only the list of features and their definitions are given.

% ── 1.3.1 Persistence-Based Features ─────────────────────────
\subsection{Persistence-Based Features}
\label{subsubsec:ph_features}

The topological representation of each window yields two sets of birth--death pairs: one for connected clusters in the trajectory ($H_0$) and one for closed loops ($H_1$). Ten scalar statistics are extracted from each set, giving 20 features in total. Table~\ref{tab:ph_features} lists these statistics and their definitions.

\begin{table}[h!]
\centering
\caption{Ten persistence-based statistics computed per homological dimension.
  Each row is computed twice: once for $H_0$ (features 1--10) and once for $H_1$
  (features 11--20).  Persistence of a topological feature is defined as its
  death scale minus its birth scale.}
\label{tab:ph_features}
\small
\begin{tabular}{clp{8cm}}
\toprule
\textbf{Index} & \textbf{Feature} & \textbf{Definition} \\
\midrule
1 / 11  & Bar count          & Number of features born and died, normalised by window length \\
2 / 12  & Entropy            & Shannon entropy of the normalised persistence values \\
3 / 13  & Maximum            & Persistence of the longest-lived feature \\
4 / 14  & Energy             & Sum of squared persistence values \\
5 / 15  & Mean               & Arithmetic mean of all persistence values \\
6 / 16  & Std deviation      & Standard deviation of the persistence values \\
7 / 17  & Median             & Median persistence \\
8 / 18  & First quartile     & 25th percentile of the persistence distribution \\
9 / 19  & Third quartile     & 75th percentile of the persistence distribution \\
10 / 20 & Third moment       & Third central moment (related to skewness) of the persistence values \\
\bottomrule
\end{tabular}
\end{table}

% ── 1.3.2 Signal-Level Features ──────────────────────────────
\subsection{Signal-Level Features}
\label{subsubsec:signal_features}

Four time-domain statistics are computed directly from the accelerometer SMV$(t)$ values within the window and appended to the feature vector, giving features 21--24. These capture the raw intensity profile of the motion event without any topological processing.

\begin{table}[h!]
\centering
\caption{Four signal-level features derived from the accelerometer magnitude
  SMV$(t)$ within each window.}
\label{tab:sig_features}
\small
\begin{tabular}{clp{8cm}}
\toprule
\textbf{Index} & \textbf{Feature} & \textbf{Definition} \\
\midrule
21 & Maximum       & Highest value of SMV$(t)$ in the window \\
22 & Std deviation & Standard deviation of SMV$(t)$ \\
23 & Mean          & Arithmetic mean of SMV$(t)$ \\
24 & Range         & Maximum minus minimum of SMV$(t)$ \\
\bottomrule
\end{tabular}
\end{table}

\begin{figure}[H]
\centering
\begin{tikzpicture}[
  node distance = 0.5cm,
  io/.style       = {rectangle, draw=black!65, rounded corners=4pt,
                     minimum width=9cm, text width=8.8cm, minimum height=0.65cm,
                     align=center, font=\small, fill=gray!10, thick},
  flowstep/.style = {rectangle, draw=black, rounded corners=3pt,
                     minimum width=9cm, text width=8.8cm, minimum height=0.65cm,
                     align=center, font=\small, fill=white, thick},
  arr/.style      = {->, >=Stealth, thick}
]
\node[io] (in)
  {Window $w\in\R^{L\times 2}$\;($\mathrm{SMV}_\mathrm{acc}$,\;$\mathrm{SMV}_\mathrm{gyr}$)};
\node[flowstep, below=0.5cm of in] (emb)
  {\textbf{Step 1\;---\;Delay Embedding}\;$(m,\tau,r)$\;$\to P_\lambda(w)\subset\R^m$};
\node[flowstep, below=0.5cm of emb] (filt)
  {\textbf{Step 2\;---\;Filtration}\;$(c,\varepsilon)$\;$\to K_c(\varepsilon)$\;
   (\texttt{Alpha} / \texttt{Rips} / \texttt{SparseRips})};
\node[flowstep, below=0.5cm of filt] (ph)
  {\textbf{Step 3\;---\;Persistent Homology}\;$\to\mathrm{PD}_0\;(H_0)$,\;$\mathrm{PD}_1\;(H_1)$};
\node[flowstep, below=0.5cm of ph] (vec)
  {\textbf{Step 4\;---\;Vectorisation}\;$\to$\;feat.\ 1--10\;($\mathrm{PD}_0$),
   \;11--20\;($\mathrm{PD}_1$),\;21--24\;(signal)};
\node[io, below=0.5cm of vec] (out)
  {Feature vector $\phi_\lambda(w)\in\R^{24}$};
\draw[arr] (in)--(emb);
\draw[arr] (emb)--(filt);
\draw[arr] (filt)--(ph);
\draw[arr] (ph)--(vec);
\draw[arr] (vec)--(out);
\end{tikzpicture}
\caption{Feature extraction pipeline (Steps 1--4). Detailed description of each step
  is given in Section~\ref{subsec:features}.}
\label{fig:features}
\end{figure}

% ─────────────────────────────────────────────────────────────
\section{Bayesian Hyperparameter Optimisation}
\label{subsec:hparam}

The topological feature computation depends on several configuration choices: the window length, the phase-space embedding parameters, and the type of topological complex. Rather than fixing these by hand or by exhaustive grid search, a Bayesian optimisation procedure was used to search for the combination that maximises fall detection performance on each dataset.

A single Optuna study was run per dataset, yielding 200 TPE iterations per dataset and 600 in total across the three datasets. Each trial proposes a TDA configuration $\lambda_t$; features are extracted under that configuration from a balanced training subset (all fall windows plus a random sample of three times as many ADL windows) and evaluated with a classifier-agnostic objective: both a logistic regression proxy and a linear SVM proxy are scored with three-fold stratified cross-validation and their $F_2$ scores are averaged. Classifier hyperparameters are kept fixed during this stage and are subsequently optimised by grid search (Section~\ref{sec:exp:classifiers}).

The search is guided by a Tree-structured Parzen Estimator (TPE). After each evaluation, the history of observations $\mathcal{H}_t=\{(\lambda_i,\hat{y}_i)\}$ is split at a quantile $q_\gamma$ into a good subset ($\hat{y}\geq q_\gamma$) and a bad subset ($\hat{y}<q_\gamma$). Kernel density estimates $\ell(\lambda)$ and $g(\lambda)$ are fitted to the configurations in the good and bad subsets respectively. The next candidate is proposed as $\lambda_{t+1}=\argmax_\lambda\,\ell(\lambda)/g(\lambda)$, the configuration most probable under the good density relative to the bad density. This ratio is proportional to the expected improvement over the current best observation (theoretical derivation in the Methodology section). The mechanism progressively focuses the search on regions of $\Lambda$ associated with high $F_2$, without evaluating all combinations exhaustively.

The averaged objective ensures that the recovered configuration $\lambda^*$ is a representation that performs well across classifier families rather than being tuned to a single model's inductive bias. A representation that simultaneously supports linear separability (beneficial for LR) and a geometrically structured metric space (beneficial for SVM with an RBF kernel) is a stronger finding than one optimised for a single classifier.

Table~\ref{tab:search_space} lists the parameters and the values considered during the search.

\begin{table}[h!]
\centering
\caption{Parameters and candidate values explored during Bayesian optimisation.
  A dash (---) indicates parameters not applicable for the Alpha complex type.}
\label{tab:search_space}
\small
\begin{tabular}{ll}
\toprule
\textbf{Parameter} & \textbf{Candidate values} \\
\midrule
Window length (\texttt{win\_sec}, seconds)
  & 1.0,\; 1.5,\; 2.0,\; 2.5,\; 3.0,\; 3.5,\; 4.0,\; 4.5,\; 5.0 \\
Embedding dimension (\texttt{dim})
  & 2,\; 3,\; 4,\; 5 \\
Embedding delay (\texttt{delay})
  & 1,\; 2,\; 3,\; 4,\; 5 \\
Stride factor (\texttt{stride\_factor})
  & 1,\; 2,\; 4 \\
Complex type (\texttt{complex\_type})
  & Alpha,\; Rips,\; SparseRips \\
Distance measure (\texttt{metric})
  & Euclidean,\; Cosine,\; Manhattan,\; Chebyshev \\
Filtration scale (\texttt{eps\_percentile})
  & 20\,\%,\; 40\,\%,\; 60\,\%,\; 80\,\% \\
\bottomrule
\end{tabular}
\end{table}

\begin{figure}[H]
\centering
\begin{tikzpicture}[
  node distance = 0.6cm,
  lbl/.style  = {font=\footnotesize\bfseries, text=blue!50!black},
  term/.style = {rectangle, draw=black!65, rounded corners=8pt,
                 minimum width=8cm, text width=7.8cm, minimum height=0.65cm,
                 align=center, font=\small, fill=gray!10, thick},
  proc/.style = {rectangle, draw=black, rounded corners=3pt,
                 minimum width=8cm, text width=7.8cm, minimum height=0.65cm,
                 align=center, font=\small, fill=white, thick},
  cond/.style = {rectangle, draw=black!60, rounded corners=3pt,
                 minimum width=8cm, text width=7.8cm, minimum height=0.65cm,
                 align=center, font=\small, fill=yellow!8, dashed, thick},
  arr/.style  = {->, >=Stealth, thick},
  larr/.style = {->, >=Stealth, thick, blue!55!black}
]
\node[term] (init)
  {Initialise:\;$T_0$ random configs\;$\to H=\{(\lambda_i,\hat{y}_i)\}$};
\node[proc, below=1.4cm of init] (eval)
  {Evaluate:\;$\hat{y}_t=\mathrm{mean}(\widehat{F}_2^{\mathrm{LR}},\widehat{F}_2^{\mathrm{SVM}})$,\;3-fold CV,\;update $H$};
\node[proc, below=0.6cm of eval] (update)
  {Update TPE:\;fit $\ell(\lambda)$ (good) and $g(\lambda)$ (bad)};
\node[proc, below=0.6cm of update] (propose)
  {Propose:\;$\lambda_{t+1}=\argmax_\lambda\,\ell(\lambda)/g(\lambda)$,\;
   $t\!\leftarrow\!t{+}1$};
\node[term, below=1.4cm of propose] (out)
  {Output:\;$\lambda^*=\argmax_t\,\hat{y}_t$};
%% flow
\draw[arr] (init)--(eval);
\draw[arr] (eval)--(update);
\draw[arr] (update)--(propose);
\draw[arr] (propose)--(out);
%% loop-back from propose to eval
\draw[larr] (propose.east)--++(1.2,0)
  |- node[right,font=\footnotesize\itshape,pos=0.25]{repeat}(eval.east);
%% loop background
\begin{pgfonlayer}{background}
  \node[rounded corners=5pt, draw=blue!30, fill=blue!3, inner sep=8pt,
        label={[lbl]above:{Optimisation loop ($T=200$ iterations per dataset)}},
        fit=(eval)(propose)] {};
\end{pgfonlayer}
\end{tikzpicture}
\caption{Bayesian hyperparameter optimisation loop (TPE, $T=200$ iterations per dataset).
  Each trial evaluates both an LR proxy and a linear SVM proxy; the objective is
  the mean of their $F_2$ scores, yielding a classifier-agnostic $\lambda^*$.
  Detailed description in Section~\ref{subsec:hparam}.}
\label{fig:bayesian}
\end{figure}

% ─────────────────────────────────────────────────────────────
\section{Classifiers and Training Setup}
\label{sec:exp:classifiers}

Two classifiers were compared throughout the experiments: logistic regression (LR) and a support vector machine (SVM). Both were configured with balanced class weights so that the minority fall class receives proportionally more influence during training, compensating for the natural class imbalance between fall and daily-activity windows.

Both classifiers operated on the same TDA feature representation extracted under the single optimal configuration $\lambda^*$. This shared representation reflects the classifier-agnostic optimisation objective: $\lambda^*$ is the configuration whose features support good performance across both LR and SVM rather than being tailored to either one individually.

Classifier hyperparameters were selected by \texttt{GridSearchCV} with three-fold stratified cross-validation, scoring by $F_2$. For logistic regression, the grid covered regularisation strength $C \in \{0.01,\,0.1,\,1,\,10,\,100\}$ and penalty type $\{\ell_1,\,\ell_2\}$. For the SVM, the grid covered $C \in \{0.1,\,1,\,10,\,100\}$ and kernel $\{\text{linear},\,\text{rbf}\}$; probability outputs were obtained via Platt scaling. All feature vectors were standardised (zero mean, unit variance, parameters estimated from training data only) within each cross-validation fold before fitting.

After grid-search selects the best estimator, the decision threshold $\tau$ is set by inspecting the precision--recall curve on the training fold: the threshold $\tau^* = \argmax_\tau F_2(\tau)$ that maximises $F_2$ along the curve replaces the default 0.5 cut-off. This calibration requires no additional model fitting and adapts the operating point to the recall-weighted $F_2$ objective.

% ─────────────────────────────────────────────────────────────
\section{Results}
\label{subsec:results}

The evaluation protocols (Naive, General, and Personal) and the performance metrics ($\Acc$, $\Rec$, $\Prec$, $F_1$, $F_2$) are defined in full in the Methodology section. Briefly, the General protocol uses leave-one-subject-out cross-validation and is the primary evaluation scenario; the Naive protocol uses a random 70/30 split and serves as an upper-bound reference; and the Personal protocol trains and tests on each subject's own data. Results under the General and Personal protocols are averaged over 15 repetitions and reported with 95\,\% confidence intervals. $F_2$ is the primary metric throughout.

\subsection{Optimal Hyperparameter Configurations}
\label{subsubsec:optimal_params}

Table~\ref{tab:optimal_params} shows the best configuration found by the Bayesian search for each dataset, together with the corresponding cross-validated $F_2$ score. The three datasets converged on markedly different settings.

\begin{table}[h!]
\centering
\caption{Best configuration selected by Bayesian optimisation for each dataset
  (200 iterations per dataset).  A dash (---) indicates a parameter not
  applicable for the Alpha complex.}
\label{tab:optimal_params}
\small
\begin{tabular}{lccc}
\toprule
\textbf{Parameter} & \textbf{MobiFall} & \textbf{SisFall} & \textbf{FAD\_40Hz} \\
\midrule
Best $F_2$ (search objective) & 0.9746      & 0.9733   & 0.9860      \\
Window length (s)             & 5.0         & 1.0      & 1.0         \\
Complex type                  & SparseRips  & Alpha    & Rips        \\
Embedding dimension           & 5           & 5        & 3           \\
Embedding delay               & 2           & 4        & 2           \\
Stride factor                 & 2           & 4        & 1           \\
Distance measure              & Chebyshev   & ---      & Manhattan   \\
Filtration scale              & 80\,\%      & ---      & 40\,\%      \\
\bottomrule
\end{tabular}
\end{table}

MobiFall required a long window (5.0~s) and a SparseRips complex, whereas both SisFall and FAD\_40Hz were best described by a short 1.0~s window. SisFall uniquely preferred the Alpha complex; FAD\_40Hz favoured the standard Rips construction. These differences suggest that the three datasets exhibit genuinely different motion geometries rather than a single universal pattern, which is a direct instance of Hypothesis~2 in the Methodology section.
% TODO: [Cross-reference to Discussion section once written.]

\subsection{MobiFall Results}
\label{subsubsec:mobifall_results}

Table~\ref{tab:mobifall_results} reports LR and SVM performance across the three evaluation protocols on MobiFall.

\begin{table}[h!]
\centering
\caption{MobiFall: mean performance across subjects under the three evaluation
  protocols.  General and Personal entries are mean $\pm$ 95\,\% CI over
  $n = 15$ repetitions per subject; Naive entries are means over 10 repetitions.}
\label{tab:mobifall_results}
\setlength{\tabcolsep}{3pt}\footnotesize
\begin{tabular}{llccccc}
\toprule
\textbf{Classifier} & \textbf{Protocol}
  & \textbf{Acc} & \textbf{Rec} & \textbf{Pre}
  & \textbf{$F_1$} & \textbf{$F_2$} \\
\midrule
\multirow{3}{*}{LR}
 & Naive
   & 0.9722 & 0.9512 & 0.6992 & 0.8042 & 0.8857 \\
 & General
   & $0.9774 \pm 0.0083$ & $0.9574 \pm 0.0449$ & ---
   & $0.8400 \pm 0.0487$ & $0.9049 \pm 0.0393$ \\
 & Personal
   & $0.9713 \pm 0.0098$ & $0.7319 \pm 0.0309$ & ---
   & $0.7488 \pm 0.0526$ & $0.7325 \pm 0.0295$ \\
\midrule
\multirow{3}{*}{SVM}
 & Naive
   & 0.9605 & 0.9558 & 0.6177 & 0.7461 & 0.8573 \\
 & General
   & $0.9734 \pm 0.0080$ & $0.9543 \pm 0.0404$ & ---
   & $0.8176 \pm 0.0455$ & $0.8918 \pm 0.0327$ \\
 & Personal
   & $0.9635 \pm 0.0096$ & $0.5822 \pm 0.0432$ & ---
   & $0.6371 \pm 0.0611$ & $0.5981 \pm 0.0479$ \\
\bottomrule
\end{tabular}
\end{table}

Under the General protocol LR reached $F_2 = 0.9049 \pm 0.0393$ and SVM $F_2 = 0.8918 \pm 0.0327$. The Personal protocol caused a substantial drop for both classifiers, with SVM falling to $F_2 = 0.5981$. This behaviour indicates that MobiFall provides too few trials per subject to support a reliable individual model; the subject-level sample count is simply insufficient to estimate a stable decision boundary from within-subject data alone.

\subsection{SisFall Results}
\label{subsubsec:sisfall_results}

Table~\ref{tab:sisfall_results} presents SisFall performance.

\begin{table}[h!]
\centering
\caption{SisFall: mean performance across subjects.  Format as in
  Table~\ref{tab:mobifall_results}.}
\label{tab:sisfall_results}
\setlength{\tabcolsep}{3pt}\footnotesize
\begin{tabular}{llccccc}
\toprule
\textbf{Classifier} & \textbf{Protocol}
  & \textbf{Acc} & \textbf{Rec} & \textbf{Pre}
  & \textbf{$F_1$} & \textbf{$F_2$} \\
\midrule
\multirow{3}{*}{LR}
 & Naive
   & 0.9788 & 0.9821 & 0.9308 & 0.9557 & 0.9713 \\
 & General
   & $0.9776 \pm 0.0045$ & $0.9836 \pm 0.0130$ & ---
   & $0.9529 \pm 0.0100$ & $0.9709 \pm 0.0108$ \\
 & Personal
   & $0.9656 \pm 0.0056$ & $0.9496 \pm 0.0073$ & ---
   & $0.9284 \pm 0.0113$ & $0.9405 \pm 0.0080$ \\
\midrule
\multirow{3}{*}{SVM}
 & Naive
   & 0.9790 & 0.9822 & 0.9314 & 0.9561 & 0.9715 \\
 & General
   & $0.9832 \pm 0.0040$ & $0.9742 \pm 0.0138$ & ---
   & $0.9642 \pm 0.0089$ & $0.9700 \pm 0.0111$ \\
 & Personal
   & $0.9604 \pm 0.0059$ & $0.9231 \pm 0.0096$ & ---
   & $0.9162 \pm 0.0110$ & $0.9196 \pm 0.0082$ \\
\bottomrule
\end{tabular}
\end{table}

SisFall yielded the highest General $F_2$ of the three datasets: LR achieved $0.9709 \pm 0.0108$ and SVM $0.9700 \pm 0.0111$. The two classifiers are nearly indistinguishable on this dataset, reflecting the large and well-balanced subject pool (24 subjects, 15 fall types). In the Personal protocol performance dropped moderately but remained above 0.93 for LR ($F_2 = 0.9405$) and above 0.91 for SVM ($F_2 = 0.9196$), showing that SisFall subjects contribute enough individual data to support within-subject training --- a contrast with MobiFall that illustrates how dataset size governs the viability of subject-specific models.

% TODO: [Per-subject breakdown tables for SisFall (24 subjects) exist in
%  TEZ_DataSet_TR.tex.  Confirm whether to include a condensed per-subject
%  table here or in an appendix.]

\subsection{FAD\_40Hz Results}
\label{subsubsec:fad_results}

Table~\ref{tab:fad_results} presents FAD\_40Hz performance.

\begin{table}[h!]
\centering
\caption{FAD\_40Hz: mean performance across subjects.  Format as in
  Table~\ref{tab:mobifall_results}.}
\label{tab:fad_results}
\setlength{\tabcolsep}{3pt}\footnotesize
\begin{tabular}{llccccc}
\toprule
\textbf{Classifier} & \textbf{Protocol}
  & \textbf{Acc} & \textbf{Rec} & \textbf{Pre}
  & \textbf{$F_1$} & \textbf{$F_2$} \\
\midrule
\multirow{3}{*}{LR}
 & Naive
   & 0.9802 & 0.9918 & 0.8145 & 0.8941 & 0.9502 \\
 & General
   & $0.9826 \pm 0.0042$ & $0.9929 \pm 0.0065$ & ---
   & $0.8856 \pm 0.0653$ & $0.9429 \pm 0.0363$ \\
 & Personal
   & $0.9878 \pm 0.0032$ & $0.9235 \pm 0.0450$ & ---
   & $0.9111 \pm 0.0457$ & $0.9168 \pm 0.0454$ \\
\midrule
\multirow{3}{*}{SVM}
 & Naive
   & 0.9857 & 0.9887 & 0.8624 & 0.9209 & 0.9603 \\
 & General
   & $0.9890 \pm 0.0027$ & $0.9878 \pm 0.0109$ & ---
   & $0.9236 \pm 0.0457$ & $0.9592 \pm 0.0234$ \\
 & Personal
   & $0.9839 \pm 0.0035$ & $0.8721 \pm 0.0771$ & ---
   & $0.8778 \pm 0.0592$ & $0.8724 \pm 0.0696$ \\
\bottomrule
\end{tabular}
\end{table}

On FAD\_40Hz, SVM outperformed LR in the General protocol ($F_2 = 0.9592$ vs.\ $0.9429$), reversing the ranking seen on MobiFall and SisFall. Recall remained near perfect for both classifiers ($\Rec \geq 0.987$). The wide $F_1$ and $F_2$ confidence intervals in the General protocol reflect high variability in precision across subjects, even though recall was stable --- the model consistently detects falls but the false-alarm rate differs substantially from one subject to the next.

\subsection{Cross-Dataset Performance Summary}
\label{subsubsec:summary}

Table~\ref{tab:summary} collects the mean General $F_2$ for all datasets and classifiers in a single view.

\begin{table}[h!]
\centering
\caption{Subject-level mean $F_2$ under the General (LOSO) protocol.
  Values are mean $\pm$ 95\,\% CI.}
\label{tab:summary}
\begin{tabular}{lcc}
\toprule
\textbf{Dataset} & \textbf{LR $F_2$} & \textbf{SVM $F_2$} \\
\midrule
MobiFall   & $0.9049 \pm 0.0393$ & $0.8918 \pm 0.0327$ \\
SisFall    & $0.9709 \pm 0.0108$ & $0.9700 \pm 0.0111$ \\
FAD\_40Hz  & $0.9429 \pm 0.0363$ & $0.9592 \pm 0.0234$ \\
\bottomrule
\end{tabular}
\end{table}

SisFall ranked highest for both classifiers, followed by FAD\_40Hz and then MobiFall. The relative ordering was consistent across classifiers, but the preferred classifier changed by dataset: LR was better on MobiFall and SisFall, while SVM was better on FAD\_40Hz. This suggests that the choice of classifier is secondary to dataset geometry and size, consistent with the representation-dominance hypothesis.

% TODO: [Comparison with published baseline methods to be inserted once
%  literature survey is complete --- see proposal Section 5.4 for the
%  planned baseline families (threshold-based, statistical-feature,
%  spectral-feature).]

\subsection{Effect of Decision Threshold Optimisation}
\label{subsubsec:threshold_effect}

For each held-out subject under the General protocol the decision threshold was swept between 0.05 and 0.95 (19 values). Table~\ref{tab:threshold_gains} shows the average per-subject gain in $F_2$ obtained by choosing the subject-specific optimal threshold rather than the trained default $\hat\tau_0(s)$ (the threshold selected by the PR-curve procedure on the corresponding training fold; see Section~\ref{subsubsec:threshold} of the Methodology).

\begin{table}[h!]
\centering
\caption{Average subject-level $F_2$ gain from threshold optimisation.
  $\Delta F_2(s) = F_2(\tau^*(s)) - F_2(\hat\tau_0(s))$ where $\hat\tau_0(s)$ is the
  PR-curve trained threshold for subject $s$; values averaged over held-out subjects.}
\label{tab:threshold_gains}
\small
\begin{tabular}{lcc}
\toprule
\textbf{Dataset} & \textbf{$\Delta F_2$ (LR)} & \textbf{$\Delta F_2$ (SVM)} \\
\midrule
MobiFall   & $+0.0348$ & $+0.0474$ \\
SisFall    & $+0.0079$ & $+0.0042$ \\
FAD\_40Hz  & $+0.0250$ & $+0.0121$ \\
\bottomrule
\end{tabular}
\end{table}

The gains were largest on MobiFall --- where the default threshold was furthest from optimal --- and smallest on SisFall, where the default was already near-optimal. The gains were not uniform across subjects within any dataset. On FAD\_40Hz, subject 11 showed the largest gain observed in any experiment: $\Delta F_2 = +0.1165$ for LR, lifting that subject's score from $0.8046$ to $0.9211$ by adjusting the threshold alone, without retraining the underlying model. This illustrates that a single shared threshold cannot account for the variation in decision boundary location across individuals, and that personalised threshold calibration can meaningfully recover performance for outlier subjects.

\paragraph{Computational cost of threshold personalisation.}
The trained threshold $\hat\tau_0$ enters only the final thresholding step
and plays no role in feature extraction, classifier training, or Platt calibration.
Adjusting $\hat\tau_0$ for a held-out subject therefore requires \emph{no model retraining}:
the probability scores $\{\hat{p}_i\}$ are computed once under the trained model and
the sweep evaluates $F_2(\tau,s)$ for each $\tau\in\mathcal{T}$ by a single pass over
the precomputed score vector.  The additional computational overhead per subject is
$O\bigl(|\mathcal{T}|\cdot|W_s|\bigr)$ comparisons, where $|\mathcal{T}|=19$ and
$|W_s|$ denotes the per-subject test-window count.  This is negligible against the
Bayesian search cost, which scales as
$O\bigl(T\cdot k\cdot N_{\mathrm{tr}}\cdot C_{\mathrm{PH}}\bigr)$
where $T=200$, $k$ is the number of inner-CV folds, $N_{\mathrm{tr}}$ is the training
window count, and $C_{\mathrm{PH}}$ is the per-window persistent homology cost.
On the UHeM cluster the Bayesian search consumed approximately 75~minutes per dataset;
the full threshold sweep over all subjects in a dataset ran in under one second of wall
time---a runtime ratio exceeding $10^3$.  This asymmetry makes threshold calibration an
attractive lightweight personalisation mechanism: it is applicable at inference time
from any small per-subject calibration set, incurs no retraining cost, and requires no
access to the original training data.

\paragraph{Statistical analysis of threshold heterogeneity.}
The per-subject gains $\Delta F_2(s)$ exhibit substantial inter-subject variability within
each dataset.  On FAD\_40Hz the subject-level range extends from near-zero to
$+0.1165$ (Subject~11, LR), and on MobiFall the inter-subject standard deviation of
$\Delta F_2$ is comparable to the dataset mean gain for SVM.  A purely mathematical
argument for threshold personalisation would predict a uniform benefit: if the posterior
probability function $f_\theta\circ\phi_{\lambda^*}$ is systematically miscalibrated in
the same direction for all subjects, a single shared adjustment of $\hat\tau_0$
would lift all subjects equally.  The observed heterogeneity refutes this: subjects with
near-zero gain already sit at an optimal threshold close to $0.50$, while subjects with
large gain sit far from it---a pattern inconsistent with any single global miscalibration.
The statistical evidence is that the \emph{distribution} of $\Delta F_2(s)$ is
right-skewed with a heavy upper tail (cf.\ Subject~11), and the inter-subject standard
deviation is not small relative to the mean.  Mathematically, this implies that
$\tau^*(s) = \argmax_{\tau\in\mathcal{T}}F_2(\tau,s)$ is a genuinely subject-specific
quantity: the geometry of the decision boundary $\{f_\theta(\phi_{\lambda^*}(w))=\tau\}$
varies across individuals in a way that the shared classifier $f_\theta$ cannot
fully absorb.  This is the statistical signature of Hypothesis~4 in the proposal and
motivates the distributional reporting of $\Delta F_2$ in the Results \& Discussion
section.

% TODO: [A per-subject $\Delta F_2$ distribution figure (one boxplot panel
%  per dataset) would visualise the heterogeneity of threshold gains and
%  support the personalisation hypothesis.]

% ═════════════════════════════════════════════════════════════
\section{Statistical Inference}
\label{subsec:stat_inference}

Section~\ref{subsubsec:stat-inference} of the Methodology defines, in general and dataset-independent terms, the two sources of variability tracked throughout this thesis (repetition-level and subject-level), the repetition-level $t$-confidence interval of Theorem~\ref{thm:t-ci}, subject-level paired comparisons, the percentile bootstrap, and the paired sign-flipping permutation test. This section instantiates that machinery with the concrete settings used for every dataset in this thesis, then applies it to two comparisons specific to the experimental design: cross-dataset transplantation of hyperparameters and the distribution of personalisation gains.

Every repetition-level interval below uses $R=15$ repetitions of the class-balancing procedure, giving $t_{0.025,14}\approx 2.145$ in~\eqref{eq:t-ci} at the standard $\alpha=0.05$ level. Every bootstrap interval uses $B=10{,}000$ resamples at the same $\alpha=0.05$. Every permutation test is computed exactly when $|\Sset_K|\leq 24$ (true for all three datasets) and would otherwise fall back to a $10{,}000$-sample Monte Carlo estimate. Subject-level averages, dispersion, and these uncertainty estimates take precedence over headline maxima throughout, for the reason given in Remark~\ref{rem:headline-bias}.

\subsection{Distributional analysis of personalisation}
\label{subsubsec:dist-pers}

The personalisation question requires a distributional rather than a single-number answer. The per-subject gain is
\[
  \Delta_{K,s}=F_{2,K,s}^{\mathrm{pers}}-F_{2,K,s}^{\mathrm{gen}}.
\]
The thesis reports the four summaries
\[
  \overline{\Delta}_K,
  \qquad
  \mathrm{med}(\Delta_{K,s}),
  \qquad
  \Var(\Delta_{K,s}),
  \qquad
  \widehat{\Prob}(\Delta_{K,s}>0)
\]
together with bootstrap percentile confidence intervals for each. The joint pattern of these four quantities takes precedence over any single summary, because personalisation is interpreted as a heterogeneous subject-level intervention rather than a constant treatment effect.

\subsection{Cross-applied hyperparameters}
\label{subsubsec:cross-applied}

To evaluate dataset-specific topological geometry, optimal tuples are cross-applied across datasets. Define the cross-applied estimator
\[
  \widehat F_2(K\mid\lambda^*_{K'})=\E\bigl[F_2\mid K,\lambda^*_{K'}\bigr],
\]
where the expectation is taken over the LOSO folds of dataset $K$ with the hyperparameter tuple optimised on a different dataset $K'\neq K$. The off-diagonal degradations
\[
  \mathrm{Drop}_{K\to K'}=\widehat F_2(K\mid\lambda^*_K)-\widehat F_2(K\mid\lambda^*_{K'})
\]
quantify the cost of transplanting a hyperparameter tuple from $K'$ to $K$ and are reported with bootstrap percentile intervals.

% ═════════════════════════════════════════════════════════════
\section{Computational Infrastructure}
\label{subsec:compute}

All experiments were executed on the National HPC Centre (UHeM, Istanbul), allocation 4025462026, under the SLURM workload manager. Compute nodes were equipped with dual Intel Xeon Gold 6148 processors (2.40~GHz base clock, 20 physical cores per socket, 40 cores per node) and 190~GB of RAM. No GPU acceleration was used: the persistent homology reduction of Algorithm~\ref{alg:reduction} is intrinsically sequential at the point cloud sizes considered, and the gain of porting to GPU is dominated by the constant factor of the optimised CPU implementation.

The Bayesian optimisation ran sequentially across $T=200$ TPE trials per dataset, with $15$ parallel inner-CV workers per trial. Across the three datasets the total wall-clock time was approximately $3$~h $45$~min, with peak memory usage close to $10$~GB per node. Persistent homology was computed with the GUDHI library~\parencite{gudhi}; classification models were implemented in scikit-learn~\parencite{scikit-learn}. All randomness was seeded for reproducibility, and the seed list, the search history, and the resulting per-subject scores are archived alongside the thesis source.


% ============================================================
%  END ch4.tex
% ============================================================

% ============================================================
%  BEGIN ch5.tex
% ============================================================
% ─────────────────────────────────────────────────────────────
\phantomsection
\chapter{Results and Discussion}
\label{sec:results}

This section evaluates the four research hypotheses (H1--H4) introduced in the
Methodology against the experimental outcomes reported in the Experiments section.
The \emph{General protocol} (leave-one-subject-out, LOSO) is the primary evaluation
criterion; the \emph{Naive} and \emph{Personal} protocols serve as an upper-bound
reference and an idealised personalisation scenario, respectively.  The
decision-threshold sweep provides a retrospective diagnostic of subject-level
heterogeneity.

All numerical claims in this section refer to the General-protocol results in
Tables~4--9 of the Experiments section.  The F$_2$ score is the headline metric
throughout; recall and accuracy are reported as supporting evidence.

% ─────────────────────────────────────────────────────────────
\section{H1 — Representation Effectiveness}
\label{subsec:H1}

\textbf{Hypothesis.} There exist hyperparameters $\lambda\in\Lambda$ and classifiers
$f_\theta$ such that the mean subject-general $F_2$ is at least $0.93$ on SisFall and
FAD\_40Hz and at least $0.80$ on MobiFall, and the mean subject-general recall is high
across all three datasets.

\subsection{F$_2$ outcomes}

Table~1 of the cross-dataset summary (Experiments, Tab.~9) gives the following
General-protocol results.

\begin{table}[h!]
\centering
\caption{General-protocol $F_2$ outcomes against H1 targets.  CI denotes
  95\,\% repetition-level confidence interval ($R=15$).}
\label{tab:H1-summary}
\small
\begin{tabular}{llccc}
\toprule
\textbf{Dataset} & \textbf{Clf} & \textbf{$F_2$} & \textbf{CI} & \textbf{H1 target met?} \\
\midrule
SisFall   & LR  & 0.9709 & $\pm0.0108$ & Yes ($\geq0.93$) \\
SisFall   & SVM & 0.9700 & $\pm0.0111$ & Yes              \\
FAD\_40Hz & LR  & 0.9429 & $\pm0.0363$ & Marginal         \\
FAD\_40Hz & SVM & 0.9592 & $\pm0.0234$ & Yes              \\
MobiFall  & LR  & 0.9049 & $\pm0.0393$ & Yes ($\geq0.80$) \\
MobiFall  & SVM & 0.8918 & $\pm0.0327$ & Yes              \\
\bottomrule
\end{tabular}
\end{table}

SisFall achieves the highest absolute $F_2$, exceeding the $0.93$ threshold for both
classifiers with a narrow confidence interval ($\pm0.011$), indicating low
repetition-to-repetition variability.  FAD\_40Hz meets the target with SVM ($F_2 =
0.9592$); the LR value of $0.9429\pm0.0363$ lies marginally below $0.93$ at the point
estimate but the lower confidence bound ($0.9066$) suggests that performance is not
materially below target.  MobiFall clearly satisfies the weaker $0.80$ criterion for
both classifiers ($0.9049$ and $0.8918$, respectively).

The recall values under the General protocol reinforce the recall-oriented design:
SisFall LR achieves $\mathrm{Rec}=0.9836$, FAD\_40Hz LR achieves
$\mathrm{Rec}=0.9929$, and MobiFall LR achieves $\mathrm{Rec}=0.9574$.  No classifier
on any dataset falls below $0.95$ recall under the General protocol.  This is a direct
consequence of the $F_2$ weighting: Proposition~2 of the Methodology states that the
marginal $F_2$ gain per unit of recall is four times the gain per unit of precision near
$\mathrm{Rec}=\mathrm{Prec}$, so the Bayesian optimisation is systematically steered
toward high-recall configurations.

\begin{figure}[!ht]
\centering
\includegraphics[width=\textwidth]{figures/fig_f2_bars}
\caption{Mean $F_2$ by dataset, classifier (LR/SVM), and evaluation protocol (Naive,
  General, Personal).  Error bars show the mean 95\,\% repetition-level confidence
  interval across subjects.  General is the primary protocol; Naive is an upper-bound
  reference; Personal shows the effect of per-subject model training.  MobiFall
  Personal collapses due to the small subject pool ($|\mathcal{S}|=9$).}
\label{fig:f2-bars}
\end{figure}

\subsection{Naive vs.\ General protocol gap}

Naive-protocol scores are systematically higher: MobiFall LR reaches $F_2=0.8857$
(Naive) versus $F_2=0.9049$ (General).  Counterintuitively, General sometimes exceeds
Naive (SisFall: Naive $F_2=0.9713$ vs.\ General $0.9709$; FAD: Naive $0.9502$ vs.\
General $0.9429$).  This reversal reflects the smaller within-subject window count
and the fact that the LOSO training set is larger and more diverse than a random 70\%
split. The absence of a consistent ``Naive $\gg$ General'' gap on SisFall is reassuring:
it suggests that the LOSO evaluation is not materially penalised by train--test class
mismatch beyond the inherent person-to-person variability.

\subsection{Mathematical basis}

The feature map $\phi_\lambda\colon\R^{n\times 3}\to\R^{24}$ extracts persistence
statistics from the delay-embedded signal.  The Cohen--Steiner stability theorem
guarantees that
\[
  d_B(\phi_\lambda(w),\,\phi_\lambda(w'))\;\leq\;\|\mathrm{SMV}(w)-\mathrm{SMV}(w')\|_\infty,
\]
so that small sensor noise perturbations of a window $w$ produce at most equally
small perturbations of its topological descriptor.  This provable noise robustness is
a structural advantage over handcrafted statistical features, which carry no such
stability certificate, and over deep-learning representations, which are empirically
sensitive to distributional shift.

\subsection{Limitation: absence of baseline comparison}

H1 is conditionally validated in the sense that the system meets its own pre-specified
performance targets.  An absolute claim---``TDA outperforms existing methods''---requires
a comparison against representative baselines: fixed-threshold methods, statistical-feature
pipelines, and spectral-feature pipelines, as planned in the proposal ($\S$5.4).  These
experiments have not been conducted.  The conclusion that TDA is competitive is therefore
conditional on the assumption that the targets themselves are meaningful.  Section~\ref{subsec:limitations}
discusses this gap and its implications.

% ─────────────────────────────────────────────────────────────
\section{H2 — Dataset-Specific Topological Geometry}
\label{subsec:H2}

\textbf{Hypothesis.}  The optimal hyperparameter tuple $\lambda^*_K$ depends on the
dataset $K$.  At least two of the three datasets have distinct optimisers, and
transplanting an optimiser from one dataset to another incurs a measurable drop in
General-protocol $F_2$.

\subsection{Distinct optimal configurations}

Table~5 of the Experiments (Tab.~\ref{tab:optimal_params_ref} below, reproduced for
convenience) shows that all three datasets converged to different complex types,
window lengths, and distance measures.

\begin{table}[h!]
\centering
\caption{Optimal TDA configurations $\lambda^*_K$ per dataset (from Experiments
  Tab.~\ref{tab:optimal_params}).  A dash (---) indicates parameters not applicable for
  the Alpha complex.}
\label{tab:optimal_params_ref}
\small
\begin{tabular}{lccc}
\toprule
\textbf{Parameter} & \textbf{MobiFall} & \textbf{SisFall} & \textbf{FAD\_40Hz} \\
\midrule
Best $F_2$ (Optuna objective)  & 0.9746     & 0.9733   & 0.9860      \\
Window length (s)              & 5.0        & 1.0      & 1.0         \\
Complex type                   & SparseRips & Alpha    & Rips        \\
Embedding dimension            & 5          & 5        & 3           \\
Embedding delay                & 2          & 4        & 2           \\
Stride factor                  & 2          & 4        & 1           \\
Distance measure               & Chebyshev  & ---      & Manhattan   \\
Filtration scale               & 80\,\%     & ---      & 40\,\%      \\
\bottomrule
\end{tabular}
\end{table}

Every hyperparameter dimension exhibits variation across datasets.  The complex type
differs across all three: SparseRips for MobiFall, Alpha for SisFall, and standard
Rips for FAD\_40Hz.  The window length separates MobiFall (5.0~s) from the other two
(1.0~s), and the distance measures are distinct where applicable (Chebyshev vs.\
Manhattan).  H2 is qualitatively confirmed: no single universal configuration exists
for these three datasets.

\subsection{Interpretation of configuration differences}

\paragraph{Window length.}
MobiFall records both a pre-fall phase (normal activity) and a post-fall phase (lying),
with the impact event compressed into a brief sub-second spike.  A 5.0~s window captures
the full topology of the trajectory: the pre-impact cloud geometry, the sharp impulsive
excursion, and the post-fall stationary phase, which are each topologically distinguishable.
SisFall and FAD\_40Hz subjects execute controlled laboratory falls that are sufficiently
discriminable from ADL in a 1.0~s window, suggesting less temporal spread in the
distinguishing geometry.

\paragraph{Complex type.}
The Alpha complex is the natural choice when the point cloud geometry is close to
Euclidean: its filtration is the Čech filtration restricted to the Delaunay triangulation,
which SisFall's well-structured sensor data supports (24 older adults executing 15 fall
types under controlled conditions produce geometrically regular clouds).  SparseRips
trades exactness for efficiency by subsampling the cloud; it is appropriate when the
cloud is large and noisy, which is the case for MobiFall's continuous-wear sensor data.
Standard Rips uses all pairwise distances and suits the moderate-size, moderately-noisy
FAD\_40Hz clouds.

\paragraph{Distance measure.}
Chebyshev distance ($\ell_\infty$) emphasises the maximum coordinate deviation and is
sensitive to sharp acceleration spikes---a characteristic feature of the fall impact in
MobiFall.  Manhattan distance ($\ell_1$) reduces the influence of outliers and is
appropriate for FAD\_40Hz, whose 40~Hz sampling produces sparser, more variable clouds
than the implicit 50~Hz resampled signals.  The Alpha complex on SisFall does not
require a user-specified distance because it is defined by the Euclidean Delaunay
triangulation.

\subsection{Sensitivity of $F_2$ to $\lambda$}

Although cross-application drops $\mathrm{Drop}_{K\to K'}$ were not computed (see
Section~\ref{subsec:limitations}), the Bayesian search trial histories provide indirect
evidence for sensitivity.  Across 200 trials per dataset, the $F_2$ values ranged from
$0.7958$ to $0.9746$ on MobiFall (range 0.179), from $0.9198$ to $0.9733$ on SisFall
(range 0.054), and from $0.8785$ to $0.9860$ on FAD\_40Hz (range 0.108).  The large
trial-to-trial variation on MobiFall and FAD implies that a mis-specified $\lambda$
(such as a configuration optimised for a different dataset) would very likely incur
a substantial $F_2$ penalty.  The narrow trial range on SisFall suggests that SisFall
is less sensitive to $\lambda$, consistent with the clean Euclidean geometry and the
strong performance even at the worst trial ($F_2=0.9198$).

\begin{figure}[!ht]
\centering
\includegraphics[width=0.75\textwidth]{figures/fig_bayesian_convergence}
\caption{Bayesian optimisation convergence: best-so-far averaged $F_2$ proxy versus
  Optuna trial number (200 trials per dataset, TPE sampler).  FAD\_40Hz achieves the
  highest proxy ($0.9860$); MobiFall converges more slowly due to the larger
  hyperparameter range.  The wide trial range on MobiFall (0.179) and FAD (0.108)
  relative to SisFall (0.054) is direct evidence that $\lambda$ sensitivity differs
  across datasets (H2).}
\label{fig:bayesian-convergence}
\end{figure}

% ─────────────────────────────────────────────────────────────
\section{H3 — Representation Dominates Classifier Choice}
\label{subsec:H3}

\textbf{Hypothesis.}  After tuning the topological representation, the $F_2$ gap
between the best LR model and the best SVM model is small relative to the $F_2$
variation attributable to the choice of $\lambda$.  Most of the predictive signal
resides in $\phi_\lambda$, not in the classifier family.

\subsection{Classifier gap}

Under the General protocol, the absolute difference $|F_2^{\mathrm{LR}} -
F_2^{\mathrm{SVM}}|$ is:
\begin{itemize}
  \item MobiFall: $|0.9049 - 0.8918| = 0.0131$
  \item SisFall:  $|0.9709 - 0.9700| = 0.0009$
  \item FAD\_40Hz: $|0.9429 - 0.9592| = 0.0163$
\end{itemize}
All three gaps are below $0.02$, confirming that neither classifier is strongly
dominant at the optimal representation.  On MobiFall and SisFall, LR is marginally
better; on FAD\_40Hz, SVM is marginally better.  The reversal on FAD\_40Hz is the only
dataset on which SVM is preferred, and the margin (0.0163) is within the confidence
interval overlap ($\pm0.0363$ for LR, $\pm0.0234$ for SVM).

\subsection{$\lambda$-tuning variation}

Table~\ref{tab:H3-comparison} quantifies the variation in $F_2$ attributable to the
choice of $\lambda$ using the Optuna trial histories (200 trials per dataset).

\begin{table}[h!]
\centering
\caption{Classifier gap vs.\ $\lambda$-tuning variation.  The trial range is
  $F_2^{\max} - F_2^{\min}$ across all 200 Optuna trials.  The ratio is
  $\text{range}/\text{gap}$.}
\label{tab:H3-comparison}
\small
\begin{tabular}{lccc}
\toprule
\textbf{Dataset} & \textbf{Classifier gap} & \textbf{Trial range} & \textbf{Ratio} \\
\midrule
MobiFall  & 0.0131 & 0.1788 & $\mathbf{13.7\times}$ \\
SisFall   & 0.0009 & 0.0534 & $\mathbf{59\times}$   \\
FAD\_40Hz & 0.0163 & 0.1075 & $\mathbf{6.6\times}$  \\
\bottomrule
\end{tabular}
\end{table}

The classifier gap is between $6.6\times$ and $59\times$ smaller than the variation
induced by choosing a poor $\lambda$.  Selecting the wrong topological configuration
degrades $F_2$ by up to $0.18$ units; switching classifier families at the optimal
configuration changes $F_2$ by at most $0.016$ units.

\subsection{Mathematical interpretation}

The feature map $\phi_{\lambda^*}\colon\R^{n\times 3}\to\R^{24}$ with the optimised
hyperparameters creates a representation in which the fall vs.\ ADL decision problem
is approximately linearly separable.  In such a regime, both logistic regression (which
fits a linear decision boundary directly) and SVM with a calibrated kernel (which finds
a maximum-margin linear separator or uses a kernel to lift to higher dimensions) operate
near their theoretical limit.  The near-equivalence of LR and SVM confirms that the
nonlinear geometric complexity of the fall detection problem has been absorbed by
$\phi_{\lambda^*}$, leaving behind a well-conditioned linear classification task.

The lone exception---SVM$>$LR on FAD\_40Hz---is consistent with this interpretation:
FAD\_40Hz has the smallest per-subject window count and the widest LR confidence
interval, suggesting residual complexity that the RBF kernel can partially resolve.
Even so, the margin is below $0.02$.

% ─────────────────────────────────────────────────────────────
\section{H4 — Heterogeneous Threshold Personalisation}
\label{subsec:H4}

\textbf{Hypothesis.}  The subject-level sweep gain
$\Delta F_2(s) = F_2(\tau^*(s),s) - F_2(\hat\tau_0(s),s)$
is heterogeneous across subjects within each dataset, and the mean gain
$\overline{\Delta F_2}$ differs across datasets.

\subsection{Per-subject sweep gain distribution}
\label{subsubsec:sweep-dist}

Table~\ref{tab:sweep-dist} reports the distribution of $\Delta F_2(s)$ computed from
the threshold sweep stored in the Results\_V15 database.  The baseline
$\hat\tau_0(s)$ is the PR-curve optimal trained threshold for subject $s$
(typically $\hat\tau_0\in[0.55,0.88]$), so the gain measures improvement over an
already-calibrated threshold rather than over the uninformative $\tau=0.50$.

\begin{table}[h!]
\centering
\caption{Per-subject threshold sweep gain $\Delta F_2(s)$ statistics.
  \emph{Near-zero}: $|\Delta F_2|\leq0.001$.
  All gains are non-negative by construction ($\tau^*$ is the sweep argmax).}
\label{tab:sweep-dist}
\small
\begin{tabular}{llccccc}
\toprule
\textbf{Dataset} & \textbf{Clf} & $n$ & $\overline{\Delta F_2}$ &
  $\sigma_{\Delta}$ & $\max\,\Delta F_2$ & Near-zero \\
\midrule
MobiFall   & LR  &  9 & 0.0348 & 0.0186 & 0.0615 (s8)  & 1/9  \\
MobiFall   & SVM &  9 & 0.0474 & 0.0350 & 0.1013 (s9)  & 1/9  \\
SisFall    & LR  & 24 & 0.0079 & 0.0109 & 0.0527 (s29) & 4/24 \\
SisFall    & SVM & 24 & 0.0042 & 0.0047 & 0.0169 (s29) & 9/24 \\
FAD\_40Hz  & LR  & 13 & 0.0250 & 0.0323 & 0.1165 (s11) & 2/13 \\
FAD\_40Hz  & SVM & 13 & 0.0121 & 0.0172 & 0.0691 (s11) & 2/13 \\
\bottomrule
\end{tabular}
\end{table}

The mean gains in Table~\ref{tab:sweep-dist} reproduce the dataset averages reported
in Experiments Tab.~8 ($+0.0348/+0.0474$ for MobiFall, $+0.0079/+0.0042$ for SisFall,
$+0.0250/+0.0121$ for FAD\_40Hz).  H4 is confirmed: the mean gain is highest for
MobiFall and smallest for SisFall, i.e., the mean differs across datasets.

\begin{figure}[!ht]
\centering
\includegraphics[width=\textwidth]{figures/fig_sweep_boxplot}
\caption{Per-subject threshold sweep gain $\Delta F_2(s) = F_2(\tau^*(s)) -
  F_2(\hat\tau_0(s))$ for LR (left) and SVM (right).  Jittered points show
  individual subjects; box shows quartiles.  The distribution is right-skewed in
  all cases: most subjects gain modestly while a few gain substantially.  FAD\_40Hz
  subject~11 achieves the largest gain (+0.1165, LR), visible as the highest point
  in the FAD panel.  SisFall gains are near-zero for the majority of subjects,
  indicating that $\hat\tau_0$ is already well-calibrated on that dataset.}
\label{fig:sweep-boxplot}
\end{figure}

\paragraph{Distribution shape.}
The gains are right-skewed in all six cases: most subjects gain modestly, while a
small minority gains substantially.  On FAD\_40Hz, subject~11 gains $\Delta F_2 =
+0.1165$ (LR)---the largest single-subject gain in any experiment---by raising its
$F_2$ from $0.8046$ to $0.9211$ through threshold adjustment alone.  On MobiFall,
subject~9 gains $+0.1013$ (SVM).  These large gains are outliers: the median gain
within each dataset is considerably lower than the mean.

\paragraph{Near-zero gainers.}
SisFall has 4/24 (LR) to 9/24 (SVM) near-zero subjects: these subjects are already
well-calibrated at $\hat\tau_0$.  Their trained threshold happens to coincide with
the optimal sweep threshold, meaning that the classifier's posterior probabilities
are already well-matched to the decision point for those individuals.  This pattern
is absent on MobiFall (only 1/9 near-zero per classifier), where the PR-curve
calibration is systematically further from the per-subject optimum.

\paragraph{Heterogeneity metric.}
The coefficient of variation ($\sigma_\Delta / \overline{\Delta F_2}$) measures
relative heterogeneity: MobiFall LR $=53\%$, MobiFall SVM $=74\%$, SisFall
LR $=138\%$, FAD\_40Hz LR $=129\%$.  All values exceed $50\%$, confirming that
the gain distribution is not concentrated.  A single global threshold offset
cannot replicate this distribution: if every subject were miscalibrated in the
same direction by the same amount, the coefficient of variation would approach
zero.  The high coefficients of variation are the statistical signature that
$\tau^*(s)$ is a genuinely subject-specific quantity.

\paragraph{Mathematical explanation.}
The decision boundary of $f_\theta\circ\phi_{\lambda^*}$ is the hyperplane
$\{x\in\R^{24}\colon f_\theta(x)=\hat\tau_0\}$.  Each subject's test windows form
a cloud in $\R^{24}$ whose local geometry around this hyperplane differs from subject
to subject: some subjects' fall windows sit close to the boundary (requiring a lower
threshold to capture them), while others' fall windows are well-separated (the default
threshold is sufficient).  The threshold sweep detects this per-subject boundary
geometry retrospectively.  A prospective personalisation scheme would use a small
held-out calibration set per subject to identify $\tau^*(s)$ before deployment
(see Section~\ref{subsec:deployment}).

\paragraph{Engineering contribution.}
This personalisation is achieved without any model retraining.  The probability
scores $\{\hat p_i\}$ are computed once under the trained model; the sweep evaluates
$F_2(\tau, s)$ by a single pass over the precomputed score vector for each $\tau
\in\T$.  The additional computational cost is $O(|\T|\cdot|W_s|)$ per subject,
where $|\T|=19$ and $|W_s|$ is the per-subject window count---less than one second
of wall time for all subjects in any dataset.  This is a runtime ratio exceeding
$10^3$ compared to the ${\sim}75$~min Bayesian search (Experiments,
Section~4.3.4).

\subsection{Personal training protocol}
\label{subsubsec:personal-protocol}

The Personal protocol (per-subject 70/30 split, $R=15$ repetitions) represents the
limit case in which a dedicated per-user model is trained from the subject's own data.
Table~\ref{tab:personal-delta} reports the aggregate personal--general $F_2$ difference.

\begin{table}[h!]
\centering
\caption{Personal vs.\ General protocol aggregate $F_2$.
  $\Delta = F_2^{\mathrm{pers}} - F_2^{\mathrm{gen}}$ (negative = General is better).}
\label{tab:personal-delta}
\small
\begin{tabular}{llccc}
\toprule
\textbf{Dataset} & \textbf{Clf} & $F_2^{\mathrm{gen}}$ & $F_2^{\mathrm{pers}}$ & $\Delta$ \\
\midrule
MobiFall  & LR  & 0.9049 & 0.7325 & $-0.172$ \\
MobiFall  & SVM & 0.8918 & 0.5981 & $-0.294$ \\
SisFall   & LR  & 0.9709 & 0.9405 & $-0.030$ \\
SisFall   & SVM & 0.9700 & 0.9196 & $-0.050$ \\
FAD\_40Hz & LR  & 0.9429 & 0.9168 & $-0.026$ \\
FAD\_40Hz & SVM & 0.9592 & 0.8724 & $-0.087$ \\
\bottomrule
\end{tabular}
\end{table}

At the aggregate level, General $>$ Personal on all datasets and both classifiers.
The magnitude of the shortfall increases dramatically as the subject pool shrinks:

\begin{itemize}
\item \textbf{SisFall ($|\Sset|=24$)}: modest drop ($\Delta \approx -0.03$ to
  $-0.05$).  Twenty-four subjects contribute enough per-subject windows
  (${\sim}50$--$150$ falls per subject) for a reliable within-subject model.
\item \textbf{FAD\_40Hz ($|\Sset|=13$)}: moderate drop ($\Delta \approx -0.03$
  to $-0.09$).  Fewer subjects and fewer falls per subject; personal models are
  noisier.
\item \textbf{MobiFall ($|\Sset|=9$)}: severe drop ($\Delta=-0.172$ for LR;
  $\Delta=-0.294$ for SVM).  Nine subjects with $\lesssim30$ fall instances per
  subject is insufficient for stable personal model estimation.  The SVM personal
  recall collapses to $0.5822$, indicating that the personal classifier is barely
  above chance on many subjects.
\end{itemize}

\begin{figure}[!ht]
\centering
\includegraphics[width=0.90\textwidth]{figures/fig_general_vs_personal}
\caption{General (LOSO) vs Personal (per-subject training) $F_2$ per subject, for
  LR (left) and SVM (right).  Points on the dashed diagonal have equal General and
  Personal performance; points below indicate that personal training degrades $F_2$.
  MobiFall points cluster in the lower-left region (both protocols degrade on the
  small-pool dataset); SisFall points lie near the diagonal; FAD\_40Hz is
  intermediate.  The shaded region marks Personal $<$ General.}
\label{fig:general-vs-personal}
\end{figure}

This pattern reveals a critical engineering constraint: personal model training is
only beneficial when per-subject calibration data is abundant.  In typical
deployment scenarios (a user who has experienced very few confirmed falls), full
personal training is counterproductive.  The threshold sweep---which requires no
retraining and adds negligible computational cost---is the appropriate
personalisation mechanism when calibration data is sparse.  The mean sweep gain for
MobiFall LR ($+0.0348$) and SVM ($+0.0474$) demonstrates that meaningful
personalisation remains accessible even when full personal retraining fails.

% ─────────────────────────────────────────────────────────────
\section{Ablation Study: Feature-Block Contributions}
\label{subsec:ablation}

\textbf{Motivation.}  The $24$-dimensional representation $\phi_{\lambda^*}$ is the
concatenation of three semantically distinct blocks: ten statistics of the
dimension-zero persistence diagram (\emph{H0}, summarising the persistence of
connected components of the delay-embedded point cloud), ten statistics of the
dimension-one diagram (\emph{H1}, summarising loops), and four signal-level
descriptors of the raw acceleration magnitude (\emph{Signal}: maximum, standard
deviation, mean, and range).  The headline results of Section~\ref{subsec:H1}
establish that the full vector is effective but do not reveal \emph{which} block
carries the discriminative power.  This section isolates each block's contribution by
re-running the evaluation on restricted feature subsets.

\textbf{Protocol.}  Each subset is evaluated under the identical General (LOSO)
protocol used throughout: a single global grid search selects the classifier
hyperparameters on balanced-undersampled data, and for each left-out subject the model
is refitted over $15$ repetitions with a precision--recall threshold chosen on the
training fold; the reported $F_2$ is the cross-subject mean with a $95\,\%$ $t$-confidence
interval.  Only the feature columns presented to the classifier change; the topological
configuration $\lambda^*$ and the delay embedding are held fixed.  As a validation
check, the full-vector rows of Table~\ref{tab:ablation_blocks} reproduce the headline
General-protocol $F_2$ of Section~\ref{subsec:H1} to within the confidence interval for
both classifiers on all three datasets.

\begin{table}[h!]
\centering
\caption{Feature-block ablation under the General (LOSO) protocol.  Entries are
  cross-subject mean $F_2 \pm 95\,\%$ CI.  Block dimensions: Full~$=24$, H0~$=10$,
  H1~$=10$, H0+H1~$=20$, Signal~$=4$.  The Full column reproduces the headline results
  of Tables~\ref{tab:mobifall_results}--\ref{tab:fad_results}.}
\label{tab:ablation_blocks}
\setlength{\tabcolsep}{3.5pt}\footnotesize
\begin{tabular}{llccccc}
\toprule
\textbf{Dataset} & \textbf{Clf}
  & \textbf{Full} & \textbf{H0} & \textbf{H1} & \textbf{H0+H1} & \textbf{Signal} \\
\midrule
\multirow{2}{*}{MobiFall}
 & LR  & $\mathbf{0.9046}$ & $0.8842$ & $0.6172$ & $0.8892$ & $0.6218$ \\
 & SVM & $\mathbf{0.8913}$ & $0.8685$ & $0.6523$ & $0.8784$ & $0.7241$ \\
\midrule
\multirow{2}{*}{SisFall}
 & LR  & $\mathbf{0.9710}$ & $0.9266$ & $0.6580$ & $0.9271$ & $0.9589$ \\
 & SVM & $\mathbf{0.9701}$ & $0.9277$ & $0.6456$ & $0.9280$ & $0.9627$ \\
\midrule
\multirow{2}{*}{FAD\_40Hz}
 & LR  & $\mathbf{0.9426}$ & $0.8082$ & $0.4247$ & $0.8044$ & $0.9505$ \\
 & SVM & $\mathbf{0.9592}$ & $0.8128$ & $0.4605$ & $0.8210$ & $0.9504$ \\
\bottomrule
\end{tabular}

\vspace{2pt}
{\footnotesize Representative $95\,\%$ confidence half-widths (LR Full):
MobiFall $\pm0.0392$, SisFall $\pm0.0107$, FAD\_40Hz $\pm0.0367$; the H1-only columns
carry the widest intervals (up to $\pm0.077$).}
\end{table}

\subsection{Within-topology contribution: H0 dominates H1}

For both classifiers and across all three datasets, the H0 block alone attains $F_2$
between $0.81$ and $0.93$, whereas the H1 block alone is the weakest configuration in
the entire study ($0.42$--$0.66$).  Adding H1 to H0 changes $F_2$ negligibly---for LR,
from $0.8842$ to $0.8892$ on MobiFall, $0.9266$ to $0.9271$ on SisFall, and $0.8082$
to $0.8044$ on FAD\_40Hz (a slight \emph{decrease}); the SVM rows show the same pattern.
The discriminative topological information is therefore concentrated almost entirely in
dimension zero.  This is consistent with the structure of the delay-embedded
acceleration signal: a fall is a transient, high-energy excursion whose phase-space
reconstruction is characterised by the spread and clustering of points (an H0
phenomenon), while one-dimensional loops carry comparatively little fall-versus-ADL
information and are more sensitive to noise---which also explains the wide confidence
intervals of the H1-only columns.  A practical corollary is that the H1 statistics
could be dropped with little loss, roughly halving the persistent-homology computation.

\subsection{Topology versus signal: a dataset-dependent division of labour}

The relative value of the topological blocks and the four signal descriptors differs
sharply across datasets.  On MobiFall---a smartphone carried loosely in a trouser
pocket---the signal block alone reaches only $F_2 = 0.6218$ (LR) or $0.7241$ (SVM),
while the topological blocks reach $0.88$; here the topological representation is
indispensable and accounts for most of the performance.  On SisFall and FAD\_40Hz---%
waist-mounted devices with cleaner, more repeatable kinematics---the four signal
descriptors alone already achieve $F_2 \approx 0.95$--$0.96$, matching or exceeding the
topological blocks on those datasets.  In every case, however, the full vector matches
or exceeds the best single block, and no single block is reliably strong on all three
datasets whereas the full vector is.  The blocks are thus complementary: the
topological features supply the discriminative signal where raw statistics fail (noisy,
loosely mounted sensors), and the combination provides the cross-dataset robustness
that motivates the representation.  This refines the H1 narrative: topology is not a
uniform performance multiplier but a robustness mechanism whose marginal value is
largest precisely where simple descriptors are weakest.

\subsection{Classifier invariance}

The block ordering is essentially independent of the classifier.  Logistic regression
and the SVM agree on the ranking within every dataset---Full $\gtrsim$ best topological
or signal block $\gg$ H1---and their per-block $F_2$ values differ by less than the
confidence intervals in almost all cells.  This reinforces the H3 finding
(Section~\ref{subsec:H3}) that the representation, not the classifier family, governs
performance: ablating a feature block changes $F_2$ far more than switching classifiers.

\subsection{Scope of the ablation}

This study ablates the three \emph{feature blocks} while holding the delay embedding
and the optimal topological configuration $\lambda^*$ fixed.  It therefore isolates the
value of each summary block but does not, on its own, separate the contribution of the
phase-space reconstruction itself; a complementary ablation that removes the delay
embedding requires re-extracting features under a modified pipeline and is left to the
future work of Section~\ref{subsec:future-work}.

% ─────────────────────────────────────────────────────────────
\section{Cross-Dataset Synthesis}
\label{subsec:synthesis}

\subsection{Performance ranking and subject pool}

Under the General protocol, the consistent ranking by $F_2$ is:
\[
  \text{SisFall} > \text{FAD\_40Hz} > \text{MobiFall.}
\]
This ranking holds for both LR and SVM.  Three factors explain it jointly.

\textbf{Subject pool size and data volume.}  SisFall has the largest pool (24 subjects,
15 fall types), providing the most diverse and representative training signal for
LOSO.  MobiFall has the smallest (9 subjects), making each LOSO fold highly sensitive
to the particular subject removed.

\textbf{Fall type diversity.}  SisFall covers 15 scripted fall types executed by
elderly participants, producing a rich topology of fall trajectories.  FAD covers
diverse activities of daily life as negatives, which increases the specificity
challenge.  MobiFall's smaller scale limits the variety of fall geometries seen by
each LOSO model.

\textbf{Optimised representation.}  The Bayesian proxy $F_2$ at the optimal $\lambda^*$
is $0.9746$ (MobiFall), $0.9733$ (SisFall), $0.9860$ (FAD\_40Hz)---suggesting that
the optimisation found similarly good configurations for all three, and that the
General-protocol gap between datasets is not due to a representation quality
difference but to the structural differences above.

\subsection{Classifier preference}

LR is preferred on MobiFall (gap $+0.0131$) and SisFall (gap $+0.0009$); SVM on
FAD\_40Hz (gap $+0.0163$).  All gaps are $<0.02$; classifier choice is secondary
to representation quality (H3).  The marginal SVM advantage on FAD\_40Hz may
reflect the RBF kernel's ability to resolve residual nonlinear structure in the
40~Hz-resampled feature cloud, but the evidence is insufficient to assert a
mechanistic explanation at the current sample size.

\subsection{Computational profile}

Table~\ref{tab:compute} summarises the wall-clock cost of each pipeline stage on the
UHeM HPC cluster (dual Intel Xeon Gold 6148, 40 cores, no GPU).

\begin{table}[h!]
\centering
\caption{Wall-clock cost per dataset stage.  All computations are CPU-only.}
\label{tab:compute}
\small
\begin{tabular}{lll}
\toprule
\textbf{Stage} & \textbf{Tool} & \textbf{Cost} \\
\midrule
Bayesian optimisation (Stage~1, 200 trials) & Optuna/TPE & ${\sim}75$~min \\
Persistent homology (per trial, inner CV)   & GUDHI      & ${\sim}20$~s (MobiFall) \\
GridSearchCV (Stage~2)                      & scikit-learn & ${\sim}$minutes \\
Decision-threshold sweep (all subjects)     & NumPy      & $<1$~s \\
\bottomrule
\end{tabular}
\end{table}

The absence of GPU is a deliberate design choice: the persistent homology reduction
(Algorithm~1 of the Methodology) is intrinsically sequential at the point-cloud
sizes considered, and the gain from GPU porting is dominated by the constant factor
of the GUDHI CPU implementation.  This has a direct consequence for deployment:
the pipeline can be executed on commodity hardware and embedded processors that lack
GPU acceleration (Section~\ref{subsec:deployment}).

% ─────────────────────────────────────────────────────────────
\section{Deployment Pathway}
\label{subsec:deployment}

The experimental results point toward a concrete deployment pathway for the proposed
pipeline on consumer wearable and smartphone hardware.  This section articulates the
evidence-based engineering argument.

\subsection{Sensor context}

All experiments in this thesis use accelerometer data collected from a sensor mounted
at the waist or carried in the trouser pocket.  This placement is both practically
common (smartphones are carried in the pocket by a large fraction of the target
population) and metrologically appropriate: the waist sensor records the acceleration
of the body's centre of mass, which undergoes the most discriminative trajectory
during a fall.  The SMV transformation $\mathrm{SMV}(t)=\|\mathbf{x}(t)\|_2$ provides
orientation invariance, eliminating the need for fixed sensor orientation or a
gyroscope---a significant simplification for consumer deployment.

\subsection{CPU-only operation}

The pipeline's entire compute budget is CPU-bound.  Persistent homology at the
window sizes and embedding dimensions considered here ($m=3$--$5$,
$|P_\lambda|\leq150$ points) completes in tens of milliseconds per window on a
desktop-class CPU core.  The threshold sweep adds negligible cost.  This profile is
compatible with:
\begin{itemize}
  \item \textbf{Smartphones}: ARM Cortex-A-class processors (e.g., Qualcomm
    Snapdragon 8xx) comfortably execute GUDHI-class computations in real time.
    The pipeline can run as a background service consuming a modest fraction of
    one core.
  \item \textbf{Smartwatches}: ARM Cortex-M-class MCUs (e.g., Wear OS devices)
    are more constrained, but the pipeline's window-level computation---feature
    extraction plus a matrix--vector multiply for the classifier---is within the
    budget of modern smartwatch processors.  A dedicated ASIC or DSP block for the
    boundary matrix reduction would further reduce energy consumption.
\end{itemize}

\subsection{Prospective threshold personalisation}

The retrospective sweep of Section~\ref{subsubsec:sweep-dist} establishes that
subject-specific thresholds improve $F_2$ by up to $+0.1165$ without retraining.
A prospective deployment would translate this into a small on-device calibration
procedure:
\begin{enumerate}
  \item During an initial wearing period, a small number of confirmed falls (or
    fall-simulation events during setup) are used as a calibration set $\mathcal{C}_s$
    for subject $s$.
  \item The precomputed probability scores on $\mathcal{C}_s$ are swept over
    $\T$ to identify $\tau^*(s)$.
  \item Subsequent detections use $\tau^*(s)$ as the operating threshold.
\end{enumerate}
This procedure requires no transmission of raw sensor data, no server round-trip for
retraining, and no access to the original training data.  It is executable on-device
in under one second once the scores are available, making it compatible with
privacy-preserving on-device deployment.  The MobiFall results illustrate the value
of this mechanism: full personal retraining degrades $F_2$ from $0.9049$ to $0.7325$
(LR) because 9 subjects provide insufficient training data, but threshold sweep
still recovers $+0.0348$ from an already-calibrated baseline.

\subsection{Sensor placement and orientation robustness}

A limitation of the current work is that sensor placement is fixed (waist/pocket) and
orientation invariance is achieved solely through SMV.  The datasets include no
variation in placement.  In a consumer deployment scenario, users may wear a device
at the wrist (smartwatch), chest, or thigh, or orient a pocket phone differently from
day to day.  The stability theorem guarantees noise robustness within a placement, but
does not extend across fundamentally different fall trajectory geometries induced by
different placements.  Addressing this is an open engineering problem and a priority
for future work (Section~\ref{subsec:future-work}).

% ─────────────────────────────────────────────────────────────
\section{Limitations}
\label{subsec:limitations}

\paragraph{1. No baseline comparison.}
The planned comparison against threshold-based, statistical-feature, and spectral-feature
baselines (proposal $\S$5.4) was not executed.  Consequently, H1 is validated only
against its own pre-specified targets; the claim that TDA is superior to simpler
methods cannot be made from the current evidence.  Absolute performance numbers are
meaningful only relative to a reference, and this reference is absent.  Establishing
this comparison is the highest-priority open task before thesis finalisation.

\paragraph{2. Cross-application drops not computed.}
H2 is confirmed qualitatively (all optimal configurations differ) but the quantitative
prediction---transplanting $\lambda^*_K$ to a different dataset $K'$ incurs a
measurable $F_2$ drop (Methodology, Section~4.8)---was not tested.  Computing
$\mathrm{Drop}_{K\to K'}$ requires re-extracting features for each off-diagonal
$(\lambda^*_K, K')$ pair, which represents a substantial additional compute budget.

\paragraph{3. No ablation study.}
The feature vector $\phi_\lambda\in\R^{24}$ contains eight persistence statistics
for each of three homology groups ($H_0$, $H_1$) and two signal statistics.  The
relative importance of these components is unknown: it is possible that a subset of
$d<24$ features carries most of the discriminative information, and that the
remaining features add noise.  An ablation or feature-importance study would
characterise which sub-map $\phi_\lambda^{(d)}$ is load-bearing.

\paragraph{4. Fixed sensor placement.}
All experiments assume the sensor is located at the waist or in the trouser pocket.
No data with wrist, chest, or thigh placement is evaluated.  The deployment argument
of Section~\ref{subsec:deployment} assumes this placement is maintained; robustness
to placement variation is untested.

\paragraph{5. Retrospective threshold sweep.}
The sweep identifies $\tau^*(s)$ on the same test windows on which it is evaluated,
introducing an optimism bias.  The reported $\overline{\Delta F_2}$ values are
retrospective diagnostics of heterogeneity, not valid prospective performance
estimates.  A prospective validation would require a separate held-out calibration
set per subject.

% ─────────────────────────────────────────────────────────────
\section{Future Work}
\label{subsec:future-work}

\subsection{Online threshold adaptation}

The most immediate extension is to replace the retrospective sweep with a prospective
online mechanism.  As a deployed device accumulates confirmed event labels (from user
feedback, caregiver verification, or clinical annotations), the per-subject threshold
$\tau^*(s)$ can be updated in a streaming fashion.  Because $\tau^*(s)$ is a scalar
parameter estimated from a small set of probability scores, even $5$--$10$ labelled
fall events are sufficient for a reliable estimate.  A natural formulation is a
Bayesian update on the threshold parameter with a conjugate prior, requiring no
gradient computation and suitable for on-device execution.

\subsection{Embedded real-time benchmark}

The CPU-only pipeline provides a plausible path to embedded deployment, but no
real-time benchmark on target hardware has been performed.  Profiling the pipeline
on an ARM Cortex-A processor (smartphone-class) and an ARM Cortex-M processor
(smartwatch-class) would establish latency and energy consumption figures.  Key
bottlenecks to characterise are the boundary matrix reduction in GUDHI (the
most compute-intensive step) and memory allocation for the feature matrix.  A
pruned, fixed-point implementation of the 24-dimensional feature extraction---
bypassing Python overhead---would likely achieve real-time performance ($\leq50$~ms
per window at 50~Hz) on a mid-range smartphone processor.

\subsection{Multi-placement sensor fusion}

Current work assumes a single fixed sensor placement.  A natural extension is to
incorporate multiple simultaneous placements (waist + wrist + thigh) and learn a
combined topological representation.  The delay-embedding and filtration steps apply
independently to each placement's signal; the resulting persistence diagrams can be
fused either by concatenation of feature vectors or by a learned weighting.  The
classifier-agnostic Bayesian optimisation framework developed here extends
straightforwardly to this setting by augmenting $\Lambda$ with a placement-selection
dimension.

\subsection{Cross-dataset transfer of $\lambda^*$}

The cross-application drop matrix $M_{K,K'}$ (Methodology, Section~4.8)---not
computed in the current work---would reveal how much configuration knowledge
transfers across datasets.  If $M_{K,K'}\approx M_{K,K}$ for some off-diagonal
pair, the Bayesian search budget could be substantially reduced for new datasets
by warm-starting from an existing $\lambda^*_K$.  This is practically relevant for
clinical deployments where collecting a full new dataset is expensive.

\subsection{Extension to other wearable time series}

The TDA pipeline introduced here---delay embedding of a magnitude signal,
persistence homology filtration, vectorised 24-dimensional descriptor,
classifier-agnostic Bayesian optimisation---is not specific to fall detection.
Any wearable time series in which discriminative information is encoded in the
\emph{shape} of short signal segments is a candidate:

\begin{itemize}
  \item \textbf{Gait analysis}: distinguishing Parkinsonian gait from healthy gait
    or post-surgical recovery via topological features of the stride cycle.
  \item \textbf{Gesture recognition}: mapping gesture trajectories to topological
    descriptors for device-control interfaces on smartwatches.
  \item \textbf{Activity recognition}: generalising beyond binary (fall/ADL) to
    multiclass (walk, run, climb, sit, fall) classification by extending the
    threshold sweep to a multiclass decision boundary.
  \item \textbf{Physiological monitoring}: ECG or PPG segments contain loop
    structures ($H_1$ generators) whose persistence captures morphological
    variability; TDA could complement classical morphological features.
  \item \textbf{Industrial IoT anomaly detection}: vibration signals from rotating
    machinery exhibit characteristic $H_0$/$H_1$ patterns at fault inception;
    the same pipeline could detect bearing or gear anomalies without labelled fault data
    (using one-class formulations).
\end{itemize}

In each case, the key requirement is that the signal windows be short enough to
delay-embed at modest dimension ($m\leq5$, $|P|\leq200$ points) and that the
discriminative information be geometrically structured.  The stability theorem
guarantees robustness to sensor noise in all these settings.

\subsection{One-class fall detection from ADL-only supervision}
\label{subsubsec:one-class}

The binary classifier developed in this thesis requires labelled fall windows for
training.  In a consumer deployment scenario, labelled falls are expensive to obtain:
they require a controlled data-collection session or a sufficiently long natural
observation period.  A newly installed device has no fall history at all---the
\emph{cold-start problem}.

An anomaly-detection formulation eliminates this requirement.  Rather than learning to
separate falls from ADL, a one-class model learns the boundary of the user's
\emph{normal} ADL behaviour exclusively from unlabelled, continuously collected sensor
data.  Any window whose topological descriptor $\phi_\lambda(w)\in\R^{24}$ lies far
from the learned ADL manifold is flagged as a candidate fall.  Concretely, a
one-class SVM~\parencite{scholkopf2001estimating} or an isolation forest trained on
$\{\phi_\lambda(w)\colon w\in\mathcal{W}_{\mathrm{ADL}}^{(s)}\}$ defines a
decision boundary in $\R^{24}$ without any fall-class supervision.

Three properties make the TDA feature space particularly well-suited for this
formulation.

\begin{enumerate}
\item \textbf{Geometric separability.}  Fall windows exhibit pronounced $H_0$
  structure (a sudden large cluster spread caused by the impulsive excursion in the
  delay-embedded cloud) that is structurally absent in ADL windows.  The Cohen--Steiner
  stability theorem ensures that this geometric difference is preserved robustly
  across subjects and sensor noise levels.

\item \textbf{Inherent personalisation.}  Each user's ADL cloud in $\R^{24}$ reflects
  that individual's movement style, pace, and activity repertoire.  A one-class model
  trained on a user's own ADL is a personalised model by construction, with no explicit
  personalisation step required.

\item \textbf{Threshold-sweep compatibility.}  The anomaly score (distance to the ADL
  boundary) is a scalar, and the decision-threshold sweep mechanism of
  Section~\ref{subsubsec:sweep-dist} applies directly: sweeping the boundary standoff
  distance over a grid $\mathcal{T}$ trades recall against precision in exactly the
  same way as the probabilistic threshold.
\end{enumerate}

The expected engineering trade-off is an elevated false-positive rate compared to
the binary classifier: vigorous ADL activities (running, stair climbing, sudden
direction changes) may appear anomalous relative to a training distribution consisting
mainly of slow or sedentary activities.  Mitigations include diverse ADL coverage
during the initial calibration period and an activity-aware normalization that
conditions the anomaly score on the subject's current motion intensity level.

From a deployment perspective, the one-class approach enables a qualitatively
different product workflow: a user installs the application, wears the device for
two to three days of ordinary activity, and the personalised fall detector is
ready---without any fall simulation, clinical annotation, or server-side retraining.
This is the most direct path from the current thesis to a consumer-ready product.

\subsection{Few-shot personalisation beyond threshold}

An alternative personalisation axis beyond threshold tuning is to adapt the feature
map $\phi_\lambda$ itself to the individual subject via meta-learning.  A model-agnostic
meta-learning (MAML) formulation could learn a $\lambda$ initialisation that adapts
quickly to a new subject with $k=5$--$10$ labelled examples.  This extends H4 from
threshold personalisation to full representation personalisation, and remains
compatible with the CPU-only constraint if the inner loop is restricted to scalar
hyperparameter updates.


% ============================================================
%  END ch5.tex
% ============================================================

% ============================================================
%  BEGIN ch6.tex
% ============================================================
% ─────────────────────────────────────────────────────────────
\phantomsection
\chapter{Conclusion}
\label{sec:conclusion}

This thesis investigated whether persistent homology, applied to short wearable
accelerometer windows via a delay-embedding pipeline, can support recall-oriented
fall detection that is simultaneously rigorous in its evaluation, principled in its
feature construction, and lightweight enough for edge deployment.  The investigation
was organised around four research hypotheses and carried out on three public
benchmark datasets---MobiFall, SisFall, and FAD\_40Hz---under a three-protocol
subject-aware evaluation framework.  The findings are summarised below at three
levels of description: mathematical, statistical, and engineering.

% ─────────────────────────────────────────────────────────────
\section{Mathematical Perspective}
\label{subsec:conclusion-math}

\paragraph{Representation effectiveness (H1).}
The feature map $\phi_\lambda\colon\R^{n\times 3}\to\R^{24}$, with hyperparameters
$\lambda^*_K$ identified by Bayesian optimisation, achieves mean subject-general
$F_2$ of $0.9709$ (LR) and $0.9700$ (SVM) on SisFall, $0.9429$ (LR) and
$0.9592$ (SVM) on FAD\_40Hz, and $0.9049$ (LR) and $0.8918$ (SVM) on MobiFall
under the leave-one-subject-out protocol.  All values exceed the pre-specified
thresholds ($\geq 0.93$ on SisFall and FAD\_40Hz; $\geq 0.80$ on MobiFall) for at
least one classifier, and recall is above $0.95$ across all dataset-classifier
combinations.  H1 is conditionally confirmed: the representation meets its own
targets, subject to the caveat that no baseline comparison has been conducted.

The theoretical basis for this effectiveness is the Cohen--Steiner stability
theorem, which guarantees that the bottleneck distance between two persistence
diagrams is bounded by the supremum norm of the perturbation between the
corresponding input functions.  This bound is a mathematical theorem, not an
empirical observation, and it provides a formal noise-robustness certificate that
no competing feature family---handcrafted statistical features, spectral
coefficients, or deep-network embeddings---possesses in comparable form.

\paragraph{Dataset-specific topological geometry (H2).}
The three datasets converged to structurally distinct optimal configurations:
MobiFall selected a SparseRips complex at a 5.0~s window with Chebyshev distance,
SisFall selected an Alpha complex at a 1.0~s window, and FAD\_40Hz selected a
standard Rips complex at a 1.0~s window with Manhattan distance.  Every dimension of
the hyperparameter tuple $\lambda$ varies across at least two datasets.  H2 is
qualitatively confirmed: no universal configuration exists for these three datasets.

Mathematically, this finding is consistent with the geometry of the underlying
delay-embedded point clouds.  The Alpha complex is derived from the Euclidean
Delaunay triangulation and is appropriate when the point cloud is convex and
well-conditioned; SisFall's controlled, elderly-subject falls produce exactly this
geometry.  The SparseRips complex handles large, noisy clouds efficiently by
subsampling; MobiFall's continuous-wear sensor data requires this treatment.  The
differing window lengths encode a genuine difference in fall dynamics: MobiFall falls
require a 5.0~s window to capture the pre-impact, impact, and post-fall phases
jointly, while SisFall and FAD\_40Hz falls are discriminable in 1.0~s.

\paragraph{Representation dominates classifier choice (H3).}
The absolute difference in General-protocol $F_2$ between LR and SVM is $0.0131$
on MobiFall, $0.0009$ on SisFall, and $0.0163$ on FAD\_40Hz---all below $0.02$.
By contrast, the $F_2$ range across 200 Optuna trials is $0.179$ on MobiFall,
$0.053$ on SisFall, and $0.108$ on FAD\_40Hz.  The ratio of $\lambda$-tuning
variation to classifier gap is $13.7\times$ on MobiFall, $59\times$ on SisFall,
and $6.6\times$ on FAD\_40Hz.

H3 is strongly confirmed.  The mathematical interpretation is that $\phi_{\lambda^*}$
with the optimal representation transforms the nonlinear fall-versus-ADL boundary
into an approximately linearly separable problem in $\R^{24}$.  Once this
transformation is achieved, both logistic regression (linear boundary) and SVM
(maximum-margin boundary, with or without an RBF kernel) operate near their
theoretical limit, producing near-equivalent results.  The signal resides in the
map $\phi_{\lambda^*}$, not in the classifier class.

% ─────────────────────────────────────────────────────────────
\section{Statistical Perspective}
\label{subsec:conclusion-stat}

\paragraph{Threshold personalisation as a distributional phenomenon (H4).}
The subject-level sweep gain $\Delta F_2(s) = F_2(\tau^*(s),s) - F_2(\hat\tau_0(s),s)$,
where $\hat\tau_0(s)$ is the PR-curve trained threshold, has a right-skewed
distribution in all six dataset-classifier combinations.  The mean gains are
$0.035$--$0.047$ on MobiFall, $0.025$--$0.035$ on FAD\_40Hz, and
$0.004$--$0.008$ on SisFall.  The coefficient of variation exceeds $50\%$ in all
cases, reaching $138\%$ for SisFall LR---indicating that the distribution is highly
dispersed relative to its mean.

H4 is confirmed in a refined form.  The proposal predicted that the distribution
of $\Delta_s$ would contain both positive and negative values.  In practice, all
gains are non-negative by construction (the sweep is an argmax), but the
distribution is far from uniform: it concentrates near zero for the majority of
subjects while a small minority achieves gains exceeding $0.10$ units of $F_2$.
Subject~11 on FAD\_40Hz (LR) gains $+0.1165$, lifting $F_2$ from $0.8046$ to
$0.9211$ by adjusting the threshold alone.  This heterogeneity is incompatible with
a uniform-shift explanation: if the classifier's posterior probabilities were
globally miscalibrated in the same direction for all subjects, a single global
threshold offset would equalise all gains.  The observed distribution implies that
the optimal decision boundary location is genuinely subject-specific.

\paragraph{Evaluation integrity under LOSO.}
The leave-one-subject-out protocol with $R=15$ independent class-balancing
repetitions separates two sources of variability that are conflated in random
train-test splits: (i) within-subject repetition variance, captured by the
repetition-level $t$-confidence intervals, and (ii) across-subject heterogeneity,
visible in the per-subject tables.  The Naive protocol consistently produces results
comparable to or marginally lower than the General protocol on SisFall (Naive
$F_2=0.9713$ vs.\ General $0.9709$ for LR), suggesting that LOSO is not
materially penalised by train-test class mismatch on large, diverse datasets.  On
MobiFall, where the subject pool is small (9 subjects), the repetition-level
confidence intervals are wider ($\pm0.039$ for LR), correctly reflecting the
limited statistical power of a 9-fold LOSO evaluation.

The Personal protocol reveals a further finding: full per-subject model training
degrades $F_2$ severely on MobiFall ($\Delta = -0.172$ for LR) but only modestly on
SisFall ($\Delta = -0.030$ for LR), with the magnitude of degradation inversely
correlated with the subject pool size.  This confirms that personalised model
training is only beneficial when sufficient per-subject data is available, and that
the threshold sweep is the appropriate personalisation mechanism when data is sparse.

% ─────────────────────────────────────────────────────────────
\section{Engineering Perspective}
\label{subsec:conclusion-eng}

\paragraph{Inference-time personalisation without retraining.}
The central engineering contribution of the thesis is the identification of
decision-threshold optimisation as a first-class personalisation mechanism that
is qualitatively different from all prior alternatives.  The threshold
$\tau\in\mathcal{T}=\{0.05,\ldots,0.95\}$ is the only parameter that can be
adjusted at inference time, without gradient computation, without access to the
training data, and without transmission of data to a server.  The computational
cost of a full per-subject sweep is $O(|\mathcal{T}|\cdot|W_s|)$ comparisons,
amounting to less than one second of wall time for all subjects in any of the three
datasets.  The ratio to the Bayesian optimisation cost exceeds $10^3$.

This asymmetry has a direct consequence for product design.  A wearable fall
detection system based on the proposed pipeline can be shipped with a population-level
model ($\phi_{\lambda^*}$, $f_\theta$) and personalised to each user at installation
time by sweeping the threshold over a small held-out calibration set.  No server
round-trip is required.  No labeled falls are required from the target user in the
first instance---the sweep can be initialised from simulated events during setup and
refined as real events accumulate.  This is the only personalisation mechanism among
those surveyed in this thesis that simultaneously satisfies the zero-retraining,
privacy-preserving, and inference-time requirements.

\paragraph{CPU-only deployment profile.}
The entire pipeline---delay embedding, simplicial filtration, boundary matrix
reduction, feature vectorization, linear classification, and threshold sweep---is
executable on commodity CPU hardware.  Persistent homology computation on the
window sizes and embedding dimensions considered here ($m \leq 5$,
$|P_\lambda| \leq 150$ points) completes in tens of milliseconds per window on a
desktop-class processor core.  The Bayesian optimisation consumed ${\sim}75$~min
of CPU time on a 40-core HPC node; this cost is paid once per dataset, not per
user.

The ARM Cortex-A processors present in current mid-range smartphones are
comfortably capable of executing the per-window inference pipeline in real time at
50~Hz.  Modern smartwatch SoCs are more constrained, but the per-window
computation---feature extraction plus a matrix-vector multiply---is within their
budget.  The absence of a GPU requirement eliminates the most significant barrier
to embedded deployment and makes the pipeline uniquely well-suited to the
battery-constrained wearable context.

\paragraph{Failure mode analysis and minimum calibration set size.}
The most significant failure mode identified in this thesis is the collapse of
the Personal protocol on MobiFall: $F_2$ drops from $0.9049$ (General) to
$0.7325$ (Personal) for LR, and from $0.8918$ to $0.5981$ for SVM.  The cause is
statistical rather than algorithmic: MobiFall contains 9 subjects with
approximately 20--30 fall instances per subject, which is insufficient for stable
within-subject classifier estimation.  The implication for system designers is
concrete: a minimum calibration set size is required before full personal model
training is beneficial.  For the datasets studied here, that threshold lies somewhere
between the 9 subjects of MobiFall and the 13 subjects of FAD\_40Hz.  Below this
threshold, the threshold sweep is the safer personalisation mechanism.

This failure mode is not a weakness of TDA specifically; it would affect any
classifier trained on sparse per-subject data.  It is, however, an important
calibration of expectations for real deployment: the naive assumption that
``more personalization is always better'' is empirically refuted by the MobiFall
Personal protocol results.

\paragraph{Path to consumer deployment.}
Taken together, the three engineering findings---inference-time personalisation,
CPU-only operation, and the identified minimum calibration requirement---define a
concrete deployment pathway:

\begin{enumerate}
  \item A population-level model ($\phi_{\lambda^*}$, $f_\theta$) is trained once
    on a large labeled dataset and shipped with the device.
  \item At installation, the user performs a brief simulated fall session
    (or the device waits to accumulate confirmed fall events from normal wear).
    The per-subject threshold $\tau^*(s)$ is identified from the calibration
    scores in under one second.
  \item Subsequent detections use $\tau^*(s)$.  As more events accumulate, the
    threshold estimate can be refined.
  \item No server communication is required for any of steps 1--3; all computation
    is on-device.
\end{enumerate}

The natural next step beyond threshold personalisation is one-class anomaly detection:
training the model exclusively on the user's own activities of daily living,
treating falls as out-of-distribution events without requiring any labeled fall data.
The Cohen--Steiner stability of the TDA features ensures that the ADL manifold in
$\R^{24}$ is robustly estimated from normal wear data, and the threshold sweep
mechanism extends directly to the standoff distance from the one-class boundary.
This approach eliminates the cold-start problem entirely and represents the most
direct path from the current thesis to a consumer-ready, zero-label deployment.

% ─────────────────────────────────────────────────────────────
\section{Limitations and Open Questions}
\label{subsec:conclusion-limits}

Four limitations of the current work constrain the strength of the conclusions.

First, the absence of a baseline comparison means that H1 is validated only against
its own pre-specified targets.  The claim that TDA is superior to statistical-feature
or threshold-based baselines on these datasets cannot be made from the current
evidence, and establishing this comparison is the highest-priority open task.

Second, the cross-application drop matrix $M_{K,K'}$ was not computed.  H2 is
therefore confirmed only qualitatively (distinct optimal configurations) and not
quantitatively (measured performance penalty from transplanting $\lambda^*_K$ to
dataset $K'$).

Third, no ablation study was conducted.  The 24-dimensional feature vector
$\phi_\lambda$ contains $H_0$, $H_1$, and signal statistics; the relative
contribution of each subvector is unknown, and it is possible that a smaller
sub-map would achieve comparable performance with lower computational cost.

Fourth, all experiments assume a fixed sensor placement at the waist or in the
trouser pocket.  The stability theorem guarantees noise robustness within a
placement, but does not extend to fundamentally different placements (wrist, chest,
thigh).  Robustness to placement variation is an important practical question that
remains unaddressed.

% ─────────────────────────────────────────────────────────────
\section{Closing Remarks}

This thesis set out to answer a focused question: can a TDA pipeline for wearable
fall detection be simultaneously principled, lightweight, and subject-aware?  The
experimental evidence across three datasets and two classifiers supports an
affirmative answer on all three dimensions.  The pipeline is principled in the sense
that its key property---feature stability---is guaranteed by a mathematical theorem.
It is lightweight in the sense that it runs without GPU acceleration at latencies
compatible with real-time embedded operation.  It is subject-aware in the sense that
the decision threshold can be personalised at inference time, at negligible
computational cost, without retraining.

The broader lesson from H3 is perhaps the most transferable finding.  The gap between
a well-tuned and a poorly tuned topological representation ($0.10$--$0.18$ units of
$F_2$) dwarfs the gap between the best and worst classifier at the optimal
representation ($< 0.02$ units).  In the fall detection context, and likely in other
wearable time-series classification problems where the discriminative signal is
encoded in short-window shape geometry, choosing the right representation is the
dominant engineering decision.  The choice of classifier is, in an important sense,
secondary.

% ============================================================
%  END ch6.tex
% ============================================================


% ── Bibliography ────────────────────────────────────────────
\printbibliography[heading=bibintoc,title={REFERENCES}]

% ── Appendices (disabled: no appendix content yet) ──────────
% To enable: set the chapter count, then re-add the three lines below.
% \eklerbolum{0}
% 
\singlespacing
% List appendices here, e.g.:
% \textbf{APPENDIX A:} Title\\

% % Appendix content goes here.
% To add an appendix, set \eklerbolum{<n>} in tez.tex, re-enable the
% 
\singlespacing
% List appendices here, e.g.:
% \textbf{APPENDIX A:} Title\\
 / % Appendix content goes here.
% To add an appendix, set \eklerbolum{<n>} in tez.tex, re-enable the
% 
\singlespacing
% List appendices here, e.g.:
% \textbf{APPENDIX A:} Title\\
 / % Appendix content goes here.
% To add an appendix, set \eklerbolum{<n>} in tez.tex, re-enable the
% \input{eklerkapak} / \input{ekler} lines, list the appendices in
% eklerkapak.tex, and add \chapter{...} blocks below.
 lines, list the appendices in
% eklerkapak.tex, and add \chapter{...} blocks below.
 lines, list the appendices in
% eklerkapak.tex, and add \chapter{...} blocks below.


% ── CV (loaded automatically at \AtEndDocument by itutez.cls)
\ozgecmis{%
% ---------- BEGIN ozgecmis.tex ----------
\vspace{10mm}

\textbf{Name SURNAME:} Kenan Evren BOYABATLI \\

\vspace{-3mm}
\textbf{E-Mail:} kbybtli@gmail.com \\

\textbf{EDUCATION:}
\vspace{-3mm}
\begin{itemize}
  \item \textbf{B.Sc.:} 2023, Istanbul Medeniyet University, Faculty of Engineering and Natural Sciences, Department of Mathematics (GPA 3.25, ranked second in the department).
\end{itemize}

\textbf{PROFESSIONAL EXPERIENCE AND AWARDS:}
\vspace{-3mm}
\begin{itemize}
  \item Junior Data Scientist, Mapin Data (December 2023 -- August 2024); developed and implemented the TDR feature-rating algorithm on a national education dataset and built Python-based process-automation tools.
  \item Researcher, T\"UB\.ITAK Project 119F279, ``Investigation of the Effects of Surface Curvature on Interface Heat Transfer for Thermally Optimized Applications'' (April -- September 2022); implemented and optimised the numerical and machine-learning algorithms used in the project.
  \item Ranked second in the Department of Mathematics, Istanbul Medeniyet University (2023).
\end{itemize}

\textbf{PUBLICATIONS, PRESENTATIONS AND PATENTS ON THE THESIS:}
\vspace{-3mm}
\begin{itemize}
  \item ---
\end{itemize}

\newpage

\textbf{OTHER PUBLICATIONS, PRESENTATIONS AND PATENTS:}
\vspace{-3mm}
\begin{itemize}
  \item Topological data analysis study on feature rating (TDR algorithm), applied to a Ministry of National Education (MEB) school dataset, 2024.
\end{itemize}
\vspace{-3mm}

% -----------------------------------------------------
% ---------- END ozgecmis.tex ----------
}

\end{document}
